% Lesson 27 — Vocabulary: Words and expressions related to keeping informed about health problems related to drugs, tobacco, alcohol abuse, and aging — Anglais Tle A — Cours signé OnBuch+
% Programme officiel MINESEC. Compiler avec : tectonic EN27-vocabulary-words-and-expressions-related-to-keeping-inf.tex
\newcommand{\DOCMATIERE}{Anglais}
\newcommand{\DOCNIVEAU}{Tle A}
\newcommand{\DOCMODULE}{Chapter 3 : Environment, Well-being and Health}
\newcommand{\DOCLECON}{EN27}
\newcommand{\DOCTITRE}{Lesson 27 — Vocabulary: Words and expressions related to keeping informed about health problems related to drugs, tobacco, alcohol abuse, and aging}
% OnBuch+ — préambule « cours signés » ENRICHI (compilé avec tectonic / XeLaTeX)
% Système visuel « Pop » d'OnBuch+ : fond crème (jamais blanc), bordures encre
% épaisses, ombres « dures » décalées, titres Archivo Black en capitales.
% Version MATHÉMATIQUES : graphes (pgfplots/tikz), boîtes pédagogiques
% (définition, propriété, méthode, attention, le savais-tu, exercice résolu,
% exercice, TP, bilan), page de garde et signature OnBuch+ ancrée en bas.
%
% Macros attendues dans main.tex AVANT \input{preamble} :
%   \DOCMATIERE  \DOCNIVEAU  \DOCMODULE  \DOCLECON (numéro)  \DOCTITRE
\documentclass[11pt]{article}

\usepackage{fontspec}
\usepackage[english,french]{babel} % français = langue principale ; l'anglais via \eng{..} / textebox
\usepackage[a4paper,margin=2cm,top=3.4cm,bottom=2.8cm,headheight=1.4cm,footskip=1.3cm]{geometry}
\usepackage{amsmath,amssymb,amsfonts}
\usepackage{mathtools} % \xmapsto, \coloneqq... souvent utilisés par les modèles
\usepackage{cancel} % simplifications barrées (bilans de forces)
% \abs{z} (module, valeur absolue) et \norm{u} (norme), utilisés par les modèles
\providecommand{\abs}{}\let\abs\relax\DeclarePairedDelimiter{\abs}{\lvert}{\rvert}
\providecommand{\norm}{}\let\norm\relax\DeclarePairedDelimiter{\norm}{\lVert}{\rVert}
\usepackage[table]{xcolor} % [table] : fournit aussi \rowcolors (alternance de lignes)
\usepackage{pagecolor}
\usepackage{tikz}
\usetikzlibrary{arrows.meta,positioning,calc,shapes.geometric,shapes.misc,shapes.arrows,decorations.pathmorphing,decorations.pathreplacing,decorations.markings,patterns,shadows,mindmap,trees,fit,backgrounds,matrix,intersections,angles,quotes}
\usepackage{pgfplots}
\usepackage[version=4]{mhchem} % équations nucléaires \ce{^{235}_{92}U}
\usepackage[european,siunitx]{circuitikz} % schémas électriques
\pgfplotsset{compat=1.18}
\usepgfplotslibrary{fillbetween,groupplots} % aires entre courbes (calcul intégral)
\usepackage{enumitem}
\usepackage{fancyhdr}
% [most] charge toutes les bibliothèques tcolorbox (listings, minted, raster,
% documentation...) et gonfle inutilement la mémoire TeX sur un document long
% (~300 boîtes) : on ne charge que ce qui est réellement utilisé.
\usepackage[skins,breakable]{tcolorbox}
\usepackage{graphicx}
\usepackage{colortbl}
\usepackage{booktabs}
\usepackage{tabularx}
\usepackage{diagbox} % case d'en-tête diagonale des tableaux à double entrée
% Colonnes extensibles centrée / gauche / droite (C, L, R), souvent utilisées par les modèles
\newcolumntype{C}{>{\centering\arraybackslash}X}
\newcolumntype{L}{>{\raggedright\arraybackslash}X}
\newcolumntype{R}{>{\raggedleft\arraybackslash}X}
\usepackage{array}
\usepackage{multicol}
\usepackage{multirow}
\usepackage{siunitx}
% parse-numbers=false : accepte aussi \SI{2,5 \times 10^{-3}}{...}
\sisetup{locale=FR,output-decimal-marker={,},group-minimum-digits=4,parse-numbers=false}
\DeclareSIUnit\ohm{\ensuremath{\symup{\Omega}}} % U+2126 absent de la police
\DeclareSIUnit\division{div} % réglages d'oscilloscope (ms/div, V/div)
\DeclareSIUnit\tr{tr} % tours (vitesse de rotation, ex. \SI{1500}{\tr\per\minute})
\DeclareSIUnit\molar{M} % concentration molaire (ex. \SI{3}{\molar}, \SI{1}{\micro\molar})
\DeclareSIUnit\an{an} % année (le modèle confond parfois avec \year, un compteur TeX) — ex. \SI{5}{\kilo\meter\per\an}
\DeclareSIUnit\FCFA{FCFA} % franc CFA (ex. \SI{500}{\FCFA\per\kilo\gram})
\DeclareSIUnit\degreeBrix{\textdegree Brix} % teneur en sucre d'un sirop/jus (réfractométrie)
\DeclareSIUnit\month{mois} % le modèle utilise parfois \month, un compteur TeX, comme unité de durée
\usepackage{qrcode}
% \degree : commande fréquemment utilisée par les modèles pour un angle ou une
% température en dehors de \SI{}{} (non fournie par les paquets chargés).
\providecommand{\degree}{^{\circ}}
% \qed (fin de démonstration), utilisé par les modèles sans amsthm
\providecommand{\qed}{\hfill$\square$}
% \barre : double barre des valeurs interdites dans un tableau de variations (tkz-tab)
\providecommand{\barre}{\ensuremath{\|}}
% environnement proof (amsthm non chargé) utilisé par les modèles
\ifdefined\proof\else\newenvironment{proof}[1][Démonstration]{\par\noindent\textit{#1.}\ }{\qed\par}\fi
\usepackage{lastpage}
\usepackage{hyperref}
\hypersetup{hidelinks}

% ============================================================
% Palette « Pop » OnBuch+ — mêmes hex que la classe P (plus_ui.dart)
% ============================================================
\definecolor{poporange}{HTML}{F59321}
\definecolor{popdark}{HTML}{E07A0C}
\definecolor{popcream}{HTML}{FFF6E8}
\definecolor{popink}{HTML}{1C1714}
\definecolor{poppurple}{HTML}{7A5AE0}
\definecolor{popgreen}{HTML}{2E9E6A}
\definecolor{poppink}{HTML}{E0457B}
\definecolor{popblue}{HTML}{2F7DE1}
\definecolor{popgold}{HTML}{C98A12}
\definecolor{popmuted}{HTML}{8A7F76}
\definecolor{popline}{HTML}{EADFD2}
\definecolor{popgray}{HTML}{8A7F76}
\colorlet{popgrey}{popgray}
% Alias tolérants (le modèle nomme parfois les couleurs différemment)
\colorlet{popwhite}{white}
\colorlet{popblack}{popink}
\colorlet{popred}{poppink}
\colorlet{popyellow}{popgold}
% Teintes claires (fonds de tableaux/encadrés) — le modèle nomme parfois
% les couleurs avec un suffixe L (ex. popblueL) sans les avoir définies.
\colorlet{poporangeL}{poporange!20}
\colorlet{popdarkL}{popdark!20}
\colorlet{poppurpleL}{poppurple!20}
\colorlet{popgreenL}{popgreen!20}
\colorlet{poppinkL}{poppink!20}
\colorlet{popblueL}{popblue!20}
\colorlet{popgoldL}{popgold!20}
\colorlet{popredL}{popred!20}
\colorlet{popgrayL}{popgray!20}
% Teintes claires pour fonds de tableaux / zones
\colorlet{popblueL}{popblue!10!white}
\colorlet{poporangeL}{poporange!14!white}
\colorlet{popgreenL}{popgreen!12!white}
\colorlet{poppurpleL}{poppurple!12!white}
\colorlet{poppinkL}{poppink!12!white}
\colorlet{popgoldL}{popgold!14!white}

\pagecolor{popcream}

% ============================================================
% Typographie
% ============================================================
\defaultfontfeatures{Ligatures=TeX}
\setmainfont{PlusJakartaSans-Medium.ttf}[Path=../../../fonts/,BoldFont=PlusJakartaSans-Bold.ttf]
% Symboles, API et emojis absents de la police principale : repli sur DejaVu Sans (caractères actifs)
\newfontface\symfont{DejaVuSans.ttf}[Path=../../../fonts/]
\makeatletter
\newcommand\symactive[1]{\catcode#1=13 \begingroup\lccode`\~=#1 \lowercase{\endgroup\def~}{{\symfont\char#1}}}
\newcommand\symblank[1]{\catcode#1=13 \begingroup\lccode`\~=#1 \lowercase{\endgroup\def~}{}}
\newcommand\symhyph[1]{\catcode#1=13 \begingroup\lccode`\~=#1 \lowercase{\endgroup\def~}{\mbox{-}}}
\newcount\symc
\newcommand\symrange[2]{\symc=#1 \@whilenum\symc<#2 \do{\expandafter\symactive\expandafter{\the\symc}\advance\symc by 1 }}
\newcommand\symblankrange[2]{\symc=#1 \@whilenum\symc<#2 \do{\expandafter\symblank\expandafter{\the\symc}\advance\symc by 1 }}
\makeatother
\symrange{"0250}{"02FF}\symactive{"2713}\symactive{"2714}\symactive{"2717}\symactive{"2718}\symactive{"2610}\symactive{"2611}\symactive{"2612}
\symhyph{"2011}\symhyph{"2E1A}\symblankrange{"1F300}{"1FAFF}\symblank{"FE0F}
% Maths en sans-serif (Fira Math) avec chiffres et lettres droites de Plus
% Jakarta, pour que les formules se fondent dans le texte.
\usepackage[math-style=ISO,bold-style=ISO]{unicode-math}
\setmathfont{FiraMath-Regular.otf}[Path=../../../fonts/]
\setmathfont{PlusJakartaSans-Medium.ttf}[Path=../../../fonts/,range={up/num,up/{Latin,latin},\mathup}]
\usepackage{icomma} % virgule décimale sans espace en mode math
% Caractères mathématiques Unicode absents des polices de texte (Plus Jakarta,
% Archivo Black) : rendus par leur équivalent mathématique, en texte comme en
% formule — sinon ils s'affichent en carré vide (ex. « Arithmétique dans ℤ »).
\usepackage{newunicodechar}
\newunicodechar{ℝ}{\ensuremath{\mathbb{R}}}
\newunicodechar{ℤ}{\ensuremath{\mathbb{Z}}}
\newunicodechar{ℂ}{\ensuremath{\mathbb{C}}}
\newunicodechar{ℕ}{\ensuremath{\mathbb{N}}}
\newunicodechar{ℚ}{\ensuremath{\mathbb{Q}}}
\newunicodechar{𝔻}{\ensuremath{\mathbb{D}}}
\newunicodechar{α}{\ensuremath{\alpha}}
\newunicodechar{β}{\ensuremath{\beta}}
\newunicodechar{γ}{\ensuremath{\gamma}}
\newunicodechar{δ}{\ensuremath{\delta}}
\newunicodechar{ε}{\ensuremath{\varepsilon}}
\newunicodechar{θ}{\ensuremath{\theta}}
\newunicodechar{λ}{\ensuremath{\lambda}}
\newunicodechar{μ}{\ensuremath{\mu}}
\newunicodechar{σ}{\ensuremath{\sigma}}
\newunicodechar{τ}{\ensuremath{\tau}}
\newunicodechar{φ}{\ensuremath{\varphi}}
\newunicodechar{ψ}{\ensuremath{\psi}}
\newunicodechar{ω}{\ensuremath{\omega}}
\newunicodechar{Σ}{\ensuremath{\Sigma}}
% majuscules produites par \MakeUppercase dans les titres (α-aminés -> Α)
\newunicodechar{Α}{\ensuremath{\alpha}}
\newunicodechar{Β}{\ensuremath{\beta}}
\newunicodechar{Γ}{\ensuremath{\Gamma}}
\newunicodechar{Δ}{\ensuremath{\Delta}}
\newunicodechar{Ε}{\ensuremath{\varepsilon}}
\newunicodechar{Θ}{\ensuremath{\Theta}}
\newunicodechar{Λ}{\ensuremath{\Lambda}}
\newunicodechar{Μ}{\ensuremath{\mu}}
\newunicodechar{Τ}{\ensuremath{\tau}}
\newunicodechar{Φ}{\ensuremath{\Phi}}
\newunicodechar{Ψ}{\ensuremath{\Psi}}
\newunicodechar{Ω}{\ensuremath{\Omega}}
\newunicodechar{↦}{\ensuremath{\mapsto}}
\newunicodechar{∈}{\ensuremath{\in}}
\newunicodechar{∉}{\ensuremath{\notin}}
\newunicodechar{∧}{\ensuremath{\wedge}}
\newunicodechar{⇔}{\ensuremath{\Leftrightarrow}}
\newunicodechar{⟺}{\ensuremath{\Longleftrightarrow}}
\newunicodechar{⇒}{\ensuremath{\Rightarrow}}
\newunicodechar{⟹}{\ensuremath{\Longrightarrow}}
\newunicodechar{∘}{\ensuremath{\circ}}
\newunicodechar{∪}{\ensuremath{\cup}}
\newunicodechar{∩}{\ensuremath{\cap}}
\newunicodechar{≡}{\ensuremath{\equiv}}
\newunicodechar{⊂}{\ensuremath{\subset}}
\newunicodechar{⊥}{\ensuremath{\perp}}
\newunicodechar{∃}{\ensuremath{\exists}}
\newunicodechar{∀}{\ensuremath{\forall}}
\newunicodechar{∛}{\ensuremath{\sqrt[3]{\,}}}
\newunicodechar{∜}{\ensuremath{\sqrt[4]{\,}}}
\newunicodechar{⃗}{\ensuremath{{}^{\to}}}
\newunicodechar{ⁿ}{\ifmmode^{n}\else\textsuperscript{\ensuremath{n}}\fi}
\newunicodechar{ˣ}{\ifmmode^{x}\else\textsuperscript{\ensuremath{x}}\fi}
\newunicodechar{ᵇ}{\ifmmode^{b}\else\textsuperscript{\ensuremath{b}}\fi}
\newunicodechar{ʳ}{\ifmmode^{r}\else\textsuperscript{\ensuremath{r}}\fi}
\newunicodechar{ᵘ}{\ifmmode^{u}\else\textsuperscript{\ensuremath{u}}\fi}
\newunicodechar{ʸ}{\ifmmode^{y}\else\textsuperscript{\ensuremath{y}}\fi}
\newunicodechar{ᵃ}{\ifmmode^{a}\else\textsuperscript{\ensuremath{a}}\fi}
\newunicodechar{ᵉ}{\ifmmode^{e}\else\textsuperscript{\ensuremath{e}}\fi}
\newunicodechar{ᵏ}{\ifmmode^{k}\else\textsuperscript{\ensuremath{k}}\fi}
\newunicodechar{ᵖ}{\ifmmode^{p}\else\textsuperscript{\ensuremath{p}}\fi}
\newunicodechar{ᴺ}{\ifmmode^{N}\else\textsuperscript{\ensuremath{N}}\fi}
\newunicodechar{ᶜ}{\ifmmode^{c}\else\textsuperscript{\ensuremath{c}}\fi}
\newunicodechar{ʰ}{\ifmmode^{h}\else\textsuperscript{\ensuremath{h}}\fi}
\newunicodechar{ᵠ}{\ifmmode^{q}\else\textsuperscript{\ensuremath{q}}\fi}
\newunicodechar{ₙ}{\ifmmode_{n}\else\textsubscript{\ensuremath{n}}\fi}
\newunicodechar{ₐ}{\ifmmode_{a}\else\textsubscript{\ensuremath{a}}\fi}
\newunicodechar{ᵢ}{\ifmmode_{i}\else\textsubscript{\ensuremath{i}}\fi}
\newunicodechar{ᵣ}{\ifmmode_{r}\else\textsubscript{\ensuremath{r}}\fi}
\newunicodechar{ₖ}{\ifmmode_{k}\else\textsubscript{\ensuremath{k}}\fi}
\newunicodechar{ᵦ}{\ifmmode_{\beta}\else\textsubscript{\ensuremath{\beta}}\fi}
\newunicodechar{ₓ}{\ifmmode_{x}\else\textsubscript{\ensuremath{x}}\fi}
\newfontfamily\popdisplay{ArchivoBlack-Regular.ttf}[Path=../../../fonts/]
\newfontfamily\poplogo{SpaceGrotesk-Bold.ttf}[Path=../../../fonts/]
\newfontfamily\popsemi{PlusJakartaSans-SemiBold.ttf}[Path=../../../fonts/]
\newfontfamily\popxbold{PlusJakartaSans-ExtraBold.ttf}[Path=../../../fonts/]
\newfontfamily\popxboldls{PlusJakartaSans-ExtraBold.ttf}[Path=../../../fonts/,LetterSpace=3.0]

\setlist[enumerate]{leftmargin=1.7em,itemsep=5pt,topsep=4pt}
\setlist[itemize]{leftmargin=1.7em,itemsep=4pt,topsep=4pt,label={\color{poporange}\textbullet}}
\setlength{\parindent}{0pt}
\setlength{\parskip}{5pt}
\renewcommand{\arraystretch}{1.25}

% pgfplots : style Pop par défaut
\pgfplotsset{
  every axis/.append style={axis line style={popink,line width=1pt},tick style={popink},
    grid style={popline,line width=0.5pt},label style={font=\small},tick label style={font=\footnotesize}},
  popcurve/.style={poporange,line width=1.8pt,no markers,smooth},
}
% Même style disponible dans un \draw TikZ simple (hors pgfplots)
\tikzset{popcurve/.style={poporange,line width=1.8pt}}
\tikzset{popgrid/.style={popline,very thin}}
\colorlet{popcurve}{poporange} % le nom du style est parfois utilisé comme couleur

% ============================================================
% Pill / squiggle
% ============================================================
\newcommand{\poppill}[3]{%
  \tikz[baseline=-0.55ex]\node[draw=popink,line width=1.2pt,rounded corners=2.6pt,%
    fill=#2,inner xsep=6.5pt,inner ysep=3pt]{%
    \popxboldls\fontsize{7}{7}\selectfont\color{#3}\MakeUppercase{#1}};%
}
\newcommand{\popsquiggle}[1]{%
  \tikz[baseline=-0.4ex]{
    \draw[#1,line width=2.6pt,line cap=round]
      plot [smooth] coordinates {(0,0.09) (0.34,0.01) (0.68,0.21) (1.02,0.01) (1.36,0.10)};
  }%
}
% Mot-clé surligné (marqueur orange sous le texte)
\newcommand{\cle}[1]{{\popxbold\color{popdark}#1}}

% ============================================================
% En-tête / pied
% ============================================================
\pagestyle{fancy}
\fancyhf{}
\renewcommand{\headrulewidth}{0pt}
\renewcommand{\footrulewidth}{0pt}
\lhead{\raisebox{-0.15ex}{\poplogo\fontsize{13}{13}\selectfont\color{popink}OnBuch\color{poporange}+}}
\rhead{\poppill{\DOCMATIERE\ \textbullet\ \DOCNIVEAU\ \textbullet\ Leçon \DOCLECON}{white}{popink}}
\lfoot{\popxbold\fontsize{7.6}{7.6}\selectfont\color{popmuted}\MakeUppercase{Cours signé OnBuch+}\ \textbullet\ {\popsemi\DOCTITRE}}
\rfoot{\tikz[baseline=-0.6ex]\node[rounded corners=8.5pt,fill=popink,minimum height=17pt,minimum width=17pt,inner xsep=5pt,inner sep=0]{\popxbold\fontsize{8}{8}\selectfont\color{popcream}\thepage\,/\,\pageref*{LastPage}};}

% ============================================================
% Titres
% ============================================================
\newcommand{\chaptitre}[2]{%
  \begin{tcolorbox}[enhanced,colback=poporange,colframe=popink,boxrule=2.6pt,arc=11pt,
      drop shadow=popink,left=15pt,right=15pt,top=12pt,bottom=11pt,before skip=3pt,after skip=18pt]
    \poppill{#2}{popink}{popcream}\par\vspace{9pt}
    {\raggedright\hyphenpenalty=10000\exhyphenpenalty=10000\popdisplay\fontsize{22}{24}\selectfont\color{popcream}\MakeUppercase{#1}\par}
  \end{tcolorbox}
}
% Section numérotée : pastille orange avec le numéro + titre Archivo
\newcounter{popsec}
\newcommand{\coursec}[1]{%
  \stepcounter{popsec}\setcounter{popsub}{0}%
  \par\vspace{16pt}\noindent
  \begin{minipage}{\linewidth}
  \tikz[baseline=-0.7ex]\node[draw=popink,line width=1.6pt,fill=poporange,rounded corners=4pt,minimum size=22pt,inner sep=0]{\popdisplay\fontsize{12}{12}\selectfont\color{popcream}\thepopsec};%
  \hspace{8pt}{\popdisplay\fontsize{15}{17}\selectfont\color{popink}\MakeUppercase{#1}}\par
  \vspace{4pt}\noindent{\color{poporange}\rule{36pt}{3.6pt}}
  \end{minipage}\par\nopagebreak\vspace{8pt}
}
\newcounter{popsub}
\newcommand{\courssub}[1]{%
  \stepcounter{popsub}%
  \par\vspace{9pt}\noindent{\popxbold\fontsize{11.5}{13}\selectfont\color{popink}{\color{poporange}\thepopsec.\thepopsub}\hspace{6pt}#1}\par\nopagebreak\vspace{4pt}
}

% ============================================================
% Boîtes pédagogiques — même recette : fond blanc, bordure épaisse colorée,
% ombre dure encre, badge plein. Titre optionnel [#1].
% ============================================================
\tcbset{popbase/.style={breakable,enhanced,colback=white,boxrule=2.3pt,arc=9pt,
  shadow={3pt}{-3pt}{0pt}{popink},left=14pt,right=14pt,top=11pt,bottom=11pt,before skip=11pt,after skip=15pt,
  fontupper=\popsemi\fontsize{10.6}{14.5}\selectfont\color{popink},
  fontlower=\popsemi\fontsize{10.6}{14.5}\selectfont\color{popink}}}
% Coche dessinée (le glyphe ✓ n'existe pas dans les polices du document)
\newcommand{\popcheck}[1]{\tikz[baseline=-0.5ex]\draw[#1,line width=1.6pt,line cap=round,line join=round](-0.13,0.0)--(-0.04,-0.09)--(0.13,0.1);}
\providecommand{\coche}{\popcheck{popgreen}}
\providecommand{\resultat}[1]{\par\smallskip\noindent\fcolorbox{popgreen}{popgreenL}{\parbox{\dimexpr\linewidth-2\fboxsep-2\fboxrule\relax}{\textbf{Résultat.} #1}}\par\smallskip}
\newcommand{\popboxhead}[3]{\poppill{#1}{#2}{white}\ifx\relax#3\relax\else\hspace{6pt}{\popxbold\fontsize{10.6}{12}\selectfont\color{popink}#3}\fi\par\vspace{7pt}}

\newtcolorbox{aretenir}[1][]{popbase,colframe=popgreen,before upper={\popboxhead{À retenir}{popgreen}{#1}}}
% Texte en anglais (document de lecture, dialogue, extrait) : cadre
% d'étude. Un titre optionnel (source, auteur) s'affiche dans l'en-tête.
\newtcolorbox{textebox}[1][]{popbase,colframe=popblue,before upper={\popboxhead{Text}{popblue}{#1}\begin{otherlanguage}{english}},after upper={\end{otherlanguage}}}
% Marques ✓ / ✗ (forme correcte / fautive) souvent utilisées par les modèles
\providecommand{\cross}{\textcolor{poppink}{\ensuremath{\times}}}
\providecommand{\xmark}{\cross}\providecommand{\crossmark}{\cross}\providecommand{\cmark}{\ensuremath{\checkmark}}
% Ligne de réponse / justification à compléter (inventée par les modèles)
\providecommand{\justification}[1]{\par\noindent\textit{Justification~:} \dotfill\par}
% \textipa (package tipa non chargé) : la police ne couvre pas l'API -> simple passage
\providecommand{\textipa}[1]{#1}
% Soulignement (ulem non chargé) : \uline, \ul
\providecommand{\uline}[1]{\underline{#1}}\providecommand{\ul}[1]{\underline{#1}}
% \glossaire{\item ...} inventée par les modèles : liste de glossaire
\providecommand{\glossaire}[1]{\begin{itemize}#1\end{itemize}}
% \alert (beamer) inventée par les modèles : texte mis en évidence
\providecommand{\alert}[1]{\textbf{\textcolor{poppink}{#1}}}
% variantes ASCII d'ordinaux écrites en macro
\providecommand{\iere}{\textsuperscript{re}}\providecommand{\ieme}{\textsuperscript{e}}\providecommand{\ere}{\textsuperscript{re}}\providecommand{\eme}{\textsuperscript{e}}\providecommand{\er}{\textsuperscript{er}}\providecommand{\ie}{\textsuperscript{e}}
% Ordinaux français écrits en macro par les modèles : 1\ière, 3\ième
\newcommand{\ière}{\textsuperscript{re}}\newcommand{\ième}{\textsuperscript{e}}
% \exemple (inventée par les modèles) : étiquette « Exemple : »
\providecommand{\exemple}{\textbf{Exemple~:}\ }
% \square utilisable aussi hors mode math (cases à cocher)
\AtBeginDocument{\let\squareMath\square\renewcommand{\square}{\ensuremath{\squareMath}}}
% Mot ou phrase en anglais à l'intérieur d'un paragraphe français
\newcommand{\eng}[1]{\ifmmode\text{\textit{#1}}\else\begingroup\selectlanguage{english}\itshape #1\endgroup\fi}
\let\esp\eng % alias toléré
% flèches d'intonation inventées par les modèles
\providecommand{\textsearrow}{}\renewcommand{\textsearrow}{\ensuremath{\searrow}}\providecommand{\textnearrow}{}\renewcommand{\textnearrow}{\ensuremath{\nearrow}}
% \fr{..} : mot français (inventé par les modèles)
\providecommand{\fr}[1]{\textit{#1}}
\newtcolorbox{exemplebox}[1][]{popbase,colframe=poporange,before upper={\popboxhead{Exemple}{poporange}{#1}}}
\newtcolorbox{definition}[1][]{popbase,colframe=popblue,before upper={\popboxhead{Définition}{popblue}{#1}}}
\newtcolorbox{propriete}[1][]{popbase,colframe=poppurple,before upper={\popboxhead{Propriété}{poppurple}{#1}}}
\newtcolorbox{methode}[1][]{popbase,colframe=popgold,before upper={\popboxhead{Méthode}{popgold}{#1}}}
\newtcolorbox{attention}[1][]{popbase,colframe=poppink,before upper={\popboxhead{Attention}{poppink}{#1}}}
\newtcolorbox{experience}[1][]{popbase,colframe=popblue,colback=popblueL,before upper={\popboxhead{Expérience}{popblue}{#1}}}
% « Le savais-tu ? » : carte violette pleine (contexte, culture, Cameroun)
\newtcolorbox{savaistu}[1][]{popbase,colback=poppurple,colframe=popink,
  fontupper=\popsemi\fontsize{10.4}{14.2}\selectfont\color{white},
  before upper={\poppill{Le savais-tu\,?}{white}{poppurple}\ifx\relax#1\relax\else\hspace{6pt}{\popxbold\color{white}#1}\fi\par\vspace{7pt}}}
% Exercice résolu : énoncé en haut, \tcblower puis la solution sur fond crème
\newcounter{popexo}\newcounter{popexr}
\newtcolorbox{exoresolu}[1][]{popbase,colframe=poporange,colbacklower=popcream,
  segmentation style={popink,line width=1.2pt,dashed},
  before upper={\stepcounter{popexr}\popboxhead{Exercice résolu \thepopexr}{poporange}{#1}},
  before lower={\poppill{Solution}{popink}{popcream}\par\vspace{6pt}}}
% Exercice (non résolu en place) — la correction est dans la partie Corrigés
\newtcolorbox{exercice}[1][]{popbase,colframe=popink,
  before upper={\stepcounter{popexo}\popboxhead{Exercice \thepopexo}{popink}{#1}}}
% Corrigé
\newtcolorbox{corrige}[1][]{popbase,colframe=popgreen,colback=popgreenL,
  before upper={\popboxhead{Corrigé}{popgreen}{#1}}}
% Objectifs / prérequis / bilan
\newtcolorbox{objectifs}{popbase,colframe=popink,colback=poporangeL,
  before upper={\popboxhead{Objectifs de la leçon}{popink}{}}}
\newtcolorbox{prerequis}{popbase,colframe=popink,colback=popblueL,
  before upper={\popboxhead{Prérequis}{popblue}{}}}
\newtcolorbox{bilan}{popbase,colframe=popink,colback=white,boxrule=2.8pt,
  before upper={\popboxhead{Fiche bilan}{poporange}{L'essentiel à connaître pour l'examen}}}
% Figure encadrée avec légende
\newtcolorbox{popfigure}[1][]{enhanced,breakable=false,colback=white,colframe=popline,boxrule=1.4pt,arc=8pt,
  left=10pt,right=10pt,top=10pt,bottom=8pt,before skip=10pt,after skip=12pt,halign=center,
  lower separated=false,halign lower=center,
  fontlower=\popsemi\fontsize{9}{11.5}\selectfont\color{popmuted},
  #1}
\newcommand{\legende}[1]{\par\vspace{6pt}{\popsemi\fontsize{9}{11.5}\selectfont\color{popmuted}#1}}

% ============================================================
% Signature OnBuch+
% ============================================================
\newcommand{\poplogoinline}[1]{{\poplogo\fontsize{#1}{#1}\selectfont\color{popink}OnBuch\color{poporange}+}}
\newcommand{\signatureonbuch}{%
  \begin{tcolorbox}[enhanced,colback=popink,colframe=popink,boxrule=2.2pt,arc=10pt,
      drop shadow=poporange,left=14pt,right=14pt,top=10pt,bottom=10pt,before skip=0pt,after skip=0pt]
    \tikz[baseline=-0.7ex]\node[circle,fill=popgreen,draw=popcream,line width=1.3pt,minimum size=22pt,inner sep=0]{\popcheck{white}};%
    \hspace{9pt}\begin{minipage}[c]{0.62\linewidth}
      {\popxbold\fontsize{12}{13}\selectfont\color{popcream}Signé\ \textbullet\ {\poplogo OnBuch\color{poporange}+}}\par\vspace{2pt}
      {\raggedright\popsemi\fontsize{8}{10.5}\selectfont\color{popline}\MakeUppercase{Équipe pédagogique OnBuch+}\\\MakeUppercase{Conforme au programme officiel MINESEC}\par}
    \end{minipage}\hfill
    {\poplogo\fontsize{22}{22}\selectfont\color{popcream}OnBuch\color{poporange}+}
  \end{tcolorbox}%
}

% ============================================================
% Page de garde
% #1 titre · #2 matière · #3 niveau · #4 module · #5 n° leçon · #6 durée ·
% #7 description / objectifs (texte ou itemize)
% ============================================================
\newcommand{\pagegarde}[7]{%
  \thispagestyle{empty}
  \newgeometry{a4paper,margin=1.7cm,top=1.6cm,bottom=1.4cm}%
  \begin{tikzpicture}[remember picture,overlay]
    \fill[poporange!18] ([xshift=-3.2cm,yshift=-3.6cm]current page.north east) circle (4.6cm);
    \fill[poppurple!14] ([xshift=2.4cm,yshift=7.6cm]current page.south west) circle (3.8cm);
  \end{tikzpicture}%
  {\poplogo\fontsize{30}{30}\selectfont\color{popink}OnBuch\color{poporange}+}\hfill
  \raisebox{0.6ex}{\poppill{Cours signé \textbullet\ Premium}{popink}{popcream}}\par
  \vspace{1pt}\popsquiggle{poppurple}\par
  \vspace{8pt}
  \begin{tcolorbox}[enhanced,colback=poporange,colframe=popink,boxrule=2.8pt,arc=13pt,
      drop shadow=popink,left=18pt,right=18pt,top=12pt,bottom=13pt,after skip=10pt]
    \poppill{#2 \textbullet\ #3}{popink}{popcream}\hspace{5pt}\poppill{#4}{white}{popink}\par\vspace{8pt}
    {\popdisplay\fontsize{14}{14}\selectfont\color{popink}LEÇON #5}\par\vspace{4pt}
    {\raggedright\hyphenpenalty=10000\exhyphenpenalty=10000\popdisplay\fontsize{25}{27}\selectfont\color{popcream}\MakeUppercase{#1}\par}
    \vspace{8pt}
    {\popxbold\fontsize{9.5}{9.5}\selectfont\color{popink}\MakeUppercase{Durée conseillée : #6}}
  \end{tcolorbox}
  \begin{tcolorbox}[enhanced,colback=white,colframe=popink,boxrule=2.2pt,arc=10pt,
      drop shadow=popink,left=16pt,right=16pt,top=10pt,bottom=10pt,after skip=8pt]
    \poppill{Dans cette leçon}{poppurple}{white}\par\vspace{5pt}
    \popsemi\fontsize{9.3}{12.6}\selectfont\color{popink}#7
  \end{tcolorbox}
  \noindent\begin{minipage}[t]{0.487\linewidth}
    \begin{tcolorbox}[enhanced,colback=popgreen,colframe=popink,boxrule=2.2pt,arc=10pt,
        drop shadow=popink,left=11pt,right=11pt,top=10pt,bottom=10pt,height=3.1cm,valign=center]
      \tcbox[colback=white,colframe=popink,boxrule=1.3pt,arc=5pt,boxsep=0pt,nobeforeafter,
        left=3pt,right=3pt,top=3pt,bottom=3pt,valign=center]{\qrcode[height=1.9cm]{https://play.google.com/store/apps/details?id=com.luvvix.onbuch&hl=fr}}%
      \hspace{8pt}\begin{minipage}[c]{0.5\linewidth}
        {\popdisplay\fontsize{11}{12.5}\selectfont\color{white}\MakeUppercase{Télécharge OnBuch}}\par\vspace{4pt}
        {\raggedright\popsemi\fontsize{8.4}{11}\selectfont\color{white}Gratuit \textbullet\ Google Play\\Tuteur IA, annales, concours, résultats\par}
      \end{minipage}
    \end{tcolorbox}
  \end{minipage}\hfill
  \begin{minipage}[t]{0.487\linewidth}
    \begin{tcolorbox}[enhanced,colback=poppurple,colframe=popink,boxrule=2.2pt,arc=10pt,
        drop shadow=popink,left=11pt,right=11pt,top=10pt,bottom=10pt,height=3.1cm,valign=center]
      \tikz[baseline=-0.7ex]\node[circle,fill=white,draw=popink,line width=1.3pt,minimum size=1.6cm,inner sep=0]{\popdisplay\fontsize{24}{24}\selectfont\color{poppurple}+};%
      \hspace{8pt}\begin{minipage}[c]{0.6\linewidth}
        {\popdisplay\fontsize{11}{12.5}\selectfont\color{white}\MakeUppercase{Passe à OnBuch+}}\par\vspace{4pt}
        {\raggedright\popsemi\fontsize{8.4}{11}\selectfont\color{white}Tous les cours signés, Tuteur IA illimité, fiches PDF — onglet OnBuch+ de l'app\par}
      \end{minipage}
    \end{tcolorbox}
  \end{minipage}
  \vspace{8pt}
  \popcontenu
  \vfill
  \signatureonbuch
  \restoregeometry
  \newpage
}

% Rangée « Ce cours contient » (page de garde)
\newcommand{\poptile}[3]{% #1 couleur #2 gros texte #3 légende
  \begin{tcolorbox}[enhanced,colback=#1,colframe=popink,boxrule=2pt,arc=9pt,shadow={2.5pt}{-2.5pt}{0pt}{popink},
      left=6pt,right=6pt,top=9pt,bottom=9pt,halign=center,valign=center,height=1.75cm,width=\linewidth,nobeforeafter]
    {\popdisplay\fontsize{12.5}{13}\selectfont\color{white}\MakeUppercase{#2}}\par\vspace{3pt}
    {\centering\popsemi\fontsize{7.8}{9.5}\selectfont\color{white}#3\par}
  \end{tcolorbox}}
\newcommand{\popcontenu}{%
  \par\noindent{\popxboldls\fontsize{7.5}{7.5}\selectfont\color{popmuted}CE COURS SIGNÉ CONTIENT}\par\vspace{7pt}
  \noindent\begin{minipage}[t]{0.235\linewidth}\poptile{popblue}{Cours}{complet, illustré, pas à pas}\end{minipage}\hfill
  \begin{minipage}[t]{0.235\linewidth}\poptile{poporange}{Méthodes}{et exercices résolus}\end{minipage}\hfill
  \begin{minipage}[t]{0.235\linewidth}\poptile{poppink}{Exercices}{type examen + corrigés}\end{minipage}\hfill
  \begin{minipage}[t]{0.235\linewidth}\poptile{popgreen}{Fiche bilan}{l'essentiel à retenir}\end{minipage}\par}

% Sommaire
\newtcolorbox{sommaire}{enhanced,colback=white,colframe=popink,boxrule=2.2pt,arc=10pt,
  drop shadow=popink,left=18pt,right=18pt,top=13pt,bottom=13pt,before skip=6pt,after skip=20pt,
  before upper={\poppill{Sommaire}{poppurple}{white}\par\vspace{9pt}\popsemi\fontsize{10.6}{14.5}\selectfont\color{popink}}}

% Signature de fin de cours — poussée tout en bas de la dernière page
\newcommand{\findecours}{%
  \par\vfill\nopagebreak
  \begin{center}\popsquiggle{poporange}\end{center}\vspace{4pt}
  \signatureonbuch
}
\AtBeginDocument{\def\llbracket{\mathopen{[\mkern-3.5mu[}}\def\rrbracket{\mathclose{]\mkern-3.5mu]}}} % intervalles d'entiers (glyphe ⟦ absent de la police)
% Glyphes absents de Fira Math : redéfinis à partir de glyphes disponibles
\AtBeginDocument{%
  \def\setminus{\mathbin{\backslash}}\let\smallsetminus\setminus
  \def\vdots{\mathord{\vbox{\baselineskip3.2pt\lineskiplimit0pt\kern4pt\hbox{.}\hbox{.}\hbox{.}}}}%
  \def\ddots{\mathinner{\mkern1mu\raise6.5pt\hbox{.}\mkern2mu\raise3.6pt\hbox{.}\mkern2mu\raise0.7pt\hbox{.}\mkern1mu}}%
  \def\triangle{\mathord{\scalebox{0.9}{$\symup{\Delta}$}}}%
  \def\nabla{\mathord{\raisebox{\depth}{\scalebox{1}[-1]{$\symup{\Delta}$}}}}%
  \def\checkmark{\popcheck{popdark}}%
  \def\star{\mathbin{\ast}}\def\bigstar{\mathord{\ast}}%
  \def\diamond{\mathbin{\scalebox{0.6}{$\Diamond$}}}%
  \def\bigcup{\mathop{\mathchoice{\raisebox{-0.3em}{\scalebox{1.7}{$\cup$}}}{\scalebox{1.25}{$\cup$}}{\cup}{\cup}}\displaylimits}%
  \def\bigcap{\mathop{\mathchoice{\raisebox{-0.3em}{\scalebox{1.7}{$\cap$}}}{\scalebox{1.25}{$\cap$}}{\cap}{\cap}}\displaylimits}%
  \def\lesseqqgtr{\mathrel{\lessgtr}}\def\bigoplus{\mathop{\oplus}\displaylimits}%
}
\providecommand{\ssi}{si et seulement si} % abréviation fréquente
\AtBeginDocument{\providecommand{\wideparen}{\overparen}} % arc de cercle

%% FIGFIX-BEGIN v5 — correctifs globaux des figures (docs/onbuch-generation/tools/figfix)
\usepackage{adjustbox}\usepackage{etoolbox}\tcbuselibrary{raster}
% toute figure TikZ/pgfplots plus large que la ligne est réduite à la largeur disponible
\BeforeBeginEnvironment{tikzpicture}{\begin{adjustbox}{max width=\linewidth}}
\AfterEndEnvironment{tikzpicture}{\end{adjustbox}}
% légendes pgfplots lisibles par-dessus les courbes
\pgfplotsset{every axis/.append style={legend style={fill=white,fill opacity=0.92,text opacity=1,draw=black!15,inner sep=3pt}}}
% étiquettes d'arêtes (syntaxe edge["..."]) avec fond blanc
\tikzset{every edge quotes/.append style={fill=white,inner sep=1.2pt}}
%% FIGFIX-END
\begin{document}

\pagegarde{Lesson 27 — Vocabulary: Words and expressions related to keeping informed about health problems related to drugs, tobacco, alcohol abuse, and aging}{Anglais}{Tle A}{Chapter 3 : Environment, Well-being and Health}{EN27}{2 h}{Cette leçon de vocabulaire fournit le lexique thématique indispensable pour comprendre et discuter des problèmes de santé liés à la drogue, au tabac, à l'alcool et au vieillissement. Elle propose champs lexicaux, collocations, expressions idiomatiques, familles de mots et faux amis, avec des mises en contexte orales et écrites.
\vspace{4pt}
\begin{multicols}{2}\small
\begin{itemize}
\item À la fin de cette leçon, je sais identifier et définir le vocabulaire lié à la drogue, au tabac, à l'alcool et au vieillissement.
\item Je sais employer les collocations courantes du thème (to abuse drugs, to quit smoking, to die of...).
\item Je sais utiliser des expressions idiomatiques liées à la santé et à la prévention.
\item Je sais éviter les faux amis et les calques du français (drug, sensible, eventually...).
\item Je sais former des mots de la même famille (addict, addicted, addiction, addictive).
\item Je sais prononcer et orthographier correctement les mots clés du thème.
\end{itemize}\end{multicols}}

\begin{sommaire}
\begin{multicols}{2}
\begin{enumerate}
\item Drugs and drug abuse
\item Tobacco and smoking
\item Alcohol abuse
\item Aging and elderly health
\item Word families, pronunciation and spelling
\item Using the vocabulary in context
\item Activité d'intégration
\item Méthodes et exercices résolus
\item Exercices
\item Corrigés des exercices
\item Fiche bilan
\end{enumerate}
\end{multicols}
\end{sommaire}

\begin{objectifs}
\begin{itemize}
\item À la fin de cette leçon, je sais identifier et définir le vocabulaire lié à la drogue, au tabac, à l'alcool et au vieillissement.
\item Je sais employer les collocations courantes du thème (to abuse drugs, to quit smoking, to die of...).
\item Je sais utiliser des expressions idiomatiques liées à la santé et à la prévention.
\item Je sais éviter les faux amis et les calques du français (drug, sensible, eventually...).
\item Je sais former des mots de la même famille (addict, addicted, addiction, addictive).
\item Je sais prononcer et orthographier correctement les mots clés du thème.
\item Je sais réemployer ce lexique dans des phrases, des dialogues et de courts textes.
\item Je sais comprendre un article ou un message de sensibilisation sur ces problèmes de santé.
\end{itemize}
\end{objectifs}

\begin{prerequis}
\begin{itemize}
\item Vocabulaire général de la santé vu en Première (illness, disease, treatment, hospital).
\item Structures de conseil et d'obligation : should, must, have to, had better.
\item Les prépositions courantes après verbes et adjectifs (to suffer from, to be addicted to).
\item La formation des mots par préfixes et suffixes (health → healthy, care → careful).
\end{itemize}
\end{prerequis}

\newpage

\coursec{Drugs and drug abuse}

\courssub{Situation d'accroche et problématique}

Every morning in Yaoundé, Mrs Ngo Bell listens to a radio health programme warning young people about drug abuse, smoking and excessive drinking. Her son, a student in Buea, wants to discuss these issues in English with his classmates and to read international health reports from the WHO. But words like \eng{addiction}, \eng{withdrawal} or \eng{ageing} are not always easy to use correctly. How can we talk accurately about health risks linked to drugs, tobacco, alcohol and old age? This lesson gives you the essential vocabulary to stay informed and to inform others.

Dans cette première partie, nous allons construire pas à pas le \cle{champ lexical des drogues et de la toxicomanie} : les types de drogues, la dépendance, les conséquences sur la santé et la prévention. À la fin de la section, vous saurez lire un article de presse anglophone sur ce sujet et en parler sans calquer le français.

\courssub{Observation : un texte pour découvrir le vocabulaire}

Lisons d'abord un court texte de type article de presse. Soulignez mentalement tous les mots qui appartiennent au thème de la drogue : vous les retrouverez dans le glossaire.

\begin{textebox}[Drug abuse among young people: a growing concern]
In Cameroon, as in many countries, awareness campaigns warn young people against drug abuse. Health workers explain that illegal drugs such as cannabis, cocaine and heroin destroy lives. Some teenagers start taking drugs at parties because their friends pressure them, or because they want to forget their problems. Very quickly, they can get hooked: what began as a habit becomes an addiction.

A drug addict who wants to come off drugs often suffers from withdrawal symptoms and needs rehabilitation. In rehab centres, doctors and counsellors help patients to become drug-free again. Unfortunately, many families cannot afford this treatment, and some addicts suffer an overdose before they receive help.

The government has banned the sale of hard drugs, and the police regularly arrest dealers who sell drugs near schools. But trafficking continues, because it brings huge profits to criminal networks. Experts say that punishment alone is not enough: the best weapon against drug abuse is information. That is why schools, churches and radio stations organise awareness campaigns to fight this scourge and to protect the health of the nation.
\end{textebox}

\textbf{Glossaire}

\begin{center}
\begin{tabularx}{\linewidth}{|l|l|X|}
\hline
\rowcolor{popblueL} \textbf{Mot anglais} & \textbf{Nature} & \textbf{Traduction française} \\
\hline
awareness campaign & nom & campagne de sensibilisation \\
\hline
drug abuse & nom & abus de drogues, toxicomanie \\
\hline
illegal drugs & nom pluriel & drogues illicites \\
\hline
to get hooked & expression & devenir accro, devenir dépendant \\
\hline
addiction & nom & dépendance, addiction \\
\hline
to come off drugs & expression & arrêter la drogue, se sevrer \\
\hline
withdrawal symptoms & nom pluriel & symptômes de sevrage \\
\hline
rehabilitation / rehab & nom & cure de désintoxication \\
\hline
overdose & nom & surdose, overdose \\
\hline
dealer & nom & revendeur de drogue, dealer \\
\hline
trafficking & nom & trafic (de drogue) \\
\hline
to ban & verbe & interdire (officiellement) \\
\hline
drug-free & adjectif & sans drogue, libéré de la drogue \\
\hline
scourge & nom & fléau \\
\hline
\end{tabularx}
\end{center}

\textbf{Questions de compréhension}

\begin{enumerate}
\item Why do some teenagers start taking drugs, according to the text?
\item What happens to an addict who stops taking drugs?
\item What two measures does the government take against drugs?
\item According to experts, what is the best weapon against drug abuse?
\end{enumerate}

\textbf{Réponses attendues} : (1) Because their friends pressure them or because they want to forget their problems. (2) He suffers from withdrawal symptoms and needs rehabilitation. (3) It has banned the sale of hard drugs and the police arrest dealers. (4) Information.

\courssub{Le champ lexical des drogues : types et acteurs}

\begin{definition}[Drug, medicine, medication]
\cle{Drug} a deux sens en anglais : (1) un \emph{médicament} (substance qui soigne) ; (2) une \emph{drogue} (substance illicite qui modifie le comportement). \cle{Medicine} (GB) et \cle{medication} (USA, plus formel) signifient uniquement « médicament » : on ne peut jamais les employer pour une drogue illicite.
\end{definition}

Le contexte décide du sens de \eng{drug}. Comparez :

\begin{itemize}
\item \eng{The doctor prescribed a new drug for her asthma.} « Le médecin lui a prescrit un nouveau médicament pour son asthme. » (sens : médicament)
\item \eng{Cocaine is a dangerous drug.} « La cocaïne est une drogue dangereuse. » (sens : drogue illicite)
\end{itemize}

\begin{attention}[Drug ne signifie pas seulement « drogue »]
Dans \eng{drugstore} (USA) ou \eng{prescription drug} (médicament délivré sur ordonnance), il s'agit de \emph{médicaments}, pas de stupéfiants. De même, \eng{drugstore} ne se traduit pas par « magasin de drogue » mais par « drugstore » (grande pharmacie américaine). Ne traduisez jamais mot à mot : lisez la phrase entière.
\end{attention}

\begin{definition}[Types de drogues et acteurs du trafic]
\cle{Illegal drugs} (drogues illicites) se divisent traditionnellement en \cle{hard drugs} (drogues dures : \eng{heroin}, \eng{cocaine}) et \cle{soft drugs} (drogues douces : \eng{cannabis}, aussi appelé \eng{marijuana} ou, familièrement, \eng{weed}). Celui qui vend de la drogue est un \cle{dealer} (ou \cle{drug dealer}) ; l'activité s'appelle \cle{drug trafficking} (trafic de drogue) et celui qui l'organise est un \cle{trafficker}. Celui qui consomme est un \cle{drug user} (consommateur de drogue).
\end{definition}

\begin{exemplebox}[Exemples traduits]
\eng{He was arrested for dealing drugs.} « Il a été arrêté pour vente de drogue. »

\eng{Drug trafficking is severely punished by law.} « Le trafic de drogue est sévèrement puni par la loi. »

\eng{The police seized a large quantity of hard drugs at the port of Douala.} « La police a saisi une grande quantité de drogues dures au port de Douala. »

\eng{This medicine is only available on prescription.} « Ce médicament n'est délivré que sur ordonnance. »
\end{exemplebox}

\begin{attention}[Heroin ou heroine ?]
Ne confondez pas \eng{heroin} (sans \textbf{e} final : l'héroïne, la drogue) et \eng{heroine} (avec \textbf{e} final : l'héroïne d'un roman ou d'un film). Les deux se prononcent de la même façon (« hé-ro-ine »), mais l'orthographe change tout le sens : \eng{She is the heroine of the novel.} « Elle est l'héroïne du roman. »
\end{attention}

\begin{savaistu}[Chemist's ou drugstore ?]
Au Royaume-Uni, la pharmacie s'appelle le \eng{chemist's} (ou \eng{pharmacy}) et le pharmacien est le \eng{chemist}. Aux États-Unis, on dit \eng{drugstore} : on y achète des médicaments, mais aussi des produits de beauté, des journaux et des snacks. Attention donc : un Américain qui dit \eng{I'm going to the drugstore} ne va pas acheter de la drogue !
\end{savaistu}

\courssub{La dépendance : addiction, addict, dependence, habit}

\begin{definition}[Addiction et ses dérivés]
\cle{Addiction} (nom, prononcé « a-dik-cheune ») signifie « dépendance, addiction ». La personne dépendante est un(e) \cle{addict} (nom, prononcé « a-dikte » : un toxicomane). L'adjectif est \cle{addicted} : on dit \cle{to be addicted to} + nom ou verbe en -ing (être dépendant de). Le nom composé \cle{drug addiction} signifie « la toxicomanie ».
\end{definition}

\begin{propriete}[Construction de « addicted » et « dependent »]
\begin{itemize}
\item \eng{to be addicted \textbf{to} something} : \eng{He is addicted to cocaine.} « Il est dépendant de la cocaïne. »
\item \eng{to be dependent \textbf{on} something} : \eng{She is dependent on sleeping pills.} « Elle est dépendante des somnifères. »
\item \eng{dependence \textbf{on}} (nom) : \eng{his dependence on alcohol} « sa dépendance à l'alcool ».
\item \eng{drug addiction} (nom composé) : « la toxicomanie ».
\end{itemize}
Remarquez la différence de préposition : \textbf{to} après \eng{addicted}, \textbf{on} après \eng{dependent/dependence}.
\end{propriete}

\begin{attention}[Habit n'est pas une addiction]
\eng{Habit} désigne une habitude (bonne ou mauvaise) : \eng{Smoking is a bad habit.} « Fumer est une mauvaise habitude. » Une \eng{drug habit} (registre familier) signifie « une consommation régulière de drogue ». Mais une habitude n'est pas forcément une addiction : \eng{addiction} implique une dépendance physique ou psychologique. Ne confondez pas non plus \eng{habit} (habitude) et \eng{habitat} (habitat, milieu de vie).
\end{attention}

\begin{exemplebox}[Exemples traduits]
\eng{Many young people get addicted to drugs without realising it.} « Beaucoup de jeunes deviennent dépendants de la drogue sans s'en rendre compte. »

\eng{His brother is a drug addict; he needs help, not punishment.} « Son frère est toxicomane ; il a besoin d'aide, pas de punition. »

\eng{Drug addiction destroys families.} « La toxicomanie détruit les familles. »
\end{exemplebox}

\courssub{Les conséquences : overdose, withdrawal, side effects, rehabilitation}

\begin{definition}[Withdrawal symptoms]
\cle{Withdrawal symptoms} (prononcé « ouiz-drô-ol sim-tomz ») : the unpleasant physical and mental effects felt when someone stops taking a drug they are dependent on — les effets physiques et mentaux désagréables ressentis quand une personne dépendante arrête de consommer une drogue (symptômes de sevrage : tremblements, nausées, anxiété, insomnie).
\end{definition}

\begin{definition}[Overdose, side effects, rehabilitation]
\cle{Overdose} (nom et verbe, souvent abrégé \eng{OD} en langage familier) : surdose, prise d'une quantité dangereuse d'une drogue — \eng{to die of an overdose} « mourir d'une overdose ». \cle{Side effects} : effets secondaires (d'un médicament). \cle{Rehabilitation} (abrégé \cle{rehab}, « ri-hab ») : cure de désintoxication ; \eng{to go into rehab} « entrer en cure de désintoxication ». L'expression \eng{to ruin one's health} signifie « se ruiner la santé ».
\end{definition}

\begin{exemplebox}[Exemples traduits]
\eng{She went into rehab last year.} « Elle est entrée en cure de désintoxication l'année dernière. »

\eng{When he stopped taking the drug, he suffered from terrible withdrawal symptoms.} « Quand il a arrêté la drogue, il a souffert de terribles symptômes de sevrage. »

\eng{This medicine can cause side effects such as headaches.} « Ce médicament peut provoquer des effets secondaires comme des maux de tête. »

\eng{Years of drug abuse ruined his health.} « Des années d'abus de drogues ont ruiné sa santé. »
\end{exemplebox}

\courssub{Prévention et lutte contre la drogue}

\begin{definition}[Awareness campaign, to fight, to ban, drug-free]
Une \cle{awareness campaign} est une campagne de sensibilisation (\eng{to raise awareness about drugs} = « sensibiliser aux dangers de la drogue »). \cle{To fight / to combat drug abuse} signifie « lutter contre la toxicomanie ». \cle{To ban} (verbe, participe passé \eng{banned}) signifie « interdire officiellement » ; le nom est \eng{a ban} : \eng{a ban on drug advertising} « une interdiction de la publicité pour les drogues ». \cle{Drug-free} (adjectif) signifie « sans drogue » : \eng{a drug-free school} « une école sans drogue ».
\end{definition}

\begin{exemplebox}[Exemples traduits]
\eng{The government has launched an awareness campaign in secondary schools.} « Le gouvernement a lancé une campagne de sensibilisation dans les établissements secondaires. »

\eng{Selling drugs near schools is banned.} « La vente de drogue près des écoles est interdite. »

\eng{After two years in rehab, he is now completely drug-free.} « Après deux ans de cure, il est maintenant complètement libéré de la drogue. »
\end{exemplebox}

\courssub{Les collocations essentielles}

En anglais, on ne peut pas combiner n'importe quel verbe avec n'importe quel nom : ce sont les \cle{collocations}. Voici les verbes qui « vont avec » \eng{drugs}.

\begin{center}
\begin{tabularx}{\linewidth}{|X|X|X|}
\hline
\rowcolor{popblueL} \textbf{Collocation anglaise} & \textbf{Traduction française} & \textbf{Exemple} \\
\hline
to take drugs & prendre de la drogue & \eng{He started taking drugs at sixteen.} \\
\hline
to use drugs & consommer de la drogue & \eng{She has never used drugs.} \\
\hline
to abuse drugs & abuser de la drogue & \eng{People who abuse drugs need treatment.} \\
\hline
to deal drugs & vendre de la drogue & \eng{He was caught dealing drugs.} \\
\hline
to get hooked on drugs & devenir accro à la drogue & \eng{Many users get hooked very quickly.} \\
\hline
to come off drugs & arrêter la drogue & \eng{It is hard to come off drugs alone.} \\
\hline
to go into rehab & entrer en désintoxication & \eng{She went into rehab last year.} \\
\hline
to fight drug abuse & lutter contre la toxicomanie & \eng{Schools must fight drug abuse.} \\
\hline
\end{tabularx}
\end{center}

\begin{attention}[Pièges fréquents des francophones]
\begin{itemize}
\item Ne dites pas \eng{*to consume drugs} (calque de « consommer ») : dites \eng{to take drugs} ou \eng{to use drugs}.
\item Ne dites pas \eng{*a drugged person} pour « un drogué » : \eng{drugged} signifie « drogué (à son insu) » ; dites \eng{a drug addict} ou \eng{a drug user}.
\item \eng{To stop drugs} est incorrect : dites \eng{to stop taking drugs} ou \eng{to come off drugs}.
\item Le nom s'écrit \eng{dependence} avec \textbf{-ence} en anglais britannique comme en anglais américain ; l'adjectif est \eng{dependent} avec \textbf{-ent}. N'écrivez jamais \eng{*dependance} ni \eng{*dependant} (formes françaises).
\end{itemize}
\end{attention}

\courssub{Carte mentale : le vocabulaire de la drogue}

\begin{popfigure}
\begin{tikzpicture}[mindmap, grow cyclic,
  every node/.style={concept, font=\small},
  root concept/.append style={concept color=popblue, font=\bfseries},
  level 1/.append style={level distance=3.2cm, sibling angle=90},
  level 2/.append style={level distance=2.4cm, sibling angle=45, font=\scriptsize}]
\node[root concept]{DRUGS}
  child[concept color=poporange]{ node {Types}
    child { node[concept color=poporangeL, text=popdark] {hard / soft drugs} }
    child { node[concept color=poporangeL, text=popdark] {medicine, medication} }
  }
  child[concept color=poppurple]{ node {Addiction}
    child { node[concept color=poppurpleL, text=popdark] {addict, addicted to} }
    child { node[concept color=poppurpleL, text=popdark] {dependence, habit} }
    child { node[concept color=poppurpleL, text=popdark] {get hooked on} }
  }
  child[concept color=poppink]{ node {Consequences}
    child { node[concept color=poppinkL, text=popdark] {overdose} }
    child { node[concept color=poppinkL, text=popdark] {withdrawal symptoms} }
    child { node[concept color=poppinkL, text=popdark] {side effects, rehab} }
  }
  child[concept color=popgreen]{ node {Prevention}
    child { node[concept color=popgreenL, text=popdark] {awareness campaign} }
    child { node[concept color=popgreenL, text=popdark] {to ban, drug-free} }
    child { node[concept color=popgreenL, text=popdark] {fight drug abuse} }
  };
\end{tikzpicture}
\legende{Carte mentale du champ lexical de la drogue : quatre branches à mémoriser.}
\end{popfigure}

\courssub{Exercice résolu et entraînement}

\begin{exoresolu}[Complétez avec le mot qui convient]
Complétez chaque phrase avec le mot ou l'expression approprié(e) : \eng{addicted}, \eng{dealer}, \eng{withdrawal symptoms}, \eng{awareness campaign}, \eng{overdose}, \eng{rehab}, \eng{banned}.

1. The police arrested a \ldots\ldots\ldots who was selling drugs near the lycée.

2. After six months, he became \ldots\ldots\ldots to heroin.

3. When she stopped taking the drug, she had terrible \ldots\ldots\ldots.

4. The school organised an \ldots\ldots\ldots to inform students about the dangers of drugs.

5. He nearly died of an \ldots\ldots\ldots last year.

6. Her parents sent her to \ldots\ldots\ldots for three months.

7. The sale of this drug is \ldots\ldots\ldots in Cameroon.

\tcblower
\textbf{Solution détaillée}

1. \textbf{dealer} — celui qui vend de la drogue est un \eng{dealer} (ou \eng{drug dealer}).

2. \textbf{addicted} — construction : \eng{to be/become addicted \textbf{to} + nom} ; attention à la préposition \eng{to}.

3. \textbf{withdrawal symptoms} — les effets ressentis quand on arrête une drogue dont on est dépendant.

4. \textbf{awareness campaign} — une campagne de sensibilisation ; collocation : \eng{to organise / to launch an awareness campaign}.

5. \textbf{overdose} — une surdose ; collocation : \eng{to die of an overdose}.

6. \textbf{rehab} — abréviation courante de \eng{rehabilitation} ; collocation : \eng{to go into rehab}.

7. \textbf{banned} — participe passé de \eng{to ban} (interdire) ; attention à la consonne doublée devant une terminaison vocalique : \eng{ban} $\rightarrow$ \eng{banned}, \eng{banning}.
\end{exoresolu}

\begin{exercice}[À vous de jouer]
1. Traduisez en anglais : « Il est devenu accro à la drogue à dix-sept ans. » / « Le gouvernement lutte contre le trafic de drogue. » / « Elle suit une cure de désintoxication depuis janvier. »

2. Trouvez l'intrus dans chaque série et justifiez : (a) \eng{medicine / medication / heroin} ; (b) \eng{dealer / trafficker / patient}.

3. Rédigez trois phrases sur la lutte contre la drogue dans votre établissement en utilisant \eng{awareness campaign}, \eng{to ban} et \eng{drug-free}.
\end{exercice}

\begin{corrige}[Exercice « À vous de jouer »]
1. \eng{He got hooked on drugs at seventeen.} / \eng{The government is fighting drug trafficking.} / \eng{She has been in rehab since January.} (present perfect avec \eng{since} : l'action continue.)

2. (a) \eng{heroin} : c'est une drogue illicite ; \eng{medicine} et \eng{medication} désignent des médicaments. (b) \eng{patient} : c'est une personne soignée ; \eng{dealer} et \eng{trafficker} participent au trafic.

3. Exemple de production : \eng{Our school has launched an awareness campaign against drugs. Smoking and drug use are banned inside the school. We want our lycée to be a drug-free environment.}
\end{corrige}

\begin{aretenir}[L'essentiel de la section]
\begin{itemize}
\item \eng{Drug} = médicament \emph{ou} drogue selon le contexte ; \eng{medicine/medication} = médicament uniquement.
\item Dépendance : \eng{addiction} (nom), \eng{addict} (personne), \eng{to be addicted \textbf{to}} ; \eng{dependent \textbf{on}} / \eng{dependence \textbf{on}}.
\item Conséquences : \eng{overdose} (surdose), \eng{withdrawal symptoms} (symptômes de sevrage), \eng{side effects} (effets secondaires), \eng{rehab} (désintoxication).
\item Prévention : \eng{awareness campaign}, \eng{to fight drug abuse}, \eng{to ban} (\eng{banned}, \eng{banning}), \eng{drug-free}.
\item Collocations à connaître par cœur : \eng{to take/use/abuse drugs}, \eng{to deal drugs}, \eng{to get hooked on}, \eng{to come off drugs}, \eng{to go into rehab}.
\end{itemize}
\end{aretenir}

\coursec{Tobacco and smoking}

Après avoir étudié le vocabulaire de la drogue et de la toxicomanie (\eng{drugs and drug abuse}), nous abordons un autre problème de santé publique majeur, présent partout autour de nous, de Yaoundé à Douala : le \cle{tabac}. Contrairement aux drogues illégales, le tabac est un produit légal, vendu librement, mais il tue chaque année des millions de personnes dans le monde. Le vocabulaire anglais du tabagisme est très riche et très fréquent dans la presse anglophone, les campagnes de santé publique et les conversations quotidiennes. C'est donc un champ lexical indispensable pour l'élève de Terminale, tant pour la compréhension écrite que pour la composition au Baccalauréat.

\courssub{Observation : a short text to discover the vocabulary}

\begin{textebox}[A personal story — “My uncle the chain smoker”]
My uncle was a chain smoker for twenty years. He used to smoke two packets of cigarettes a day, and he would light a cigarette as soon as he woke up. After his doctor had warned him about lung cancer and serious respiratory diseases, he decided to kick the habit. It was hard: nicotine is very addictive, and the tar in cigarettes had already damaged his lungs. He coughed all the time. First, he tried to cut down on cigarettes; then he gave up smoking completely. It was difficult, but now he lives in a smoke-free home, and his children are no longer exposed to second-hand smoke. He says quitting was the best decision of his life.
\end{textebox}

\textbf{Glossaire (mots difficiles) :}

\begin{center}
\begin{tabularx}{\linewidth}{|l|l|X|}
\hline
\rowcolor{popblueL} \textbf{English} & \textbf{Nature} & \textbf{Français} \\
\hline
chain smoker & nom & gros fumeur (qui enchaîne les cigarettes) \\
\hline
packet of cigarettes & nom & paquet de cigarettes \\
\hline
to light a cigarette & verbe & allumer une cigarette \\
\hline
to warn & verbe & avertir, mettre en garde \\
\hline
lung cancer & nom & cancer du poumon \\
\hline
to kick the habit & expression & se débarrasser d'une (mauvaise) habitude \\
\hline
addictive & adjectif & qui crée une dépendance \\
\hline
tar & nom & goudron \\
\hline
to damage & verbe & endommager, abîmer \\
\hline
to cut down (on) & verbe & réduire (sa consommation de) \\
\hline
to give up smoking & verbe & arrêter de fumer \\
\hline
smoke-free & adjectif & sans tabac, non-fumeur \\
\hline
second-hand smoke & nom & fumée secondaire (tabagisme passif) \\
\hline
to quit & verbe & arrêter, cesser \\
\hline
\end{tabularx}
\end{center}

\textbf{Comprehension questions :}
\begin{enumerate}
\item How long did the narrator's uncle smoke?
\item What did the doctor warn him about?
\item What two steps did he take before quitting completely?
\item Why is his home healthier now?
\end{enumerate}

\textbf{Réponses attendues :} 1. \eng{He smoked for twenty years.} 2. \eng{The doctor warned him about lung cancer and serious respiratory diseases.} 3. \eng{First he cut down on cigarettes, then he gave up smoking completely.} 4. \eng{Because it is now smoke-free and his children are no longer exposed to second-hand smoke.}

\courssub{The lexical field of tobacco}

\begin{definition}[Tobacco and smoking — le champ lexical]
Le champ lexical du \cle{tabac} comprend les produits (\eng{tobacco}, \eng{cigarette}, \eng{cigar}, \eng{pipe}), les personnes (\eng{smoker}, \eng{non-smoker}), les actions (\eng{to smoke}, \eng{to light}, \eng{to puff}, \eng{to inhale}) et les phénomènes sociaux (\eng{passive smoking}, \eng{smoking ban}). En anglais, \eng{tobacco} est \textbf{indénombrable} : on dit \eng{some tobacco}, jamais \eng{a tobacco}.
\end{definition}

\begin{center}
\begin{tabularx}{\linewidth}{|l|X|X|}
\hline
\rowcolor{popblueL} \textbf{English} & \textbf{Français} & \textbf{Exemple en contexte} \\
\hline
tobacco & tabac (indénombrable) & \eng{Tobacco is sold in many shops.} \\
\hline
cigarette & cigarette & \eng{He lit a cigarette.} \\
\hline
cigar & cigare & \eng{My grandfather smokes a cigar on Sundays.} \\
\hline
pipe & pipe & \eng{The old man was smoking a pipe.} \\
\hline
smoker & fumeur & \eng{Smokers pay more for health insurance.} \\
\hline
heavy smoker & gros fumeur & \eng{He is a heavy smoker: a packet a day.} \\
\hline
non-smoker & non-fumeur & \eng{Non-smokers have the right to clean air.} \\
\hline
passive / second-hand smoking & tabagisme passif & \eng{Passive smoking harms children.} \\
\hline
\end{tabularx}
\end{center}

\begin{exemplebox}[Exemples traduits]
\eng{Passive smoking harms children.} « Le tabagisme passif nuit aux enfants. »

\eng{She is trying to cut down on cigarettes.} « Elle essaie de réduire sa consommation de cigarettes. »

\eng{He is a heavy smoker: he smokes a packet a day.} « C'est un gros fumeur : il fume un paquet par jour. »

\eng{Many workers in Douala are exposed to second-hand smoke in bars.} « De nombreux employés à Douala sont exposés à la fumée des autres dans les bars. »
\end{exemplebox}

\begin{definition}[Passive smoking]
\eng{Passive smoking} (ou \eng{second-hand smoking}) means \emph{breathing in other people's tobacco smoke}, which can also cause serious diseases such as lung cancer and asthma. En français : le \cle{tabagisme passif}. On dit \eng{to be exposed to second-hand smoke} = « être exposé à la fumée des autres ».
\end{definition}

\begin{attention}[“Packet” ou “pack” ?]
En anglais britannique, on dit \eng{a packet of cigarettes} ; \eng{a pack of cigarettes} est la forme américaine, très répandue aussi. Les deux sont correctes, mais au Baccalauréat, privilégiez la forme britannique \eng{packet}. En revanche, ne dites jamais \eng{a box of cigarettes} (calque du français « une boîte »).
\end{attention}

\courssub{Actions and verbs: to smoke, to quit, to cut down}

\begin{center}
\begin{tabularx}{\linewidth}{|l|X|X|}
\hline
\rowcolor{popblueL} \textbf{Verbe anglais} & \textbf{Français} & \textbf{Exemple} \\
\hline
to smoke & fumer & \eng{He smokes twenty cigarettes a day.} \\
\hline
to light a cigarette & allumer une cigarette & \eng{She lit a cigarette nervously.} \\
\hline
to puff (on) & tirer (des bouffées) & \eng{He puffed on his pipe silently.} \\
\hline
to inhale & inhaler & \eng{Do not inhale the smoke deeply.} \\
\hline
to quit / give up smoking & arrêter de fumer & \eng{He quit smoking last year.} \\
\hline
to cut down (on) & réduire & \eng{She is cutting down on cigarettes.} \\
\hline
\end{tabularx}
\end{center}

\begin{propriete}[“To give up” + nom ou verbe en -ing]
Le verbe à particule \cle{to give up} (« renoncer à, arrêter ») est suivi d'un \textbf{nom} ou d'un \textbf{verbe en -ing}, jamais d'un infinitif.

\eng{He gave up smoking.} « Il a arrêté de fumer. » (jamais \eng{gave up to smoke})

\eng{She gave up cigarettes.} « Elle a renoncé aux cigarettes. »

Même structure avec \eng{to stop} : \eng{He stopped smoking.} « Il a arrêté de fumer. »

Conjugaison : \eng{give — gave — given} ; \eng{quit — quit — quit} (verbe invariant) ; \eng{light — lit — lit} (ou \eng{lighted}).
\end{propriete}

\begin{attention}[“Stop smoking” ou “stop to smoke” ?]
Attention au piège classique ! \eng{To stop smoking} signifie « arrêter de fumer » (on cesse l'habitude), tandis que \eng{to stop to smoke} signifie « s'arrêter \textbf{pour} fumer » (on interrompt une activité afin de fumer).

\eng{He stopped smoking.} « Il a arrêté de fumer. »

\eng{He stopped to smoke.} « Il s'est arrêté pour fumer. »

Ne traduisez donc jamais « arrêter de fumer » par \eng{stop to smoke} !
\end{attention}

\courssub{Health effects and regulation}

\textbf{Effets sur la santé :} \eng{lung cancer} (cancer du poumon), \eng{respiratory diseases} (maladies respiratoires), \eng{nicotine} (nicotine, substance qui crée la dépendance), \eng{tar} (goudron), \eng{coughing} (toux), \eng{to damage one's lungs} (s'abîmer les poumons).

\textbf{Réglementation :} \eng{smoking ban} (interdiction de fumer), \eng{smoke-free area} (zone non-fumeurs), \eng{warning label} (étiquette d'avertissement), le slogan célèbre \eng{“Smoking kills”} (« Fumer tue »), et \eng{to raise taxes on tobacco} (augmenter les taxes sur le tabac).

\begin{exemplebox}[Exemples traduits]
\eng{The government has raised taxes on tobacco.} « Le gouvernement a augmenté les taxes sur le tabac. »

\eng{Every packet carries a warning label: “Smoking kills”.} « Chaque paquet porte un avertissement : “Fumer tue”. »

\eng{Restaurants in Buea are now smoke-free areas.} « Les restaurants de Buea sont désormais des zones non-fumeurs. »

\eng{Years of smoking have damaged his lungs.} « Des années de tabagisme ont abîmé ses poumons. »
\end{exemplebox}

\begin{popfigure}
\begin{center}
\begin{tikzpicture}[node distance=0.6cm, every node/.style={font=\small}]
\node[draw=popblue, fill=popblueL, rounded corners, align=center, minimum width=2.6cm] (smoke) {smoking};
\node[draw=poporange, fill=poporangeL, rounded corners, align=center, minimum width=2.6cm, right=of smoke] (addict) {nicotine\\addiction};
\node[draw=poppurple, fill=poppurpleL, rounded corners, align=center, minimum width=2.6cm, right=of addict] (damage) {lung\\damage};
\node[draw=poppink, fill=poppinkL, rounded corners, align=center, minimum width=2.6cm, right=of damage] (cancer) {cancer};
\draw[-{Stealth}, thick, popdark] (smoke) -- (addict);
\draw[-{Stealth}, thick, popdark] (addict) -- (damage);
\draw[-{Stealth}, thick, popdark] (damage) -- (cancer);
\node[draw=popgreen, fill=popgreenL, rounded corners, align=center, below=0.9cm of addict] (solutions) {solutions: quit / rehab / smoking ban};
\draw[-{Stealth}, thick, popgreen, dashed] (solutions) -- (smoke.south);
\draw[-{Stealth}, thick, popgreen, dashed] (solutions.east) -- (damage.south);
\end{tikzpicture}
\end{center}
\legende{La chaîne des effets du tabac et les solutions possibles : arrêter (\eng{quit}), se faire soigner (\eng{rehab}), interdire (\eng{ban}).}
\end{popfigure}

\begin{savaistu}[“To kick the habit”]
L'expression idiomatique \eng{to kick the habit} signifie littéralement « donner un coup de pied à l'habitude », c'est-à-dire \emph{se débarrasser d'une (mauvaise) habitude}. Elle s'emploie pour le tabac, l'alcool ou la drogue : \eng{After twenty years, he finally kicked the habit.} « Après vingt ans, il s'est enfin débarrassé de cette habitude. »
\end{savaistu}

\begin{exoresolu}[Complétez avec le vocabulaire du tabac]
Énoncé — Complétez chaque phrase avec le mot ou l'expression qui convient : \eng{chain smoker}, \eng{second-hand smoke}, \eng{gave up}, \eng{warning label}, \eng{cut down}.

1. My father used to be a \ldots\ : he smoked three packets a day.

2. Children who breathe \ldots\ can develop asthma.

3. She \ldots\ smoking when she became pregnant.

4. Every cigarette packet must carry a \ldots\ .

5. He cannot stop completely, so he is trying to \ldots\ on cigarettes.

\tcblower
Solution détaillée :

1. \eng{chain smoker} — « fumeur à la chaîne », collocation exacte pour quelqu'un qui enchaîne les cigarettes.

2. \eng{second-hand smoke} — la fumée respirée par les non-fumeurs (tabagisme passif).

3. \eng{gave up} — \eng{to give up smoking} = arrêter de fumer ; ici au prétérit : \eng{gave up}.

4. \eng{warning label} — « étiquette d'avertissement », collocation obligatoire dans ce contexte.

5. \eng{cut down} — \eng{to cut down on cigarettes} = réduire sa consommation ; attention à la préposition \eng{on}.
\end{exoresolu}

\begin{attention}[Erreurs fréquentes des francophones]
\begin{itemize}
\item Ne dites pas \eng{a tobacco} : \eng{tobacco} est indénombrable ; dites \eng{some tobacco} ou \eng{a packet of cigarettes}.
\item Ne traduisez pas « un fumeur » par \eng{a smoking man} : le nom est \eng{a smoker}.
\item Attention à la prononciation de \eng{cough} : « kof » (le \eng{-ough} se prononce « f »), et de \eng{nicotine} : « ni-ko-tine » (accent sur la première syllabe).
\item \eng{To cut down} exige la préposition \eng{on} : \eng{cut down on cigarettes}, jamais \eng{cut down cigarettes}.
\item Ne confondez pas \eng{smoke} (la fumée, indénombrable) et \eng{a smoke} qui, au sens familier, désigne « une cigarette » : \eng{He went out for a smoke.}
\end{itemize}
\end{attention}

\begin{exercice}[À vous de jouer]
1. Traduisez en anglais : « Il a arrêté de fumer il y a deux ans. » / « Le tabagisme passif est dangereux pour les bébés. » / « Elle essaie de réduire sa consommation de cigarettes. »

2. Rédigez un court paragraphe (5 lignes) intitulé \eng{Why people should kick the habit}, en utilisant au moins quatre mots du vocabulaire de cette leçon.
\end{exercice}

\begin{corrige}[Exercice — traductions]
1. \eng{He gave up smoking two years ago.} / \eng{Passive smoking is dangerous for babies.} / \eng{She is trying to cut down on cigarettes.}

2. Modèle : \eng{People should kick the habit because smoking damages their lungs and can cause lung cancer. Nicotine is very addictive, so it is hard to quit. Governments should raise taxes on tobacco and create more smoke-free areas to protect non-smokers from second-hand smoke.}
\end{corrige}

\coursec{Alcohol abuse}

Après avoir étudié le tabac et ses méfaits, nous abordons maintenant un autre fléau de santé publique : l'abus d'alcool. Au Cameroun comme dans les pays anglophones, la consommation excessive d'alcool est à l'origine de maladies graves, d'accidents de la route et de drames familiaux. Le vocabulaire de ce domaine est très riche en anglais, et il comporte plusieurs pièges pour les francophones : faux amis (\emph{spirits}, \emph{alcoholic}), adjectifs à double emploi (\emph{drunk} / \emph{drunken}) et expressions idiomatiques (\emph{to be on the wagon}). Observons d'abord un texte, puis construisons méthodiquement tout le champ lexical.

\courssub{Observation : un texte de référence}

\begin{textebox}[Health warning: alcohol and young people]
Binge drinking among teenagers is a growing concern. Drinking too much alcohol can damage the liver and lead to alcoholism. Doctors advise adults to drink in moderation or to stay sober. In many countries, the law bans the sale of alcohol to minors, and the legal drinking age is strictly controlled. Heavy drinking also causes drink-driving accidents: a driver who is under the influence of alcohol reacts slowly and puts lives in danger. After a night of drunken parties, many young people wake up with a terrible hangover: headache, nausea and tiredness. Specialists warn that a person who has a drinking problem should seek help early. Some former drinkers decide to go on the wagon and never touch beer, wine or spirits again. Their families are proud of them, because staying sober is not easy. Health campaigns on the radio, in schools and in churches remind everyone that alcohol abuse destroys the body, the mind and the family.
\end{textebox}

\textbf{Glossaire}

\begin{center}
\begin{tabularx}{\linewidth}{|l|l|X|}
\hline
\rowcolor{popblueL} English & Nature & Français \\
\hline
a concern & nom & une préoccupation, un sujet d'inquiétude \\
\hline
to damage & verbe & endommager, abîmer \\
\hline
the liver & nom & le foie \\
\hline
to ban & verbe & interdire \\
\hline
a minor & nom & un mineur \\
\hline
nausea & nom & la nausée \\
\hline
to seek help & expression & chercher de l'aide \\
\hline
a former drinker & nom & un ancien buveur \\
\hline
\end{tabularx}
\end{center}

\textbf{Questions de compréhension}
\begin{enumerate}
\item What health problems can heavy drinking cause?
\item What do doctors advise adults to do?
\item What does the law say about the sale of alcohol to minors?
\item What are the symptoms of a hangover?
\item What does it mean “to go on the wagon”?
\end{enumerate}

\begin{corrige}[Questions de compréhension]
\begin{enumerate}
\item It can damage the liver and lead to alcoholism.
\item They advise them to drink in moderation or to stay sober.
\item The law bans the sale of alcohol to minors.
\item Headache, nausea and tiredness.
\item It means to stop drinking alcohol completely.
\end{enumerate}
\end{corrige}

\courssub{Le champ lexical de l'alcool et des boissons}

\begin{definition}[Alcohol and alcoholic drinks]
\cle{Alcohol} (prononcé « al-keu-rol ») désigne l'alcool en général. Une boisson alcoolisée se dit \cle{an alcoholic drink} ou \cle{an alcoholic beverage} (terme plus soutenu). Les principales catégories sont : \cle{beer} (la bière), \cle{wine} (le vin) et \cle{spirits} (les spiritueux, les alcools forts comme le whisky ou la vodka). Une personne qui boit trop est \cle{a drunkard} (un ivrogne, péjoratif) ou \cle{an alcoholic} (un alcoolique, terme médical).
\end{definition}

\begin{attention}[Faux amis : spirits et alcoholic]
Deux pièges majeurs pour les francophones :
\begin{itemize}
\item \eng{Spirits} (au pluriel) ne signifie pas « les esprits » mais \textbf{les boissons fortement alcoolisées} (whisky, gin, vodka). Pour « l'esprit » humain, on dit \emph{mind} ou \emph{spirit} au singulier ; le sens « esprits » (fantômes) n'apparaît que dans un contexte religieux ou fantastique.
\item \eng{Alcoholic} est à la fois un adjectif (\eng{an alcoholic drink}) et un nom (\eng{He is an alcoholic} = C'est un alcoolique). On ne dit JAMAIS \eng{He is an alcohol} : \emph{alcohol} désigne la substance, pas la personne.
\end{itemize}
\end{attention}

\begin{exemplebox}[Exemples traduits]
\eng{Beer and palm wine are common drinks in many Cameroonian villages.} « La bière et le vin de palme sont des boissons courantes dans de nombreux villages camerounais. »

\eng{Spirits contain more alcohol than beer.} « Les spiritueux contiennent plus d'alcool que la bière. »

\eng{He is an alcoholic and he needs medical help.} « C'est un alcoolique et il a besoin d'une aide médicale. »

\eng{The old drunkard was sleeping near the bar.} « Le vieil ivrogne dormait près du bar. »
\end{exemplebox}

\courssub{États et actions : to drink, to get drunk, hangover}

\begin{definition}[États liés à l'ivresse]
\cle{To drink} (drank, drunk) = boire. \cle{To get drunk} = se saouler, devenir ivre. \cle{Drunk} = ivre, saoul. \cle{To be under the influence of alcohol} = être sous l'influence de l'alcool (expression officielle, policière). \cle{A hangover} (prononcé « han-gô-veur ») = la gueule de bois, le mal de tête du lendemain de beuverie.
\end{definition}

\begin{propriete}[La distinction drunk / drunken]
Les deux adjectifs viennent du participe passé de \emph{drink}, mais ils ne s'emploient pas au même endroit :
\begin{itemize}
\item \cle{Drunk} s'emploie \textbf{après le verbe} (attribut) : \eng{He is drunk.} « Il est ivre. »
\item \cle{Drunken} s'emploie \textbf{devant le nom} (épithète) : \eng{a drunken man} « un homme ivre », \eng{drunken behaviour} « un comportement d'ivrogne ».
\end{itemize}
Exception d'usage : dans la langue parlée américaine, on entend souvent \eng{drunk driving}, mais la forme britannique correcte est \eng{drink-driving}.
\end{propriete}

\begin{exemplebox}[Exemples traduits]
\eng{The driver was drunk and could not control his car.} « Le chauffeur était ivre et ne pouvait pas contrôler sa voiture. »

\eng{A drunken passenger insulted the conductor on the bus to Douala.} « Un passager ivre a insulté le contrôleur dans le bus pour Douala. »

\eng{After the party, she had a terrible hangover.} « Après la fête, elle avait une terrible gueule de bois. »

\eng{He was arrested because he was under the influence of alcohol.} « Il a été arrêté parce qu'il était sous l'influence de l'alcool. »
\end{exemplebox}

\courssub{Abus et conséquences : alcoholism, binge drinking, cirrhosis}

\begin{definition}[Binge drinking]
\cle{Binge drinking} (prononcé « bindj drin-king ») signifie \textbf{boire une grande quantité d'alcool en peu de temps} : c'est la beuverie, la consommation excessive ponctuelle, très répandue chez les adolescents et les étudiants anglophones. On dit aussi \eng{to go on a binge} (faire une beuverie).
\end{definition}

\begin{definition}[Abus et maladies]
\cle{Alcohol abuse} = l'abus d'alcool. \cle{Alcoholism} (prononcé « al-keu-ro-lizme ») = l'alcoolisme, la dépendance chronique à l'alcool. \cle{Liver disease} = la maladie du foie ; \cle{cirrhosis} (prononcé « si-rô-cisse ») = la cirrhose, destruction progressive du foie. \cle{Drink-driving} (Royaume-Uni) / \cle{drunk driving} (États-Unis) = la conduite en état d'ivresse.
\end{definition}

\begin{exemplebox}[Exemples traduits]
\eng{Drink-driving causes many accidents in Cameroon.} « La conduite en état d'ivresse cause de nombreux accidents au Cameroun. »

\eng{Binge drinking is dangerous for teenagers' brains.} « La beuverie est dangereuse pour le cerveau des adolescents. »

\eng{Years of alcohol abuse destroyed his liver; he now suffers from cirrhosis.} « Des années d'abus d'alcool ont détruit son foie ; il souffre maintenant de cirrhose. »

\eng{Alcoholism is a disease, not a habit.} « L'alcoolisme est une maladie, pas une habitude. »
\end{exemplebox}

\courssub{Prévention : moderation, legal drinking age, collocations}

\begin{definition}[Prévention et modération]
\cle{To drink in moderation} = boire avec modération. \cle{To limit one's intake} = limiter sa consommation. \cle{The legal drinking age} = l'âge légal pour boire de l'alcool. \cle{To ban the sale of alcohol to minors} = interdire la vente d'alcool aux mineurs. \cle{To stay sober} = rester sobre, ne pas boire.
\end{definition}

\begin{propriete}[Collocations essentielles]
\begin{itemize}
\item \cle{heavy drinking} = forte consommation d'alcool ; \eng{a heavy drinker} = un gros buveur.
\item \cle{to have a drinking problem} = avoir un problème avec l'alcool.
\item \cle{to stay sober} = rester sobre.
\item \cle{to be on the wagon} (idiome) = ne plus boire d'alcool du tout. Au contraire, \eng{to fall off the wagon} = rechuter, se remettre à boire.
\end{itemize}
\end{propriete}

\begin{exemplebox}[Exemples traduits]
\eng{He has been on the wagon for six months.} « Il ne boit plus d'alcool depuis six mois. »

\eng{Doctors advise people to limit their alcohol intake.} « Les médecins conseillent aux gens de limiter leur consommation d'alcool. »

\eng{She finally admitted that she had a drinking problem.} « Elle a finalement admis qu'elle avait un problème avec l'alcool. »

\eng{The government should ban the sale of alcohol to minors near schools.} « Le gouvernement devrait interdire la vente d'alcool aux mineurs près des écoles. »
\end{exemplebox}

\begin{savaistu}[Legal drinking age : 21 ou 18 ?]
Aux États-Unis, le \eng{legal drinking age} est fixé à \textbf{21 ans}, alors qu'il est de \textbf{18 ans} au Royaume-Uni. Cette différence est un sujet classique de débat au Bac : faut-il relever ou abaisser l'âge légal ? Les partisans du « 21 » invoquent la protection du cerveau des jeunes ; les partisans du « 18 » rappellent qu'à cet âge on peut voter et conduire. Au Cameroun, la loi interdit la vente d'alcool aux mineurs, mais son application reste difficile, notamment autour des écoles et des débits de boissons.
\end{savaistu}

\courssub{Tableau récapitulatif et comparaison}

\begin{center}
\begin{tabularx}{\linewidth}{|X|X|}
\hline
\rowcolor{popblueL} English & Français \\
\hline
alcohol & l'alcool \\
\hline
an alcoholic drink / beverage & une boisson alcoolisée \\
\hline
beer / wine / spirits & la bière / le vin / les spiritueux \\
\hline
a drunkard / an alcoholic & un ivrogne / un alcoolique \\
\hline
to get drunk & se saouler \\
\hline
drunk (après le verbe) / drunken (devant le nom) & ivre \\
\hline
a hangover & la gueule de bois \\
\hline
binge drinking & la beuverie, l'hyperalcoolisation ponctuelle \\
\hline
alcohol abuse / alcoholism & l'abus d'alcool / l'alcoolisme \\
\hline
liver disease / cirrhosis & la maladie du foie / la cirrhose \\
\hline
drink-driving (UK) / drunk driving (US) & la conduite en état d'ivresse \\
\hline
to drink in moderation & boire avec modération \\
\hline
to be on the wagon & ne plus boire (idiome) \\
\hline
the legal drinking age & l'âge légal de consommation d'alcool \\
\hline
\end{tabularx}
\end{center}

\begin{popfigure}
\begin{center}
\begin{tikzpicture}
\node[draw=popgreen, fill=popgreenL, rounded corners, align=center, text width=5.5cm] (mod) at (-4,0) {\textbf{Moderate drinking}\\ \eng{to drink in moderation}\\ \eng{to limit one's intake}\\ \eng{to stay in control}\\ Effets : détente, socialisation};
\node[draw=poporange, fill=poporangeL, rounded corners, align=center, text width=5.5cm] (abuse) at (4,0) {\textbf{Alcohol abuse}\\ \eng{binge drinking}\\ \eng{heavy drinking}\\ \eng{alcoholism}\\ Effets : cirrhose, accidents, dépendance};
\node[draw=popblue, fill=popblueL, rounded corners, align=center, text width=4cm] (root) at (0,2.6) {\textbf{Alcohol}\\ \eng{beer, wine, spirits}};
\draw[-{Stealth}, thick, popgreen] (root) -- (mod);
\draw[-{Stealth}, thick, poporange] (root) -- (abuse);
\end{tikzpicture}
\end{center}
\legende{Consommation modérée et abus d'alcool : deux chemins, deux conséquences.}
\end{popfigure}

\courssub{Exercice résolu et pièges à éviter}

\begin{exoresolu}[Complétez avec le mot correct]
Énoncé — Complétez chaque phrase avec le mot ou l'expression qui convient : \emph{drunk, drunken, spirits, on the wagon, binge drinking, hangover, sober, legal drinking age}.
\begin{enumerate}
\item A \ldots\ldots\ldots man was shouting in the street last night.
\item Whisky and vodka are strong \ldots\ldots\ldots .
\item He used to drink every day, but now he is \ldots\ldots\ldots .
\item \ldots\ldots\ldots at parties is very dangerous for students.
\item After drinking too much, I woke up with a bad \ldots\ldots\ldots .
\item In the United States, the \ldots\ldots\ldots is twenty-one.
\end{enumerate}
\tcblower
Solution détaillée :
\begin{enumerate}
\item \textbf{drunken} — le mot est placé devant le nom \emph{man}, donc on utilise l'épithète \emph{drunken} (et non \emph{drunk}, réservé à la position après le verbe).
\item \textbf{spirits} — le whisky et la vodka sont des alcools forts ; attention, \emph{spirits} au pluriel ne signifie pas « esprits ».
\item \textbf{on the wagon} — idiome : il ne boit plus du tout.
\item \textbf{Binge drinking} — boire beaucoup en peu de temps lors de fêtes ; majuscule en début de phrase.
\item \textbf{hangover} — la gueule de bois du lendemain.
\item \textbf{legal drinking age} — l'âge légal pour consommer de l'alcool, fixé à 21 ans aux États-Unis.
\end{enumerate}
\end{exoresolu}

\begin{attention}[Erreurs fréquentes des francophones]
\begin{itemize}
\item Ne dites jamais \eng{He is an alcohol} : dites \eng{He is an alcoholic} ou \eng{He has a drinking problem}.
\item Ne traduisez pas « spiritueux » par \eng{spiritual drinks} : le mot correct est \eng{spirits}.
\item N'écrivez pas \eng{a drunk man} dans un devoir soigné : préférez \eng{a drunken man} (épithète) ; réservez \eng{drunk} à la position attribut : \eng{The man is drunk}.
\item \eng{To be drunk} signifie « être ivre » ; ne confondez pas avec \eng{to drink} (boire, sans excès) : \eng{I drink water every day} ne veut pas dire que vous êtes ivre !
\item Prononciation : \emph{alcohol} se prononce « al-keu-rol » (l'accent sur la première syllabe) et \emph{alcoholic} « al-keu-ro-lik ».
\end{itemize}
\end{attention}

\begin{exercice}[À vous de jouer]
Traduisez en anglais : 1) La beuverie est un problème croissant chez les jeunes. 2) Il est resté sobre pendant toute la fête. 3) Un homme ivre a causé un accident près du marché. 4) Les médecins conseillent de boire avec modération. 5) Elle ne touche plus ni à la bière ni aux spiritueux : elle est « on the wagon ».
\end{exercice}

\begin{corrige}[À vous de jouer]
\begin{enumerate}
\item \eng{Binge drinking is a growing problem among young people.}
\item \eng{He stayed sober during the whole party.}
\item \eng{A drunken man caused an accident near the market.}
\item \eng{Doctors advise people to drink in moderation.}
\item \eng{She no longer touches beer or spirits: she is on the wagon.}
\end{enumerate}
\end{corrige}

\begin{aretenir}[L'essentiel de la section]
Le vocabulaire de l'abus d'alcool s'organise autour de quatre pôles : les boissons (\emph{beer, wine, spirits}), les états (\emph{drunk} après le verbe, \emph{drunken} devant le nom, \emph{hangover}), l'abus et ses conséquences (\emph{binge drinking, alcoholism, cirrhosis, drink-driving}) et la prévention (\emph{to drink in moderation, to stay sober, legal drinking age}). Retenez l'idiome \emph{to be on the wagon} et les deux grands faux amis : \emph{alcoholic} (jamais \emph{an alcohol} pour une personne) et \emph{spirits} (alcools forts, pas esprits).
\end{aretenir}

\coursec{Aging and elderly health}

After studying the dangers of \eng{alcohol abuse}, we now turn to a natural stage of life that concerns every family: \cle{growing old}. In Cameroon, our grandparents are often the pillars of the extended family; in Britain or the United States, many elderly people live in \eng{care homes}. Talking about old age in English requires a precise, respectful vocabulary — and a few traps to avoid. Let us first read a short text and observe the words in context.

\courssub{Observing the vocabulary in context}

\begin{textebox}[Growing old gracefully — a reading text]
My grandmother is eighty-five. She suffers from arthritis and hearing loss, but she stays active and eats a balanced diet. In our family, we take turns to care for her instead of sending her to a care home.

Every morning, she walks to the market in Buea with her old friend Mama Ngomba. “Walking keeps my bones strong,” she laughs. She retired at the age of sixty, after forty years as a primary school teacher, and she now receives a small pension. Her eyesight is poor, so she wears glasses, and her memory is not what it used to be, but the doctor says she shows no sign of dementia.

In many Western countries, elderly people often live in old people's homes, where professional caregivers look after them. In Cameroon, the extended family usually cares for its senior citizens at home. My grandmother says she wants to live to a ripe old age, surrounded by her grandchildren. Doctors agree that healthy ageing depends on three things: regular exercise, a balanced diet, and a medical check-up at least once a year. Growing old is not a disease; it is a stage of life that deserves respect, patience and love.
\end{textebox}

\textbf{Glossaire}\par
\begin{itemize}
\item \eng{arthritis} (n., prononcé « a-zraï-tiss ») : arthrite
\item \eng{hearing loss} (n.) : perte d'audition
\item \eng{to care for} (v.) : s'occuper de, prendre soin de
\item \eng{care home} (n.) : maison de retraite médicalisée
\item \eng{to retire} (v., prononcé « ri-taï-eur ») : prendre sa retraite
\item \eng{pension} (n., prononcé « pen-cheun ») : retraite (argent versé)
\item \eng{eyesight} (n.) : vue
\item \eng{dementia} (n., prononcé « di-men-chia ») : démence
\item \eng{caregiver} (n.) : aidant, soignant
\item \eng{to look after} (phrasal verb) : s'occuper de
\item \eng{senior citizens} (n.) : personnes âgées (terme poli)
\item \eng{a ripe old age} (expr.) : un âge avancé, une belle vieillesse
\item \eng{medical check-up} (n.) : bilan de santé, examen médical de routine
\end{itemize}

\textbf{Comprehension questions}\par
\begin{enumerate}
\item What two health problems does the grandmother suffer from?
\item At what age did she retire, and what was her job?
\item What are the three secrets of healthy ageing according to doctors?
\item How does the Cameroonian tradition differ from the Western one?
\end{enumerate}

\begin{corrige}[Comprehension questions]
1. She suffers from \eng{arthritis} and \eng{hearing loss}.\par
2. She retired \eng{at the age of sixty}; she was \eng{a primary school teacher}.\par
3. \eng{Regular exercise, a balanced diet, and a medical check-up at least once a year.}\par
4. \eng{In Cameroon, the extended family cares for elderly people at home, whereas in Western countries many of them live in old people's homes.}
\end{corrige}

\courssub{The stages of life: from youth to old age}

\begin{definition}[Ageing and its word family]
\cle{Ageing} (British spelling) or \cle{aging} (American spelling) is the process of growing old. The verb is \eng{to age} (prononcé « èidj ») : \eng{My grandfather is ageing well.} « Mon grand-père vieillit bien. » The adjective is \eng{aged} : \eng{a man aged seventy} « un homme âgé de soixante-dix ans ». Other useful expressions: \eng{to grow old} (vieillir), \eng{to get older} (prendre de l'âge), \eng{old age} (la vieillesse, nom).
\end{definition}

\begin{popfigure}
\begin{tikzpicture}[
stage/.style={rectangle, rounded corners, draw=popblue, fill=popblueL, minimum width=2.6cm, minimum height=1.1cm, align=center, font=\small\bfseries},
voc/.style={rectangle, draw=poporange, fill=poporangeL, rounded corners, align=center, font=\scriptsize},
arr/.style={-{Stealth[length=3mm]}, thick, popdark}]
\node[stage, align=center] (y) at (0,0){Youth\\ \scriptsize jeunesse};
\node[stage, align=center] (a) at (3.4,0){Adulthood\\ \scriptsize âge adulte};
\node[stage, align=center] (m) at (6.8,0){Middle age\\ \scriptsize âge mûr};
\node[stage, align=center] (r) at (10.2,0){Retirement\\ \scriptsize retraite};
\node[stage, align=center] (o) at (13.6,0){Old age\\ \scriptsize vieillesse};
\draw[arr] (y) -- (a);
\draw[arr] (a) -- (m);
\draw[arr] (m) -- (r);
\draw[arr] (r) -- (o);
\node[voc, align=center] at (0,-1.0){a teenager\\ young people};
\node[voc, align=center] at (3.4,-1.0){an adult\\ to start a family};
\node[voc, align=center] at (6.8,-1.0){middle-aged\\ in one's fifties};
\node[voc, align=center] at (10.2,-1.0){to retire at 60\\ to receive a pension};
\node[voc, align=center] at (13.6,-1.0){elderly people\\ senior citizens};
\end{tikzpicture}
\legende{The stages of life — les étapes de la vie et le vocabulaire associé.}
\end{popfigure}

\begin{propriete}[Politeness: how to talk about old people]
In English, \eng{old people} can sound rude. Prefer the polite terms \cle{elderly people} or \cle{senior citizens}. Note that \eng{elderly} is an adjective only: you can say \eng{an elderly lady} but never \eng{an elderly} alone as a noun. However, \eng{the elderly} (les personnes âgées, en général, avec un verbe au pluriel) is correct and polite : \eng{The elderly need our respect.} « Les personnes âgées ont besoin de notre respect. »
\end{propriete}

\begin{exemplebox}[Examples]
\eng{She looks after her elderly parents.} « Elle s'occupe de ses parents âgés. »\par
\eng{He retired at the age of sixty.} « Il a pris sa retraite à soixante ans. »\par
\eng{Senior citizens pay half price on some British trains.} « Les personnes âgées paient moitié prix dans certains trains britanniques. »\par
\eng{My grandmother lived to a ripe old age.} « Ma grand-mère a vécu jusqu'à un âge avancé. »
\end{exemplebox}

\courssub{Health problems related to age}

\begin{definition}[Dementia]
\cle{Dementia} is a serious illness of the mind that affects memory and thinking, mostly in old age (démence). \eng{Alzheimer's disease} (prononcé « al-zaï-meuz ») is the most common form of dementia. Related expressions: \eng{to lose one's memory} (perdre la mémoire), \eng{memory loss} (perte de mémoire).
\end{definition}

\begin{center}
\begin{tabularx}{\linewidth}{|X|X|X|}
\hline
\rowcolor{popblueL} English & Français & Example \\
\hline
arthritis & arthrite & She suffers from arthritis in her knees. \\
\hline
Alzheimer's disease & maladie d'Alzheimer & Alzheimer's disease affects the memory. \\
\hline
dementia & démence & Dementia is common in very old age. \\
\hline
hearing loss & perte d'audition & He has hearing loss and wears a hearing aid. \\
\hline
poor eyesight & mauvaise vue & Poor eyesight makes reading difficult. \\
\hline
fragile bones & os fragiles & Elderly people often have fragile bones. \\
\hline
\end{tabularx}
\end{center}

\begin{propriete}[Key collocations]
Learn these fixed combinations as whole units:
\begin{itemize}
\item \eng{to suffer from arthritis} — souffrir d'arthrite (jamais \eng{to suffer arthritis} sans \eng{from})
\item \eng{to lose one's memory} — perdre la mémoire
\item \eng{to retire at 60 / at the age of 60} — prendre sa retraite à 60 ans
\item \eng{to live to a ripe old age} — vivre jusqu'à un âge avancé
\item \eng{to wear glasses / a hearing aid} — porter des lunettes / un appareil auditif
\item \eng{to take turns to look after someone} — se relayer pour s'occuper de quelqu'un
\end{itemize}
\end{propriete}

\courssub{Caring for the elderly}

\begin{definition}[Care vocabulary]
A \cle{care home} (also \eng{old people's home}, \eng{nursing home}) is an institution where elderly people live and receive medical care. A \cle{caregiver} (British also \eng{carer}) is a person who looks after a sick or elderly person. A \cle{pension} is the money received after retirement; a \eng{pensioner} is a retired person. Do not confuse \eng{to retire} (prendre sa retraite) with \eng{retirement} (la retraite, l'état) : \eng{He is enjoying his retirement.} « Il profite de sa retraite. »
\end{definition}

\begin{exemplebox}[To look after = to care for]
\eng{We take turns to look after our grandfather.} « Nous nous relayons pour nous occuper de notre grand-père. »\par
\eng{She works as a caregiver in a care home in Douala.} « Elle travaille comme soignante dans une maison de retraite à Douala. »\par
\eng{He receives a small pension every month.} « Il reçoit une petite retraite chaque mois. »
\end{exemplebox}

\begin{attention}[False friends: sensible, eventually, assist]
\eng{Sensible} is a false friend: it means « raisonnable » in English; French « sensible » is \eng{sensitive}. \eng{My grandmother is a sensible woman.} « Ma grand-mère est une femme raisonnable. » — \eng{She is very sensitive to criticism.} « Elle est très sensible à la critique. »\par
Similarly, \eng{eventually} means « finalement », not « éventuellement » (which is \eng{possibly / maybe}). \eng{He eventually moved to a care home.} « Il a finalement emménagé en maison de retraite. »\par
Also note: \eng{to assist} someone = aider quelqu'un (formel) ; « assister à » une réunion = \eng{to attend a meeting}. Finally, \eng{a pension} en anglais ne désigne jamais une « pension de famille » (which is \eng{a guesthouse} ou \eng{a bed and breakfast}).
\end{attention}

\courssub{Healthy ageing}

\begin{aretenir}[The vocabulary of healthy ageing]
\eng{healthy ageing} (bien vieillir), \eng{to stay active} (rester actif), \eng{a balanced diet} (une alimentation équilibrée), \eng{regular exercise} (de l'exercice régulier), \eng{a medical check-up} (un bilan de santé).\par
\eng{Regular exercise and a balanced diet are the secrets of healthy ageing.} « L'exercice régulier et une alimentation équilibrée sont les secrets pour bien vieillir. »
\end{aretenir}

\begin{savaistu}[Elderly people in Cameroon and in the West]
In Cameroon, elderly people usually live with their extended family, and children and grandchildren take turns to care for them. In the United Kingdom and the United States, \eng{care homes} and \eng{nursing homes} are much more common. This is an excellent cultural comparison to use in your speaking exam: \eng{In my country, we look after our elderly relatives at home, whereas in Britain many senior citizens live in care homes.}
\end{savaistu}

\begin{exoresolu}[Gap-filling exercise]
Complete the sentences with: \eng{dementia — pension — arthritis — look after — senior citizens — check-up}.\par
1. My uncle suffers from \dots ; his hands hurt when it rains.\par
2. She retired last year and now receives a monthly \dots .\par
3. In our village, the children \dots\ their elderly parents.\par
4. The doctor found no sign of \dots ; his memory is excellent.\par
5. \dots\ should have a medical \dots\ every year.\par
\tcblower
1. \eng{arthritis} (douleurs articulaires).\par
2. \eng{pension} (argent de la retraite).\par
3. \eng{look after} (s'occuper de).\par
4. \eng{dementia} (la mémoire est bonne, donc pas de démence).\par
5. \eng{Senior citizens} ; \eng{check-up} (bilan de santé annuel ; majuscule en début de phrase).
\end{exoresolu}

\begin{exercice}[Your turn]
Translate into English: 1. Elle s'occupe de ses parents âgés. 2. Il a pris sa retraite à soixante ans. 3. Ma grand-mère souffre d'arthrite mais elle reste active. 4. Les personnes âgées ont souvent les os fragiles. 5. Il a finalement emménagé dans une maison de retraite.
\end{exercice}

\begin{corrige}[Exercice « Your turn »]
1. \eng{She looks after her elderly parents.}\par
2. \eng{He retired at the age of sixty.}\par
3. \eng{My grandmother suffers from arthritis but she stays active.}\par
4. \eng{Elderly people often have fragile bones.}\par
5. \eng{He eventually moved to a care home.} (attention : \eng{eventually} = finalement).
\end{corrige}

\coursec{Word families, pronunciation and spelling}

Dans la section précédente, nous avons étudié le vocabulaire du vieillissement et de la santé des personnes âgées : \eng{aging}, \eng{elderly people}, \eng{arthritis}, \eng{caregiver}. Vous avez peut-être remarqué que certains de ces mots se ressemblent : \eng{care} et \eng{caregiver}, \eng{health} et \eng{healthy}. Ce n'est pas un hasard. L'anglais, comme le français, construit des \cle{familles de mots} autour d'une racine commune. Apprendre ces familles — et savoir prononcer et orthographier correctement chaque membre — est la clé pour enrichir rapidement votre vocabulaire et éviter les erreurs au Bac.

\courssub{Observing word families: the case of \eng{addict}}

\begin{textebox}[A radio interview — Cameroon Health Watch]
\textbf{Journalist:} Doctor, drug \eng{addiction} is a growing problem among young people in our cities. How does someone become an \eng{addict}?

\textbf{Dr. Mbarga:} It usually starts with curiosity. A teenager tries a drug at a party. But many substances, like nicotine, are highly \eng{addictive} — the body quickly needs more. After a few months, the person is \eng{addicted} and cannot stop without help.

\textbf{Journalist:} So \eng{prevention} is essential?

\textbf{Dr. Mbarga:} Absolutely. As we say in English, “\eng{Prevention is better than cure.}” We must inform young people before they touch drugs, alcohol or tobacco. Once \eng{addiction} begins, treatment is long and difficult, and \eng{withdrawal} symptoms can be painful.

\textbf{Journalist:} What about \eng{alcoholism}? Is it also an addiction?

\textbf{Dr. Mbarga:} Yes. An \eng{alcoholic} is a person who is addicted to alcohol. \eng{Alcoholism} destroys families, careers and health — especially the liver.
\end{textebox}

\textbf{Glossaire}

\begin{center}
\begin{tabularx}{\linewidth}{|l|l|X|}
\hline
\rowcolor{popblueL} Mot & Nature & Traduction \\
\hline
addict & nom & un toxicomane \\
\hline
addicted & adjectif & dépendant, accro \\
\hline
addictive & adjectif & qui crée une dépendance \\
\hline
addiction & nom & la dépendance, la toxicomanie \\
\hline
prevention & nom & la prévention \\
\hline
withdrawal & nom & le sevrage \\
\hline
alcoholic & adjectif / nom & alcoolique ; un alcoolique \\
\hline
alcoholism & nom & l'alcoolisme \\
\hline
\end{tabularx}
\end{center}

\textbf{Observez bien} : dans ce dialogue, un seul mot de base, \eng{addict}, donne naissance à quatre mots différents. C'est exactement comme en français : \emph{toxicomane}, \emph{toxicomanie}, \emph{accro}. Voyons maintenant la mécanique de ces familles.

\courssub{Word families: form and use}

\begin{definition}[Famille de mots]
Une \cle{famille de mots} (\eng{word family}) est un ensemble de mots construits sur la même racine, dont la nature grammaticale change grâce à des \cle{suffixes} : nom, verbe, adjectif, adverbe. Exemple : \eng{addict} (nom) → \eng{addicted} (adjectif) → \eng{addiction} (nom) → \eng{addictive} (adjectif).
\end{definition}

\begin{propriete}[Les suffixes -tion / -sion et -ism]
\begin{itemize}
\item Pour former un \textbf{nom d'action ou d'état} à partir d'un verbe, on ajoute souvent \cle{-tion} ou \cle{-sion} : \eng{prevent} → \eng{prevention} (la prévention), \eng{addict} → \eng{addiction} (la dépendance), \eng{inject} → \eng{injection}, \eng{decide} → \eng{decision}.
\item Pour exprimer un \textbf{état, une maladie ou une doctrine}, on ajoute \cle{-ism} : \eng{alcohol} → \eng{alcoholism} (l'alcoolisme), \eng{smoke} → \eng{smokeless} (sans fumée) ne prend pas -ism, mais \eng{alcoholism}, \eng{racism} (le racisme) en sont des exemples classiques.
\end{itemize}
\end{propriete}

\begin{propriete}[Les suffixes -ive et -less]
\begin{itemize}
\item \cle{-ive} forme des adjectifs exprimant une qualité active : \eng{addict} → \eng{addictive} (qui crée une dépendance), \eng{prevent} → \eng{preventive} (préventif).
\item \cle{-less} forme des adjectifs exprimant l'\textbf{absence} de quelque chose : \eng{smoke} → \eng{smokeless} (sans fumée), \eng{care} → \eng{careless} (négligent, sans soin), \eng{harm} → \eng{harmless} (inoffensif). Attention au sens : \eng{careless} signifie « négligent », pas « sans soignant » !
\end{itemize}
\end{propriete}

\begin{popfigure}
\begin{center}
\begin{tikzpicture}[
grow=down,
level distance=2.2cm,
level 1/.style={sibling distance=4.6cm},
edge from parent/.style={draw=popblue, thick, -{Stealth}},
every node/.style={align=center}
]
\node[draw=popdark, fill=popgoldL, rounded corners, font=\bfseries, align=center]{addict\\ \footnotesize (nom : un toxicomane)}
  child { node[draw=popblue, fill=popblueL, rounded corners, align=center]{addicted\\ \footnotesize (adj.) \\ \footnotesize\itshape “He is addicted to drugs.”} }
  child { node[draw=poppurple, fill=poppurpleL, rounded corners, align=center]{addiction\\ \footnotesize (nom) \\ \footnotesize\itshape “She is fighting her addiction.”} }
  child { node[draw=popgreen, fill=popgreenL, rounded corners, align=center]{addictive\\ \footnotesize (adj.) \\ \footnotesize\itshape “Nicotine is addictive.”} };
\end{tikzpicture}
\legende{Arbre de la famille du mot \eng{addict} : une racine, trois dérivés, une phrase d'exemple par branche.}
\end{center}
\end{popfigure}

\courssub{Other useful word families}

\begin{center}
\begin{tabularx}{\linewidth}{|l|l|X|X|}
\hline
\rowcolor{popblueL} Racine & Dérivés & Exemple anglais & Traduction \\
\hline
smoke (v/n) & smoker, smoking, smokeless & Smoking is forbidden here. & Il est interdit de fumer ici. \\
\hline
alcohol (n) & alcoholic (adj/n), alcoholism & Alcoholism is a disease. & L'alcoolisme est une maladie. \\
\hline
health (n) & healthy, unhealthy & Junk food is unhealthy. & La malbouffe est mauvaise pour la santé. \\
\hline
care (v/n) & careful, careless, caregiver & Be careful with medicines. & Fais attention avec les médicaments. \\
\hline
prevent (v) & prevention, preventive & Prevention is better than cure. & Mieux vaut prévenir que guérir. \\
\hline
depend (v) & dependent, dependency, dependence & He is dependent on alcohol. & Il est dépendant de l'alcool. \\
\hline
\end{tabularx}
\end{center}

\begin{exemplebox}[Exemples traduits]
\begin{itemize}
\item \eng{Nicotine is highly addictive.} « La nicotine crée une forte dépendance. »
\item \eng{Prevention is better than cure.} « Mieux vaut prévenir que guérir. »
\item \eng{Many smokers want to quit smoking.} « Beaucoup de fumeurs veulent arrêter de fumer. »
\item \eng{His father is an alcoholic; he has suffered from alcoholism for years.} « Son père est alcoolique ; il souffre d'alcoolisme depuis des années. »
\item \eng{A caregiver looks after elderly people.} « Un aidant s'occupe des personnes âgées. »
\item \eng{She is dependent on painkillers.} « Elle est dépendante aux antidouleurs. »
\end{itemize}
\end{exemplebox}

\begin{attention}[La préposition avec \eng{addicted} et apparentés]
On dit toujours \cle{addicted TO}, jamais \eng{addicted of} ni \eng{addicted by}. \eng{He is addicted to drugs.} « Il est dépendant de la drogue. » De même : \eng{dependent ON}, \eng{to suffer FROM arthritis}, \eng{to be interested IN health}. Les prépositions anglaises ne se calquent pas sur le français !
\end{attention}

\courssub{Pronunciation and word stress}

En anglais, chaque mot a une \cle{syllabe accentuée} (\eng{stressed syllable}) prononcée plus fortement. Piège majeur : \textbf{l'accent change souvent de place à l'intérieur d'une même famille}, contrairement au français où l'accent est stable.

\begin{center}
\begin{tabularx}{\linewidth}{|l|l|X|}
\hline
\rowcolor{popblueL} Mot & Prononciation (lettres françaises) & Accent tonique \\
\hline
addict (nom) & « a-dikte » & \textbf{AD}dict (1\iere{} syllabe) \\
\hline
addicted & « a-dik-tid » & a\textbf{DDIC}ted (2\ieme{} syllabe) \\
\hline
addiction & « a-dik-cheun » & a\textbf{DDIC}tion (2\ieme{} syllabe) \\
\hline
alcohol & « al-keu-hol » & \textbf{AL}cohol (1\iere{} syllabe) \\
\hline
alcoholic & « al-keu-ho-lik » & alco\textbf{HOL}ic (3\ieme{} syllabe) \\
\hline
alcoholism & « al-keu-ho-li-zeum » & \textbf{AL}coholism (1\iere{} syllabe) \\
\hline
cigarette & « si-gueu-ret » & ciga\textbf{RETTE} (dernière syllabe) \\
\hline
disease & « di-ziz » & di\textbf{SEASE} (2\ieme{} syllabe) \\
\hline
elderly & « el-deu-li » & \textbf{EL}derly (1\iere{} syllabe) \\
\hline
withdrawal & « wi-th-dro-ol » & with\textbf{DRAW}al (2\ieme{} syllabe) \\
\hline
\end{tabularx}
\end{center}

\begin{attention}[Prononciation : pièges pour les francophones]
\begin{itemize}
\item \eng{disease} se prononce « di-ziz » avec un son \textbf{z}, pas « di-siss ». De même \eng{decease} mais attention à ne pas confondre avec \eng{deceased} (décédé).
\item \eng{cigarette} : l'accent est sur la dernière syllabe, comme en français, mais le « g » se prononce « gueu » (dur), jamais « j ».
\item \eng{healthy} se prononce « hel-thi » : le \textbf{th} se dit avec la langue entre les dents, jamais « z » ni « s ».
\item \eng{withdrawal} : le \textbf{th} est également interdentaire ; ne pas dire « wi-tra-oual ».
\end{itemize}
\end{attention}

\courssub{Tricky spelling}

\begin{propriete}[Orthographes à mémoriser]
\begin{itemize}
\item \cle{tobacco} : un seul \textbf{b} et deux \textbf{c} (jamais \eng{tabacco}).
\item \cle{alcoholic} : un seul \textbf{l} après « alco- » (\eng{alco-} + \eng{-holic}), mais \eng{alcohol} en a deux.
\item \cle{ageing} (britannique) / \cle{aging} (américain) : les deux graphies sont acceptées ; l'orthographe britannique \eng{ageing} est recommandée au Cameroun.
\item \cle{withdrawal} : un seul \textbf{l} à la fin (jamais \eng{withdrawall}).
\item \cle{healthy} : avec \textbf{-th-} ; le nom est \eng{health}, l'adverbe \eng{healthily} (changement de \eng{y} en \eng{i}).
\item \cle{prevention} : pas de \eng{t} après le \eng{n} (pas de \eng{preventition}).
\end{itemize}
\end{propriete}

\begin{methode}[Comment mémoriser le vocabulaire durablement]
\begin{enumerate}
\item Apprenez les mots \textbf{par famille} : retenez ensemble \eng{addict, addicted, addiction, addictive} plutôt qu'un mot isolé.
\item Apprenez chaque mot \textbf{avec sa collocation} : \eng{to suffer from arthritis}, \eng{to be addicted to drugs}, \eng{to quit smoking}, \eng{to look after elderly people}, \eng{dependent on}.
\item Réemployez chaque mot dans \textbf{une phrase personnelle} sur votre vie ou votre entourage : \eng{My grandfather suffers from arthritis.}
\item Vérifiez \textbf{l'accent tonique} en prononçant le mot à voix haute trois fois.
\end{enumerate}
\end{methode}

\begin{exoresolu}[Compléter avec le bon membre de la famille]
Énoncé — Complétez chaque phrase avec la forme correcte du mot entre parenthèses :
\begin{enumerate}
\item Cigarettes are very \ldots\ (addict).
\item His \ldots\ to alcohol destroyed his career. (addicted)
\item \ldots\ is better than cure. (prevent)
\item She is a \ldots\ ; she smokes twenty cigarettes a day. (smoke)
\end{enumerate}
\tcblower
Solution détaillée :
\begin{enumerate}
\item \eng{addictive} : il faut un adjectif après le verbe \eng{are} → « Les cigarettes créent une forte dépendance. »
\item \eng{addiction} : après l'adjectif possessif \eng{his}, il faut un nom → \eng{His addiction to alcohol\ldots} « Son addiction à l'alcool\ldots »
\item \eng{Prevention} : sujet de la phrase, donc nom en \eng{-tion} → « Mieux vaut prévenir que guérir. »
\item \eng{smoker} : après l'article \eng{a}, il faut un nom de personne → « C'est une fumeuse. »
\end{enumerate}
\end{exoresolu}

\begin{savaistu}[Un proverbe en or pour le Bac]
Le proverbe \eng{Prevention is better than cure} (« Mieux vaut prévenir que guérir ») est très apprécié des correcteurs dans les sujets d'expression écrite du Bac sur la santé. Placé en introduction ou en conclusion d'une composition sur la drogue, le tabac ou l'alcool, il montre que vous maîtrisez la culture anglophone. Retenez aussi \eng{Health is wealth} (« La santé est une richesse »).
\end{savaistu}

\begin{exercice}[À vous de jouer]
\begin{enumerate}
\item Donnez la famille complète de \eng{care} (nom, adjectifs positif et négatif, nom de personne).
\item Traduisez : « Mon oncle est dépendant du tabac ; il essaie d'arrêter de fumer. »
\item Soulignez la syllabe accentuée dans : \eng{alcoholic}, \eng{addiction}, \eng{cigarette}, \eng{elderly}, \eng{withdrawal}.
\end{enumerate}
\end{exercice}

\begin{corrige}[Exercice « À vous de jouer »]
\begin{enumerate}
\item \eng{care} (nom/verbe) → \eng{careful} (soigneux) → \eng{careless} (négligent) → \eng{caregiver} (aidant).
\item \eng{My uncle is addicted to tobacco; he is trying to quit smoking.} (Attention : \eng{addicted TO}, et \eng{quit + -ing}.)
\item alco\textbf{HOL}ic, a\textbf{DDIC}tion, ciga\textbf{RETTE}, \textbf{EL}derly, with\textbf{DRAW}al.
\end{enumerate}
\end{corrige}

\begin{aretenir}[L'essentiel de la section]
\begin{itemize}
\item Les familles de mots se construisent avec des suffixes : \eng{-tion} (action/état), \eng{-ism} (état, maladie), \eng{-ive} (qualité), \eng{-less} (absence).
\item \eng{addicted TO}, \eng{dependent ON}, jamais \eng{of} ; \eng{suffer FROM}, \eng{quit + -ing}.
\item L'accent tonique change dans la famille : \eng{ADdict} mais \eng{aDDICtion} ; \eng{ALcohol} mais \eng{alcoHOLic}.
\item Orthographes piégeuses : \eng{tobacco}, \eng{alcoholic}, \eng{ageing/aging}, \eng{withdrawal}, \eng{healthy/healthily}.
\end{itemize}
\end{aretenir}

\coursec{Using the vocabulary in context}

Dans la section précédente, nous avons appris à reconnaître les familles de mots (\eng{addict / addiction / addictive}), à prononcer correctement les mots difficiles et à éviter les pièges orthographiques. Mais connaître des mots ne suffit pas : un bon élève de Terminale doit savoir les \cle{réemployer} dans un texte, un dialogue, une phrase transformée ou une production écrite. C'est exactement l'objectif de cette dernière section : passer du \emph{savoir} au \emph{savoir-faire}.

\courssub{Guided reading: guessing meaning from context}

\begin{methode}[Comment deviner le sens d'un mot inconnu dans un texte]
Face à un mot inconnu, ne paniquez pas et ne sortez pas le dictionnaire tout de suite. Suivez ces quatre étapes :
\begin{enumerate}
\item \textbf{Lisez le titre} du texte : il annonce le thème et active votre vocabulaire connu.
\item \textbf{Repérez les mots connus} autour du mot inconnu : ils donnent des indices.
\item \textbf{Devinez le sens grâce au contexte} : observez le \cle{type du mot} (nom, verbe, adjectif), les mots qui l'entourent, et les \cle{préfixes et suffixes}. Exemple : \eng{unhealthy} = \eng{un-} (négation) + \eng{healthy} (en bonne santé) = « pas en bonne santé, malsain ».
\item \textbf{Vérifiez au glossaire} ou au dictionnaire pour confirmer votre hypothèse.
\end{enumerate}
\end{methode}

\begin{popfigure}
\begin{tikzpicture}[
  node distance=0.55cm,
  box/.style={draw=popblue, fill=popblueL, rounded corners, thick, align=center, text width=6.2cm, minimum height=1.05cm, font=\small},
  arr/.style={-{Stealth[length=3mm]}, thick, popdark}
]
\node[box, align=center] (s1){\textbf{1. Read the title}\\ Lire le titre};
\node[box, below=of s1, align=center] (s2){\textbf{2. Spot known words}\\ Repérer les mots connus};
\node[box, below=of s2, align=center] (s3){\textbf{3. Guess unknown words from context}\\ Deviner les mots inconnus par le contexte};
\node[box, below=of s3, fill=popgreenL, draw=popgreen, align=center] (s4){\textbf{4. Check the glossary}\\ Vérifier au glossaire};
\draw[arr] (s1) -- (s2);
\draw[arr] (s2) -- (s3);
\draw[arr] (s3) -- (s4);
\end{tikzpicture}
\legende{Organigramme de la méthode de lecture : du titre au glossaire.}
\end{popfigure}

\begin{textebox}[Awareness campaign in Douala — reading text]
Health experts in Douala have launched an awareness campaign against drug and alcohol abuse among students. The campaign, which started last month, targets secondary schools in the city, where cases of binge drinking and smoking have been reported.

According to a nurse who works in a rehabilitation centre, many young addicts suffer from withdrawal symptoms when they try to quit. “When a teenager stops taking drugs suddenly, his body reacts badly: he may feel sick, anxious or depressed,” she explains. “That is why quitting alone is so difficult. Victims need medical care and family support.”

The nurse also warns against passive smoking: students who do not smoke but spend time with smokers can also damage their lungs. She advises families to care for victims instead of rejecting them, and schools to teach healthy habits early. “Prevention is better than cure,” she adds. “If we keep young people informed about the dangers of tobacco, alcohol and drugs, fewer of them will become addicted.”

The campaign also includes workshops on aging well: students are encouraged to help their elderly grandparents eat healthy food, exercise regularly and go for medical check-ups. A healthy lifestyle, the experts say, protects people at every age.
\end{textebox}

\begin{definition}[Glossaire du texte]
\begin{itemize}
\item \eng{awareness campaign} (nom) : campagne de sensibilisation.
\item \eng{binge drinking} (nom) : consommation excessive d'alcool en peu de temps (« biture express »).
\item \eng{withdrawal symptoms} (nom pluriel) : symptômes de sevrage.
\item \eng{to quit} (verbe) : arrêter, renoncer (à une habitude).
\item \eng{passive smoking} (nom) : tabagisme passif.
\item \eng{care for} (verbe) : prendre soin de.
\item \eng{check-up} (nom) : bilan de santé, visite de contrôle.
\item \eng{elderly} (adjectif) : âgé ; \eng{the elderly} = les personnes âgées.
\end{itemize}
\end{definition}

\begin{exemplebox}[Deviner un mot grâce au contexte : exemple traité]
Prenons la phrase : \eng{Many young addicts suffer from withdrawal symptoms when they try to quit.}
\begin{itemize}
\item \textbf{Type du mot} : \eng{withdrawal symptoms} est un groupe nominal (adjectif + nom pluriel).
\item \textbf{Mots qui l'entourent} : \eng{suffer from} + maladie indique toujours un problème de santé (\eng{suffer from malaria, suffer from headaches}).
\item \textbf{Indice temporel} : \eng{when they try to quit} = « quand ils essaient d'arrêter ».
\item \textbf{Conclusion} : \eng{withdrawal symptoms} désigne donc les troubles de santé qui apparaissent quand un toxicomane arrête la drogue → « symptômes de sevrage ». Le contexte seul a suffi !
\end{itemize}
\end{exemplebox}

\textbf{Questions de compréhension.}\par
\begin{enumerate}
\item Who launched the awareness campaign, and where?
\item What happens to young addicts when they try to quit?
\item Why is passive smoking dangerous?
\item What two pieces of advice does the nurse give to families and schools?
\item What does the campaign say about the elderly?
\end{enumerate}

\begin{corrige}[Questions de compréhension]
\begin{enumerate}
\item Health experts in Douala launched the campaign.
\item They suffer from withdrawal symptoms: they may feel sick, anxious or depressed.
\item Because people who spend time with smokers can also damage their lungs.
\item Families should care for victims instead of rejecting them, and schools should teach healthy habits early.
\item Students should help their elderly grandparents eat healthy food, exercise and go for medical check-ups.
\end{enumerate}
\end{corrige}

\courssub{Model dialogue: interviewing a nurse}

\begin{textebox}[An interview at the school infirmary]
\textbf{Student:} Good morning, Madam. May I ask you a few questions about tobacco and alcohol?\\
\textbf{Nurse:} Of course, go ahead.\\
\textbf{Student:} Why is smoking so dangerous for young people?\\
\textbf{Nurse:} Tobacco damages the lungs and causes cancer. Worse, nicotine is very addictive: a teenager who starts smoking can quickly become a heavy smoker.\\
\textbf{Student:} And what about alcohol?\\
\textbf{Nurse:} Many students practise binge drinking at parties. It harms the liver and the brain, and drunk young people often have accidents.\\
\textbf{Student:} Can an addict simply stop?\\
\textbf{Nurse:} It is hard, because of withdrawal symptoms. But with medical care and family support, victims can recover.\\
\textbf{Student:} What can schools do?\\
\textbf{Nurse:} They should raise awareness and keep students informed about these dangers. Prevention is better than cure!\\
\textbf{Student:} Thank you very much, Madam.
\end{textebox}

\begin{attention}[Faux amis et calques du français]
\begin{itemize}
\item Ne dites pas \eng{to consume drugs} (calque de « consommer de la drogue ») : c'est rare et peu naturel. On dit \eng{to take drugs} ou \eng{to use drugs}.
\item \eng{To inform} exige un complément : on ne dit pas \eng{to inform about health problems} tout seul, mais \eng{to keep people informed about health problems} (« tenir les gens informés des problèmes de santé »).
\item \eng{To assist} ne veut pas dire « assister à » (= \eng{to attend}) mais « aider » : \eng{The nurse assists the victims.} = L'infirmière aide les victimes.
\item « Sensibiliser » se traduit par \eng{to raise awareness}, jamais par \eng{to sensitize}.
\end{itemize}
\end{attention}

\begin{exemplebox}[Exemples traduits]
\eng{The campaign keeps people informed about the dangers of smoking.} « La campagne tient les gens informés des dangers du tabac. »\par
\eng{Schools should raise awareness about alcohol abuse.} « Les écoles devraient sensibiliser à l'abus d'alcool. »\par
\eng{He decided to quit smoking last year.} « Il a décidé d'arrêter de fumer l'année dernière. »\par
\eng{The elderly need regular medical check-ups.} « Les personnes âgées ont besoin de bilans de santé réguliers. »
\end{exemplebox}

\courssub{Practice: gap-fill and rephrasing}

\begin{center}
\begin{tabularx}{\linewidth}{|X|X|X|}
\hline
\rowcolor{popblueL} English & Français & Collocation fréquente \\
\hline
addicted (to) & accro (à) & addicted to drugs / alcohol \\
\hline
withdrawal symptoms & symptômes de sevrage & suffer from withdrawal symptoms \\
\hline
binge drinking & alcoolisation excessive & practise binge drinking \\
\hline
elderly & âgé & the elderly, elderly people \\
\hline
to quit & arrêter & quit smoking / drinking \\
\hline
heavy smoker & gros fumeur & become a heavy smoker \\
\hline
awareness & sensibilisation & raise awareness about \\
\hline
to recover & se rétablir & recover from an addiction \\
\hline
\end{tabularx}
\end{center}

\begin{exoresolu}[Phrases à compléter]
Énoncé — Complétez avec : \eng{addicted, withdrawal, binge drinking, elderly, to quit}.
\begin{enumerate}
\item Many teenagers are \ldots{} to cigarettes because of nicotine.
\item When he stopped taking drugs, he suffered from \ldots{} symptoms.
\item \ldots{} at parties is very dangerous for the liver.
\item We should respect and help \ldots{} people in our community.
\item It is difficult \ldots{} smoking without support.
\end{enumerate}
\tcblower
Solution détaillée :
\begin{enumerate}
\item \eng{addicted} — structure : \eng{to be addicted to} + nom.
\item \eng{withdrawal} — \eng{withdrawal symptoms} = symptômes de sevrage.
\item \eng{Binge drinking} — sujet de la phrase (nom en \eng{-ing}), majuscule en début de phrase.
\item \eng{elderly} — \eng{elderly people} = les personnes âgées.
\item \eng{to quit} — structure : \eng{It is difficult to} + base verbale.
\end{enumerate}
\end{exoresolu}

\begin{methode}[Reformuler avec un mot de la même famille]
Pour transformer une phrase sans changer son sens :
\begin{enumerate}
\item Repérez le mot clé (souvent un verbe : \eng{smokes}).
\item Trouvez un autre mot de sa famille (\eng{smoker}, \eng{smoking}).
\item Reconstruisez la phrase autour de ce mot : \eng{He smokes a lot} → \eng{He is a heavy smoker.}
\end{enumerate}
Autres exemples : \eng{She is addicted to alcohol.} → \eng{She suffers from alcohol addiction.} ; \eng{He drinks too much.} → \eng{He is a heavy drinker.}
\end{methode}

\begin{exercice}[Reformulation]
Transformez chaque phrase avec le mot entre parenthèses :
\begin{enumerate}
\item He smokes a lot. (\eng{smoker})
\item She is addicted to drugs. (\eng{addiction})
\item They drink alcohol excessively at weekends. (\eng{binge})
\item The nurse takes care of old people. (\eng{elderly})
\end{enumerate}
\end{exercice}

\begin{corrige}[Exercice 1]
\begin{enumerate}
\item He is a heavy smoker.
\item She suffers from drug addiction.
\item They practise binge drinking at weekends.
\item The nurse takes care of the elderly / of elderly people.
\end{enumerate}
\end{corrige}

\courssub{Guided writing: “How to keep young people away from drugs”}

\begin{methode}[Rédiger un paragraphe de 8 à 10 lignes]
\begin{enumerate}
\item \textbf{Introduction} (1-2 lignes) : présentez le problème (\eng{Drug abuse is a serious problem among young people in Cameroon.}).
\item \textbf{Causes ou dangers} (2-3 lignes) : \eng{addiction, withdrawal symptoms, health problems}.
\item \textbf{Solutions} (3-4 lignes) : \eng{raise awareness, keep students informed, teach healthy habits, care for victims}.
\item \textbf{Conclusion} (1 ligne) : une phrase forte, par exemple \eng{Prevention is better than cure.}
\end{enumerate}
\end{methode}

\begin{exemplebox}[Modèle de paragraphe]
\eng{Drug abuse is a serious problem among young people in our country. Many teenagers take drugs because of peer pressure, and they quickly become addicted. When they try to quit, they suffer from painful withdrawal symptoms. To keep young people away from drugs, schools should raise awareness about the dangers of addiction and keep students informed. Families must care for victims instead of rejecting them, and the government should punish drug dealers. Finally, sport and healthy habits can help teenagers stay away from drugs. Prevention is better than cure.}
\end{exemplebox}

\begin{exercice}[À vous d'écrire]
Rédigez votre propre paragraphe de 8 à 10 lignes : \emph{How to keep young people away from drugs}. Utilisez au moins six mots de la leçon : \eng{addicted, withdrawal, awareness, to quit, healthy habits, to care for, prevention}.
\end{exercice}

\courssub{Self-assessment: check your progress}

\begin{aretenir}[Checklist d'auto-évaluation — I can...]
Cochez chaque case honnêtement :
\begin{itemize}
\item[$\square$] I can \textbf{name} ten words about drugs, tobacco, alcohol and aging.
\item[$\square$] I can \textbf{define} \eng{withdrawal symptoms}, \eng{binge drinking} and \eng{awareness}.
\item[$\square$] I can \textbf{use} \eng{addicted to}, \eng{to quit} and \eng{to raise awareness} in sentences.
\item[$\square$] I can \textbf{pronounce} \eng{withdrawal} (« ouiz-drô-al »), \eng{elderly} (« èl-deur-li ») and \eng{awareness} (« a-wère-ness ») correctly.
\item[$\square$] I can \textbf{guess} the meaning of an unknown word from the context.
\item[$\square$] I can \textbf{write} a short paragraph about keeping young people away from drugs.
\end{itemize}
Si une case reste vide, relisez la section correspondante de la leçon.
\end{aretenir}

\begin{savaistu}[Prevention is better than cure]
Le proverbe \eng{Prevention is better than cure} (« Mieux vaut prévenir que guérir ») est très utilisé dans les campagnes de santé publique anglophones. Au Cameroun, les ministères de la Santé publique et de l'Éducation de base organisent régulièrement des campagnes de sensibilisation dans les écoles contre le tabac, l'alcool et la drogue — exactement le thème de notre texte de Douala !
\end{savaistu}

\newpage

\coursec{Activité d'intégration}

\coursec{Lesson 27 — Vocabulary: Words and expressions related to keeping informed about health problems related to drugs, tobacco, alcohol abuse, and aging}

\courssub{Objectifs de l'activité}

À l'issue de cette activité d'intégration, l'élève sera capable de :

\begin{itemize}
\item réemployer à l'oral et à l'écrit le vocabulaire des quatre champs lexicaux étudiés dans la leçon : \cle{drugs and drug abuse}, \cle{tobacco and smoking}, \cle{alcohol abuse} et \cle{aging and elderly health} ;
\item utiliser correctement les collocations apprises (\eng{to quit smoking}, \eng{binge drinking}, \eng{to care for the elderly}, \eng{awareness campaign}\ldots) et au moins une expression idiomatique (\eng{to kick the habit}, \eng{prevention is better than cure}) ;
\item éviter les faux amis et les calques du français (\eng{drugs} $\neq$ « médicaments », \eng{to assist} $\neq$ « assister à ») ;
\item produire un dialogue oral structuré (interview radiophonique) avec un registre de langue adapté : questions claires, réponses développées, formules de transition ;
\item écouter attentivement la production des autres groupes et évaluer l'emploi du lexique à l'aide d'une grille d'observation.
\end{itemize}

\courssub{Situation}

Le lycée de Nkol-Eton, à Yaoundé, participe à la \emph{National Health Awareness Week}. La radio scolaire du lycée, \emph{Eagle FM}, diffuse chaque vendredi une émission en anglais suivie par tous les élèves de Terminale. Cette semaine, le thème imposé est : \eng{“Stay Healthy, Stay Informed”} « Restez en bonne santé, restez informés ». Ton groupe a été sélectionné pour préparer et présenter l'émission de vendredi prochain. Le proviseur a demandé que l'émission aborde les quatre grands problèmes de santé publique étudiés en classe : la drogue, le tabac, l'alcool et la santé des personnes âgées.

\begin{textebox}[Memo from the Principal — Eagle FM, Nkol-Eton High School, Yaoundé]
Dear Terminale A students,

As part of the National Health Awareness Week, our school radio station Eagle FM will broadcast a special five-minute programme entitled “Stay Healthy, Stay Informed” this Friday at 10 a.m.

Your group of four will present the show. One student will act as the radio host, and the three others will play the following guests: a doctor from the Yaoundé Central Hospital, a former smoker who managed to quit, and a student who takes care of an elderly grandparent.

The programme must inform our school community about the dangers of drug abuse, tobacco and binge drinking, and about the importance of caring for the elderly. Each speaker must use precise vocabulary: collocations such as “to raise awareness”, “second-hand smoke” or “to kick the habit” are expected. Remember: prevention is better than cure!

The best show, chosen by the audience, will be recorded and sent to the regional delegation of secondary education.

Good luck, and stay on air!

The Principal
\end{textebox}

\begin{definition}[Glossaire du document]
\begin{itemize}
\item \cle{to broadcast} (v.) : diffuser (à la radio) ; prononciation « brod-kaste ».
\item \cle{host} (n.) : présentateur, animateur d'une émission.
\item \cle{former smoker} (n.) : ancien fumeur (personne qui a arrêté de fumer).
\item \cle{to raise awareness} (coll.) : sensibiliser (litt. « élever la conscience »).
\item \cle{second-hand smoke} (n.) : tabagisme passif.
\item \cle{on air} (expr.) : à l'antenne, en direct.
\end{itemize}
\end{definition}

\courssub{Consignes}

Travail en groupes de quatre. L'activité se déroule en cinq étapes, de la préparation écrite à l'évaluation orale.

\begin{enumerate}
\item \textbf{Relecture et repérage (10 min).} Relisez le mémo du proviseur. Soulignez les quatre invités ou rôles à distribuer et les quatre thèmes obligatoires. Répartissez les rôles dans votre groupe : le \cle{host}, le \cle{doctor}, le \cle{former smoker}, l'\cle{student}.

\item \textbf{Banque de mots (10 min).} Sur une feuille, chaque élève liste, pour son rôle, au moins \textbf{huit mots ou collocations} de la leçon et \textbf{une expression idiomatique}. Exemples attendus : \eng{addiction}, \eng{drug trafficking}, \eng{to quit smoking}, \eng{nicotine}, \eng{binge drinking}, \eng{liver damage}, \eng{to care for the elderly}, \eng{awareness campaign}, \eng{to kick the habit}, \eng{prevention is better than cure}. Vérifiez dans votre cours qu'il n'y a aucun faux ami.

\item \textbf{Écriture du script (25 min).} Rédigez ensemble le script complet de l'émission (environ 350 à 400 mots) : une introduction du présentateur, trois interviews courtes (deux ou trois questions chacune), et une conclusion avec un message de prévention. Utilisez les structures de l'interview : \eng{Thank you for joining us\ldots}, \eng{Could you tell our listeners\ldots?}, \eng{What advice would you give to\ldots?}

\item \textbf{Répétition et représentation (5 min par groupe).} Répétez en veillant à la prononciation des mots difficiles : \eng{alcohol} « al-ko-ol », \eng{elderly} « el-deur-li », \eng{disease} « di-ziz ». Puis jouez l'émission devant la classe, ou enregistrez-la sur un téléphone.

\item \textbf{Écoute active et vote (pendant les passages).} Pendant que les autres groupes jouent, cochez sur la grille d'écoute chaque mot du thème correctement employé, notez les faux amis éventuels, puis votez pour l'émission la plus convaincante et justifiez votre choix en deux phrases en anglais.
\end{enumerate}

\begin{exemplebox}[Grille d'écoute (listening grid) à reproduire]
\begin{center}
\begin{tabularx}{\linewidth}{|X|X|X|}\hline
\rowcolor{popblueL} \textbf{Theme} & \textbf{Words heard (tick)} & \textbf{Correct?} \\ \hline
Drugs & addiction, drug abuse, trafficking\ldots & yes / no \\ \hline
Tobacco & to quit smoking, nicotine, second-hand smoke\ldots & yes / no \\ \hline
Alcohol & binge drinking, liver damage, hangover\ldots & yes / no \\ \hline
Aging & elderly, to care for, retirement\ldots & yes / no \\ \hline
Idiom & to kick the habit / prevention is better than cure & yes / no \\ \hline
\end{tabularx}
\end{center}
\end{exemplebox}

\courssub{Ressources à mobiliser}

\begin{aretenir}[Lexique utile par champ]
\begin{itemize}
\item \textbf{Drugs :} \eng{drug abuse}, \eng{addiction}, \eng{addict}, \eng{drug dealer}, \eng{overdose}, \eng{rehabilitation centre}, \eng{to get hooked on drugs}.
\item \textbf{Tobacco :} \eng{to smoke}, \eng{smoker / non-smoker}, \eng{nicotine}, \eng{lung cancer}, \eng{to quit smoking}, \eng{second-hand smoke}, \eng{to kick the habit}.
\item \textbf{Alcohol :} \eng{alcohol abuse}, \eng{binge drinking}, \eng{drunk}, \eng{hangover}, \eng{liver damage}, \eng{to drive under the influence}.
\item \textbf{Aging :} \eng{the elderly}, \eng{old age}, \eng{to care for}, \eng{retirement}, \eng{ageing population}, \eng{to look after one's grandparents}.
\item \textbf{Sensibilisation :} \eng{awareness campaign}, \eng{to raise awareness}, \eng{prevention}, \eng{health risks}, \eng{prevention is better than cure}.
\end{itemize}
\end{aretenir}

\begin{propriete}[Structures de l'interview radiophonique]
\begin{itemize}
\item Accueillir : \eng{Good morning, dear listeners, and welcome to “Stay Healthy, Stay Informed”.}
\item Présenter un invité : \eng{Our first guest today is Dr Ngo Bassa, from the Yaoundé Central Hospital.}
\item Poser une question ouverte : \eng{Could you tell our listeners why drug abuse is so dangerous for young people?}
\item Relancer : \eng{That's very interesting. And what about tobacco?}
\item Conclure : \eng{Thank you for joining us. Remember: prevention is better than cure!}
\end{itemize}
\end{propriete}

\begin{attention}[Faux amis à éviter dans l'émission]
\begin{itemize}
\item \eng{drugs} signifie « drogues » (stupéfiants) ; pour « médicaments », dites \eng{medicines} ou \eng{medication}.
\item \eng{to assist} signifie « aider » ; « assister à une émission » se dit \eng{to attend} ou \eng{to listen to}.
\item \eng{sensible} signifie « raisonnable » ; « sensible » (qui ressent facilement) se dit \eng{sensitive}.
\item On dit \eng{to be addicted \textbf{to} drugs} (jamais \eng{addicted of}).
\item \eng{the elderly} est un nom collectif pluriel : \eng{The elderly \textbf{need} our help.} (sans \emph{s} à \emph{elderly}).
\end{itemize}
\end{attention}

\courssub{Proposition de production et corrigé commenté}

\begin{exoresolu}[Script modèle de l'émission “Stay Healthy, Stay Informed”]
Énoncé : rédiger et jouer l'émission de radio en respectant toutes les contraintes (quatre rôles, huit mots du thème par intervenant, une expression idiomatique, aucun faux ami).

\tcblower

\textbf{Production modèle (script en anglais) :}

\textbf{Host :} Good morning, dear listeners, and welcome to “Stay Healthy, Stay Informed”, the Eagle FM special programme for the National Health Awareness Week. Today we are going to talk about drugs, tobacco, alcohol and the health of the elderly. Our first guest is Dr Ngo Bassa, from the Yaoundé Central Hospital. Doctor, could you tell our listeners why drug abuse is such a serious problem among young people?

\textbf{Doctor :} Thank you for inviting me. Drug abuse destroys young lives. When teenagers get hooked on drugs, they quickly develop an addiction, and some end up in a rehabilitation centre. An overdose can even kill. That is why our hospital has launched an awareness campaign in schools: we want to raise awareness before it is too late. Remember, dear students: prevention is better than cure.

\textbf{Host :} Thank you, Doctor. Our second guest is Mr Etoundi, a former smoker. Mr Etoundi, how did you manage to quit smoking?

\textbf{Former smoker :} Well, it was very hard because of the nicotine, but I kicked the habit two years ago. I used to smoke twenty cigarettes a day, and my doctor warned me about lung cancer. I also realised that my second-hand smoke was dangerous for my children. So I stopped, and today I feel much healthier. My advice is simple: never start, and if you smoke, quit now.

\textbf{Host :} That's very inspiring. Finally, we welcome Amina, a student in Terminale A who cares for her elderly grandmother. Amina, what does that mean in your daily life?

\textbf{Student :} I look after my grandmother every day: I help her take her medication, I cook for her and I keep her company. Old age can be lonely, so caring for the elderly is everybody's duty. With an ageing population, our society must not forget its grandparents.

\textbf{Host :} Thank you all for joining us. Dear listeners, stay away from drugs, tobacco and binge drinking, take care of the elderly, and stay healthy, stay informed!

\textbf{Commentaire des choix de langue :}
\begin{itemize}
\item Lexique du thème correctement employé : \eng{drug abuse}, \eng{addiction}, \eng{overdose}, \eng{rehabilitation centre}, \eng{awareness campaign}, \eng{to quit smoking}, \eng{nicotine}, \eng{lung cancer}, \eng{second-hand smoke}, \eng{binge drinking}, \eng{medication}, \eng{the elderly}, \eng{ageing population}.
\item Deux expressions idiomatiques intégrées naturellement : \eng{to kick the habit} (« se débarrasser d'une habitude ») et \eng{prevention is better than cure} (« mieux vaut prévenir que guérir »).
\item Faux amis évités : \eng{medication} (et non \emph{drugs}) pour « médicaments » ; \eng{to look after} et \eng{to care for} (et non \emph{to assist}) pour « s'occuper de ».
\item Prépositions correctes : \eng{hooked \textbf{on} drugs}, \eng{advice} reste indénombrable (\eng{my advice is\ldots}).
\item Structure d'interview respectée : formules d'accueil, questions ouvertes, relances, conclusion avec message de prévention.
\end{itemize}
\end{exoresolu}

\courssub{Bilan de l'activité}

Cette activité a permis de mobiliser l'ensemble des savoir-faire de la leçon dans une tâche de communication authentique et motivante. À l'écrit, les élèves ont rédigé un script en réemployant les collocations et les mots de la même famille étudiés ; à l'oral, ils ont travaillé la prononciation, le débit et le registre de langue propre à l'interview radiophonique ; en compréhension orale, ils ont écouté activement les autres groupes grâce à la grille d'écoute. Le vote final et sa justification en anglais développent enfin l'esprit critique et l'argumentation.

\begin{methode}[Pour réussir ce type de tâche d'intégration]
\begin{enumerate}
\item Identifier d'abord les contraintes de la situation (rôles, thèmes, durée, nombre de mots imposés).
\item Constituer une banque de lexique par rôle avant d'écrire, en vérifiant chaque mot dans le cours.
\item Rédiger un script structuré : introduction, développement en interviews, conclusion avec message.
\item Vérifier les faux amis, les prépositions et les collocations avant la représentation.
\item Parler clairement, en soignant la prononciation des mots difficiles, et écouter les autres pour évaluer.
\end{enumerate}
\end{methode}

\begin{experience}[Prolongement possible]
En devoir maison, chaque élève transforme le script de son groupe en un court article pour le journal du lycée, intitulé \eng{“Health Awareness Week at Nkol-Eton”}, de 150 à 180 mots, en réutilisant au moins dix mots du thème et une expression idiomatique. Les meilleurs articles seront affichés dans la salle d'anglais.
\end{experience}

\newpage

\coursec{Méthodes et exercices résolus}

\coursec{Lesson 27 — Vocabulary: Words and expressions related to keeping informed about health problems related to drugs, tobacco, alcohol abuse, and aging}

\begin{definition}[Objectifs de la leçon]
Cette leçon de vocabulaire thématique (Chapter 3 : Environment, Well-being and Health) vise à te permettre de :
\begin{itemize}
\item \textbf{Identifier et employer} le lexique précis des addictions (drogues, tabac, alcool) et du vieillissement.
\item \textbf{Maîtriser les collocations} (associations de mots figées) et les prépositions qui les accompagnent.
\item \textbf{Éviter les faux amis} et les calques syntaxiques du français (erreurs classiques au Bac).
\item \textbf{Dériver les mots} (word formation) pour passer d'une nature grammaticale à l'autre.
\item \textbf{Produire un discours cohérent} (oral ou écrit) de sensibilisation ou de description sur ces thèmes de santé publique.
\end{itemize}
\end{definition}

\begin{textebox}[Situation de communication : Campagne de sensibilisation au lycée de Buea]
\eng{“Good morning, students. Today’s assembly focuses on a silent epidemic: substance abuse. Whether it is tobacco, alcohol or hard drugs, the consequences on your health and future are devastating. Remember, addiction is not a moral failure but a medical condition that requires support. Let us build a drug-free generation together.”}
\end{textebox}

\courssub{Drugs and drug abuse}

\begin{methode}[Construire et utiliser le champ lexical des drogues]
Pour maîtriser ce vocabulaire au Bac, procède méthodiquement :
\begin{enumerate}
\item \textbf{Repère le mot-clé du thème} dans le texte : \eng{drug}, \eng{substance}, \eng{addiction}, \eng{abuse}. Tout le champ lexical gravite autour.
\item \textbf{Classe les mots en familles logiques} :
\begin{itemize}
\item Le produit : \eng{a drug}, \eng{a substance}, \eng{a pill}, \eng{a needle}, \eng{a dose}.
\item La personne : \eng{a drug addict} (ou \eng{an addict}), \eng{a user}, \eng{a dealer} (trafiquant), \eng{a pushers} (argot).
\item L'action : \eng{to take drugs}, \eng{to use substances}, \eng{to inject}, \eng{to sniff/snort}, \eng{to overdose}, \eng{to quit}, \eng{to kick the habit} (familier), \eng{to relapse} (rechuter).
\item La conséquence : \eng{addiction}, \eng{dependence}, \eng{overdose}, \eng{withdrawal symptoms} (symptômes de sevrage), \eng{rehab} (centre de désintoxication, abréviation de \eng{rehabilitation}).
\end{itemize}
\item \textbf{Apprends les collocations}, pas les mots isolés : \eng{to abuse drugs}, \eng{to be addicted to}, \eng{to suffer from addiction}, \eng{to die of an overdose}, \eng{to enter rehab}, \eng{a drug-related crime}.
\item \textbf{Vérifie la préposition} : \eng{addicted \textbf{to}}, \eng{dependent \textbf{on}}, \eng{to suffer \textbf{from}}, \eng{to die \textbf{of/from}}.
\item \textbf{Réemploie le lexique dans une phrase personnelle} liée au contexte camerounais (sensibilisation dans les écoles, rôle des \eng{peer educators}, campagne \eng{“No to Drugs”}).
\end{enumerate}
\end{methode}

\begin{propriete}[Tableau récapitulatif : Prépositions obligatoires du thème]
\begin{center}
\begin{tabularx}{\linewidth}{|X|X|X|}
\hline
\rowcolor{popblueL} \textbf{Français} & \textbf{English} & \textbf{Remarque} \\
\hline
être accro à & \eng{to be addicted \textbf{to}} & Jamais \eng{addicted of/on}. \\
dépendre de & \eng{to be dependent \textbf{on}} & \eng{dependent upon} (plus formel). \\
souffrir de & \eng{to suffer \textbf{from}} & \eng{He suffers from liver disease.} \\
mourir de & \eng{to die \textbf{of/from}} & \eng{of} : cause directe (\eng{of an overdose}) ; \eng{from} : cause indirecte/maladie (\eng{from cancer}). \\
lutter contre & \eng{to fight \textbf{against}} / \eng{to combat} & \eng{The government combats drug trafficking.} \\
sensibiliser à & \eng{to raise awareness \textbf{about/of}} & \eng{A campaign about drug abuse.} \\
\hline
\end{tabularx}
\end{center}
\end{propriete}

\begin{exoresolu}[Vocabulary in context — gap filling]
Énoncé : Complete the following sentences with the correct word from the box: \eng{addicted — overdose — quit — second-hand — abuse — lung}.
\begin{enumerate}
\item Many young people in urban areas \eng{............ drugs without knowing the dangers.}
\item My uncle finally decided to \eng{............ smoking after twenty years.}
\item \eng{............ smoke is dangerous for children who live with smokers.}
\item She became \eng{............ to painkillers after her accident.}
\item Smoking is the main cause of \eng{............ cancer.}
\item He was rushed to hospital after a drug \eng{............ .}
\end{enumerate}
\tcblower
Solution détaillée :
\begin{enumerate}
\item \eng{Many young people in urban areas \textbf{abuse} drugs without knowing the dangers.} — Le verbe \eng{to abuse} (abuser de) est la collocation correcte avec \eng{drugs}. \textbf{Attention} : \eng{to abuse} ne signifie pas « abuser » au sens de « profiter de » (\eng{to take advantage of}).
\item \eng{My uncle finally decided to \textbf{quit} smoking after twenty years.} — \eng{to quit} + V-ing : arrêter de. Synonyme formel : \eng{to give up smoking}.
\item \eng{\textbf{Second-hand} smoke is dangerous for children who live with smokers.} — \eng{second-hand smoke} = fumée passive. On parle aussi de \eng{passive smoking} (tabagisme passif).
\item \eng{She became \textbf{addicted} to painkillers after her accident.} — Structure : \eng{to be/become addicted \textbf{to} + nom}. Jamais « addicted of ».
\item \eng{Smoking is the main cause of \textbf{lung} cancer.} — \eng{lung cancer} = cancer du poumon. Ne pas confondre avec \eng{liver} (foie) → \eng{liver cancer}.
\item \eng{He was rushed to hospital after a drug \textbf{overdose}.} — \eng{an overdose} = une surdose. Collocation : \eng{to die of an overdose}, \eng{to survive an overdose}.
\end{enumerate}
Conclusion : la bonne réponse dépend toujours de la \cle{collocation} et de la \cle{préposition}, pas d'une traduction mot à mot du français.
\end{exoresolu}

\courssub{Tobacco and smoking}

\begin{definition}[Vocabulaire clé : Tobacco]
\begin{itemize}
\item \eng{Tobacco} (nom indénombrable) : le tabac (substance). \eng{A cigarette}, \eng{a cigar}, \eng{a pipe}, \eng{rolling tobacco}, \eng{snuff} (tabac à priser), \eng{chewing tobacco}.
\item \eng{Nicotine} : la nicotine (substance addictive).
\item \eng{Tar} : le goudron (cancérigène).
\item \eng{Carbon monoxide} : le monoxyde de carbone.
\item \eng{A smoker} / \eng{a non-smoker} / \eng{a heavy smoker} (grand fumeur) / \eng{a chain smoker} (fumeur compulsif).
\item \eng{To smoke} / \eng{to light up} (allumer une cigarette) / \eng{to put out a cigarette} (écraser).
\item \eng{Second-hand smoke} / \eng{passive smoking} : tabagisme passif.
\item \eng{Smoke-free} : sans fumée (ex. \eng{a smoke-free zone}).
\item \eng{Health warning} : avertissement sanitaire (sur les paquets).
\item \eng{Plain packaging} : paquet neutre.
\end{itemize}
\end{definition}

\begin{exemplebox}[Collocations "Tobacco"]
\begin{itemize}
\item \eng{To \textbf{ban} smoking in public places} (interdire).
\item \eng{To \textbf{increase taxes on} tobacco products} (augmenter les taxes).
\item \eng{Graphic health warnings \textbf{deter} young people from starting} (dissuader).
\item \eng{Smoking \textbf{harms} nearly every organ of the body} (nuire à).
\item \eng{A \textbf{smoking cessation} programme} (programme d'arrêt du tabac).
\end{itemize}
\end{exemplebox}

\courssub{Alcohol abuse}

\begin{definition}[Vocabulaire clé : Alcohol]
\begin{itemize}
\item \eng{Alcohol} (indénombrable) : l'alcool. \eng{An alcoholic drink/beverage}.
\item \eng{Alcohol abuse} : abus d'alcool / consommation nocive. \eng{Alcoholism} (ou \eng{alcohol use disorder}) : l'alcoolisme (maladie).
\item \eng{An alcoholic} : un alcoolique (personne). \eng{A heavy drinker} : un gros buveur.
\item \eng{Binge drinking} : beuverie express / biture (consommation massive en peu de temps), fréquent chez les étudiants (\eng{students}).
\item \eng{To drink} / \eng{to have a drink} / \eng{to booze} (familier) / \eng{to get drunk} (s'enivrer) / \eng{to be intoxicated} (être ivre — terme médical/légal).
\item \eng{A hangover} : la gueule de bois. \eng{Withdrawal symptoms} : tremblements, hallucinations (\eng{delirium tremens}).
\item \eng{Liver cirrhosis} / \eng{liver damage} : cirrhose / atteintes hépatiques.
\item \eng{Drunk driving} (US) / \eng{drink-driving} (UK) : conduite en état d'ivresse. \eng{Breathalyser test} (test d'alcoolémie).
\item \eng{Sobriety} : la sobriété. \eng{To stay sober} : rester sobre.
\end{itemize}
\end{definition}

\begin{attention}[Faux ami majeur : \eng{Intoxicated}]
\eng{Intoxicated} signifie « en état d'ivresse » (sous l'effet de l'alcool ou de la drogue), \textbf{pas} « intoxiqué » au sens médical d'empoisonnement (food poisoning = \eng{food poisoning} / \eng{intoxication alimentaire}). \\
Ex. : \eng{He was arrested for driving while intoxicated (DWI).} → Il a été arrêté pour conduite en état d'ivresse.
\end{attention}

\courssub{Aging and elderly health}

\begin{methode}[Parler de la vieillesse avec correction et respect]
\begin{enumerate}
\item \textbf{Utilise les termes respectueux} : préfère \eng{elderly people}, \eng{senior citizens}, \eng{older adults} ou \eng{the elderly} (nom collectif) à « old people » (souvent perçu comme brutal ou irrespectueux).
\item \textbf{Maîtrise le lexique médical de base} :
\begin{itemize}
\item \eng{Ageing} (UK) / \eng{Aging} (US) : le vieillissement (processus).
\item \eng{Life expectancy} : l'espérance de vie.
\item \eng{Wrinkles} : les rides. \eng{Grey hair} (UK) / \eng{Gray hair} (US) : cheveux blancs/gris.
\item \eng{Memory loss} / \eng{dementia} / \eng{Alzheimer's disease} : perte de mémoire / démence / Alzheimer.
\item \eng{Arthritis} : l'arthrite / rhumatismes. \eng{Osteoporosis} : l'ostéoporose.
\item \eng{Poor eyesight} / \eng{hearing loss} : mauvaise vue / perte auditive.
\item \eng{A walking stick} / \eng{a walking frame} / \eng{a wheelchair} : canne / déambulateur / fauteuil roulant.
\item \eng{A retirement home} / \eng{a nursing home} / \eng{a care home} : maison de retraite / EHPAD.
\item \eng{A caregiver} / \eng{a carer} (UK) : aidant / soignant.
\end{itemize}
\item \textbf{Attention à l'orthographe} : \eng{ageing} (britannique, garde le \eng{e}) / \eng{aging} (américain). Les deux sont acceptés au Bac, mais sois cohérent dans ta copie.
\item \textbf{Construis des phrases équilibrées} : \eng{As people grow older, they often suffer from arthritis.} — \eng{Elderly people need regular medical check-ups.}
\item \textbf{Relie au contexte local} : au Cameroun, les personnes âgées vivent souvent avec leur \eng{extended family} (famille élargie) plutôt qu'en \eng{retirement home}. Le respect des aînés (\eng{respect for elders}) est une valeur culturelle forte.
\end{enumerate}
\end{methode}

\begin{exoresolu}[Reading comprehension — short text on aging]
Énoncé : Read the text and answer the questions in complete sentences.

\eng{“In many Cameroonian families, grandparents play a central role: they look after grandchildren and transmit traditions. However, ageing brings health challenges such as arthritis, memory loss and poor eyesight. Unlike in Western countries, few elderly people live in retirement homes; most are cared for by their extended family. Doctors advise regular check-ups, a balanced diet and light exercise to age in good health.”}

\begin{enumerate}
\item What role do grandparents play in Cameroonian families?
\item Mention two health problems linked to ageing.
\item Where do most elderly Cameroonians live?
\item What do doctors advise?
\end{enumerate}
\tcblower
Solution détaillée :
\begin{enumerate}
\item \eng{Grandparents look after their grandchildren and transmit traditions.} — On reprend les deux éléments du texte avec le verbe \eng{to look after} (s'occuper de), \cle{phrasal verb} essentiel.
\item \eng{Two health problems linked to ageing are arthritis and memory loss.} (On peut ajouter \eng{poor eyesight}). — Le texte donne une liste explicite ; il suffit d'en extraire deux éléments.
\item \eng{Most elderly Cameroonians live with their extended family, not in retirement homes.} — La réponse complète inclut le contraste introduit par \eng{unlike}.
\item \eng{Doctors advise regular check-ups, a balanced diet and light exercise.} — Structure : \eng{to advise + nom}. On aurait aussi : \eng{Doctors advise elderly people to have regular check-ups} (\eng{advise + somebody + to}-infinitif).
\end{enumerate}
Conclusion : en compréhension écrite, réponds par une \cle{phrase complète} en reprenant le vocabulaire du texte, sans recopier tout le paragraphe.
\end{exoresolu}

\courssub{Word families, pronunciation and spelling}

\begin{propriete}[Tableau de dérivation (Word Formation) — Familles du thème]
\begin{center}
\begin{tabular}{|l|l|l|l|l|}
\hline
\rowcolor{popblueL} \textbf{Verbe} & \textbf{Nom (Personne)} & \textbf{Nom (Concept/Chose)} & \textbf{Adjectif} & \textbf{Adverbe} \\
\hline
\eng{to addict} & \eng{an addict} & \eng{addiction} & \eng{addictive} & \eng{addictively} \\
\hline
\eng{to smoke} & \eng{a smoker} & \eng{smoking} & \eng{smoky} & — \\
\hline
\eng{to drink} & \eng{a drinker} & \eng{drink/drinking} & \eng{drunken} (ivrogne) & \eng{drunkenly} \\
\hline
— & — & \eng{alcohol} & \eng{alcoholic} & \eng{alcoholically} \\
\hline
— & \eng{an alcoholic} & \eng{alcoholism} & — & — \\
\hline
\eng{to abuse} & \eng{an abuser} & \eng{abuse} & \eng{abusive} & \eng{abusively} \\
\hline
\eng{to depend} & \eng{a dependent} & \eng{dependence} / \eng{dependency} & \eng{dependent} (on) & \eng{dependently} \\
\hline
\eng{to prevent} & — & \eng{prevention} & \eng{preventive} / \eng{preventable} & \eng{preventively} \\
\hline
\eng{to treat} & \eng{a patient} / \eng{a therapist} & \eng{treatment} & \eng{treatable} & — \\
\hline
\eng{to cure} & — & \eng{a cure} & \eng{curable} / \eng{incurable} & — \\
\hline
\eng{to age} & — & \eng{age} / \eng{ageing/aging} & \eng{aged} / \eng{elderly} & — \\
\hline
\eng{to die} & — & \eng{death} & \eng{dead} / \eng{dying} & — \\
\hline
— & — & \eng{health} & \eng{healthy} / \eng{unhealthy} & \eng{healthily} \\
\hline
\eng{to harm} & — & \eng{harm} & \eng{harmful} / \eng{harmless} & \eng{harmfully} \\
\hline
\end{tabular}
\end{center}
\end{propriete}

\begin{methode}[Dériver un mot selon la nature demandée (Méthode Bac)]
\begin{enumerate}
\item \textbf{Identifie la nature grammaticale attendue} dans la phrase : nom, verbe, adjectif ou adverbe ? Regarde ce qui précède le blanc (article \eng{the/a} → nom ; sujet → verbe ; \eng{very/so} → adjectif ; verbe → adverbe).
\item \textbf{Connais les suffixes clés} :
\begin{itemize}
\item Noms : \eng{-tion/-sion} (\eng{addiction, prevention}), \eng{-ment} (\eng{treatment}), \eng{-ance/-ence} (\eng{dependence}), \eng{-ism} (\eng{alcoholism}), \eng{-er/-or} (\eng{smoker}), \eng{-ist} (\eng{specialist}).
\item Adjectifs : \eng{-ive} (\eng{addictive, preventive}), \eng{-al} (\eng{medical}), \eng{-ous} (\eng{harmful}), \eng{-ful/-less} (\eng{harmful, harmless}), \eng{-able/-ible} (\eng{treatable, curable}).
\item Adverbes : \eng{-ly} (\eng{harmfully}) — attention à \eng{healthily} (y → i).
\end{itemize}
\item \textbf{Mémorise les familles du thème} (voir tableau ci-dessus).
\item \textbf{Vérifie l'orthographe de la dérivation} : \eng{to die → dying → death} ; \eng{healthy → health} (pas de « a » final au nom) ; \eng{alcohol} (un seul \eng{c}, un seul \eng{o} initial).
\item \textbf{Relis la phrase complète} pour vérifier le sens.
\end{enumerate}
\end{methode}

\begin{exoresolu}[Word formation — give the correct form]
Énoncé : Complete each sentence with the correct form of the word in brackets.
\begin{enumerate}
\item Cigarettes are highly \eng{............} (addict).
\item The government launched a campaign for the \eng{............} of drug abuse (prevent).
\item \eng{............} is a disease that destroys families (alcoholic).
\item Passive \eng{............} kills non-smokers every year (smoke).
\item Regular exercise keeps elderly people \eng{............} (health).
\end{enumerate}
\tcblower
Solution détaillée :
\begin{enumerate}
\item \eng{Cigarettes are highly \textbf{addictive}.} — Après l'adverbe \eng{highly}, il faut un adjectif : \eng{addictive} (qui crée une dépendance). Ne pas confondre avec \eng{addicted} (qui décrit la personne dépendante : \eng{He is addicted to cigarettes}).
\item \eng{...the \textbf{prevention} of drug abuse.} — Après l'article \eng{the}, il faut un nom : \eng{prevention}. Collocation à retenir : \eng{prevention campaigns}.
\item \eng{\textbf{Alcoholism} is a disease...} — Sujet de la phrase : il faut le nom abstrait \eng{alcoholism} (l'alcoolisme). \eng{An alcoholic} désigne la personne (un alcoolique).
\item \eng{Passive \textbf{smoking} kills non-smokers...} — Après l'adjectif \eng{passive}, il faut un nom : \eng{smoking} (nom en \eng{-ing}). \eng{Passive smoking} = tabagisme passif.
\item \eng{...keeps elderly people \textbf{healthy}.} — Structure \eng{keep + complément + adjectif} : \eng{healthy}. Attention à l'orthographe : le nom est \eng{health}, l'adjectif \eng{healthy}.
\end{enumerate}
Conclusion : la dérivation exige d'abord l'analyse de la \cle{nature grammaticale} du mot attendu, puis le choix du \cle{suffixe} correct.
\end{exoresolu}

\begin{aretenir}[Prononciation : pièges fréquents]
\begin{itemize}
\item \eng{Addict} / \eng{addicted} / \eng{addictive} : accent sur la 1\up{re} syllabe \eng{AD-dict}, \eng{ad-DIC-ted}, \eng{ad-DIC-tive}. Ne pas dire « ad-dic-tif ».
\item \eng{Alcohol} : /ˈælkəhɒl/ → « al-co-ol » (3 syllabes, pas « al-col »).
\item \eng{Cigarette} : accent sur la dernière syllabe \eng{cig-a-RETTE} (pas sur la première comme en français).
\item \eng{Overdose} : /ˈəʊvədəʊs/ → « o-ve-rdose » (accent sur \eng{o} et \eng{dose}).
\item \eng{Ageing/Aging} : /ˈeɪdʒɪŋ/ → « éi-djin » (le \eng{g} est doux \eng{/dʒ/} comme dans \eng{village}).
\item \eng{Arthritis} : /ɑːˈθraɪtɪs/ → « ar-thraï-tis » (th = /θ/).
\item \eng{Healthy} : /ˈhɛlθi/ → « èl-thi » (th = /θ/, pas de son /g/).
\end{itemize}
\end{aretenir}

\courssub{Using the vocabulary in context (Speaking \& Writing)}

\begin{methode}[Rédiger un paragraphe de sensibilisation (Writing)]
\begin{enumerate}
\item \textbf{Annonce le problème} avec une phrase générale au présent de vérité générale : \eng{Drug abuse is a growing problem among young people in our towns.}
\item \textbf{Développe avec deux ou trois arguments} en utilisant le lexique étudié : \eng{Many teenagers become addicted to drugs and drop out of school.}
\item \textbf{Propose des solutions} concrètes : \eng{The government should organise prevention campaigns in schools.} / \eng{Parents must talk to their children about the dangers of smoking.}
\item \textbf{Conclus} par un appel ou une opinion personnelle : \eng{In my opinion, health education is the best weapon against addiction.}
\item \textbf{Relis} en vérifiant : prépositions (\eng{addicted to}, \eng{suffer from}), temps (présent de vérité générale), et absence de calques.
\end{enumerate}
\end{methode}

\begin{exoresolu}[Guided writing — awareness paragraph]
Énoncé : Write a paragraph of about 80 words on the following topic: \emph{“Why should young people avoid tobacco, alcohol and drugs?”} Use at least five words from the lesson (\eng{addicted, harmful, prevention, health, quit, abuse}...).
\tcblower
Solution détaillée — démarche : on suit le plan problème → arguments → solutions → conclusion, en présent de vérité générale.

Réponse rédigée possible :

\eng{“Tobacco, alcohol and drugs are extremely harmful to young people's health. A teenager who starts smoking can quickly become addicted to nicotine and may later suffer from lung cancer. Alcohol abuse damages the liver and destroys families, while drugs can lead to an overdose or even death. Moreover, addicted students often drop out of school. To solve this problem, the government should organise prevention campaigns in schools and on the radio. Parents also have a role to play: they must inform their children about these dangers. In my opinion, it is easier to avoid these products than to quit later.”}

Vérification : cinq mots du lexique sont bien employés (\eng{harmful, addicted to, abuse, prevention, quit}) ; prépositions correctes (\eng{addicted to}, \eng{suffer from}, \eng{drop out of}) ; environ 90 mots ; aucun calque du français.

Conclusion : une bonne production courte combine \cle{lexique précis}, \cle{structures simples maîtrisées} et \cle{plan clair}.
\end{exoresolu}

\begin{experience}[Activité orale : "Peer Educator Role-Play" (Binôme)]
\textbf{Contexte :} Tu es un \eng{peer educator} (éducateur pair) dans ton lycée à Yaoundé. Ton camarade hésite à fumer sa première cigarette / boire de l'alcool en soirée / essayer une "pilule" pour réviser.
\textbf{Consigne :} En 2 minutes, dissuade-le en utilisant \textbf{au moins 5 mots/expressions} de la leçon (\eng{addicted to, harmful, health, risk, cancer, liver, overdose, quit, prevention, peer pressure}).
\textbf{Structure suggérée :}
\begin{enumerate}
\item \eng{“Hey, don't do it. It's really harmful because...”}
\item \eng{“You risk becoming addicted to...”}
\item \eng{“Think about your health / your future exams.”}
\item \eng{“It's hard to quit later. Prevention is better.”}
\item \eng{“Let's go play football / eat a mango instead.”}
\end{enumerate}
\textbf{Évaluation par les pairs :} Lexique utilisé / Fluidité / Prononciation / Ton convaincant.
\end{experience}

\begin{savaistu}[Culture \& Contexte : Santé publique au Cameroun et dans le monde]
\begin{itemize}
\item \textbf{Cameroun :} Le Ministère de la Santé Publique (MINSANTE) mène des campagnes \eng{“Non au tabac”} et \eng{“Non à la drogue”} dans les établissements scolaires. La loi n°2006/018 du 29 décembre 2006 interdit de fumer dans les lieux publics. L'OMS (\eng{WHO}) estime que le tabagisme tue plus de 8 millions de personnes par an dans le monde.
\item \textbf{Anglophonie :} Au Royaume-Uni, le \eng{NHS (National Health Service)} offre des programmes gratuits \eng{Stop Smoking Services}. Aux USA, la crise des \eng{opioids} (opioïdes) est une urgence nationale (\eng{opioid epidemic}). L'âge légal pour l'alcool est de 21 ans aux USA (\eng{legal drinking age}), 18 ans au UK et au Cameroun.
\item \textbf{Vieillissement :} L'espérance de vie (\eng{life expectancy}) au Cameroun est d'environ 60 ans (Banque Mondiale, 2022). La prise en charge des \eng{elderly} repose traditionnellement sur la \eng{extended family}, mais l'urbanisation crée de nouveaux défis (\eng{isolation, poverty}).
\end{itemize}
\end{savaistu}

\begin{attention}[Les 5 erreurs qui coûtent des points au Bac]
\begin{enumerate}
\item \textbf{Préposition fautive} : écrire « addicted of » ou « dependent of » au lieu de \eng{addicted \textbf{to}} et \eng{dependent \textbf{on}}. De même : \eng{to suffer \textbf{from}}, \eng{to die \textbf{of}}, \eng{to listen \textbf{to}}.
\item \textbf{Faux ami \eng{drugs}} : utiliser \eng{drugs} pour « médicaments ». Pour un médicament, dis \eng{medicine} ou \eng{medication} ; \eng{drugs} évoque d'abord les stupéfiants.
\item \textbf{Calque du temps français} : « She smokes since ten years » est faux. Avec \eng{since/for}, l'anglais exige le present perfect : \eng{She \textbf{has smoked} for ten years} ou \eng{\textbf{has been smoking} since 2015}.
\item \textbf{Orthographe des dérivés} : « healty » au lieu de \eng{healthy}, « dieing » au lieu de \eng{dying}, « alcool » au lieu de \eng{alcohol} (un seul « c », un seul « o » initial : a-l-c-o-h-o-l).
\item \textbf{Article possessif oublié} : « Smoking is bad for the health » est un calque ; l'anglais dit \eng{Smoking is bad for \textbf{your} health}. De même : \eng{He broke \textbf{his} leg}, jamais « the leg ».
\end{enumerate}
\end{attention}

\begin{exercice}[Entraînement personnel — À faire chez soi]
\begin{enumerate}
\item \textbf{Traduction :} Traduisez en anglais en soignant le vocabulaire et la grammaire.
\begin{enumerate}
\item Mon oncle est alcoolique depuis dix ans ; il a une cirrhose du foie.
\item Les campagnes de prévention sont essentielles pour dissuader les jeunes de fumer.
\item Le tabagisme passif tue des non-fumeurs chaque année.
\item Elle a réussi à arrêter de fumer l'année dernière.
\item Les personnes âgées ont besoin de contrôles médicaux réguliers.
\end{enumerate}
\item \textbf{Word Formation :} Donnez la forme correcte du mot entre parenthèses.
\begin{enumerate}
\item \eng{The \_\_\_\_\_\_\_\_ (addict) potential of heroin is extremely high.}
\item \eng{Doctors recommend a \_\_\_\_\_\_\_\_ (balance) diet for the elderly.}
\item \eng{His \_\_\_\_\_\_\_\_ (depend) on alcohol ruined his marriage.}
\item \eng{Smoking is a \_\_\_\_\_\_\_\_ (prevent) cause of death.}
\item \eng{The patient is \_\_\_\_\_\_\_\_ (recover) well after treatment.}
\end{enumerate}
\item \textbf{Expression écrite :} Rédigez un article de 120 mots pour le journal de votre école intitulé \eng{“Say No to Drugs: Protect Your Future”}. Utilisez au moins 8 mots du vocabulaire de la leçon.
\end{enumerate}
\end{exercice}

\begin{corrige}[Exercice n°1 -- Traduction]
\begin{enumerate}
\item \eng{My uncle has been an alcoholic for ten years; he has liver cirrhosis.} (Present perfect \eng{has been} pour \eng{depuis} + durée ; \eng{liver cirrhosis}).
\item \eng{Prevention campaigns are essential to deter young people from smoking.} (\eng{deter from + V-ing} ; \eng{prevention campaigns} collocation).
\item \eng{Passive smoking kills non-smokers every year.} (\eng{Passive smoking} = tabagisme passif ; \eng{non-smokers}).
\item \eng{She managed to quit smoking last year.} (\eng{manage to} = réussir à ; \eng{quit smoking} ; \eng{last year} → prétérit).
\item \eng{Elderly people need regular medical check-ups.} (\eng{Elderly people} respectueux ; \eng{check-ups} nom composé).
\end{enumerate}
\end{corrige}

\begin{corrige}[Exercice n°2 -- Word Formation]
\begin{enumerate}
\item \eng{addictive} (adjectif après \eng{The} + adverbe \eng{extremely} → adjectif).
\item \eng{balanced} (participe passé adjectival, collocation \eng{balanced diet}).
\item \eng{dependence} (nom après possessif \eng{His} ; \eng{dependency} possible mais \eng{dependence} plus courant pour l'addiction).
\item \eng{preventable} (adjectif signifiant « qu'on peut prévenir », collocation \eng{preventable cause}).
\item \eng{recovering} (adjectif après \eng{is} ; \eng{recovering well} = en voie de guérison). Attention : \eng{recovered} = guéri (état fini).
\end{enumerate}
\end{corrige}

\begin{corrige}[Exercice n°3 -- Expression écrite (Proposition de correction)]
\eng{“Say No to Drugs: Protect Your Future}

\eng{Drug abuse is a major threat to the youth of Cameroon. Many students start smoking or drinking alcohol because of peer pressure, ignoring the terrible consequences. These substances are highly addictive and harmful. A young person who experiments with drugs risks becoming an addict, suffering from memory loss, or dying of an overdose. Alcohol abuse destroys the liver and leads to violence. Tobacco causes lung cancer and harms non-smokers through passive smoking.}

\eng{Prevention is our best weapon. Schools must organise awareness campaigns. Parents should talk openly with their children. As students, we must say NO to drugs and YES to life, sports and education. Remember: it is easier to never start than to quit later. Your health is your wealth. Protect your future today!”}

\textbf{Analyse :} ~130 mots. Vocabulaire utilisé : \eng{drug abuse, peer pressure, addictive, harmful, addict, overdose, alcohol abuse, liver, tobacco, lung cancer, passive smoking, prevention, awareness, quit, health}. Structure : Problème → Conséquences → Solutions → Appel final. Grammaire : Présent, modaux (\eng{must, should}), gérondifs, prépositions correctes.
\end{corrige}

\begin{popfigure}
\begin{tikzpicture}[
  mindmap,
  grow cyclic,
  every node/.style={concept, execute at begin node=\hskip0pt},
  concept color=popblue!80,
  level 1/.append style={level distance=4.5cm, sibling angle=90, concept color=popblue!60},
  level 2/.append style={level distance=3cm, sibling angle=45, concept color=poporange!60},
  text=white,
  font=\small\bfseries
]
\node[concept, concept color=popdark] {Health Issues}
  child[concept color=poppurple] { node {Drugs}
    child { node {Addiction} }
    child { node {Overdose} }
    child { node {Rehab} }
    child { node {Dealer} }
  }
  child[concept color=popgreen] { node {Tobacco}
    child { node {Nicotine} }
    child { node {Lung Cancer} }
    child { node {Passive Smoking} }
    child { node {Quit Smoking} }
  }
  child[concept color=poppink] { node {Alcohol}
    child { node {Alcoholism} }
    child { node {Liver Damage} }
    child { node {Binge Drinking} }
    child { node {Drunk Driving} }
  }
  child[concept color=popgold] { node {Aging}
    child { node {Arthritis} }
    child { node {Memory Loss} }
    child { node {Care Home} }
    child { node {Life Expectancy} }
  };
\end{tikzpicture}
\legende{Carte mentale récapitulative des quatre champs lexicaux de la leçon 27.}
\end{popfigure}

\newpage

\coursec{Exercices}

\courssub{Je vérifie mes connaissances}

\begin{exercice}[QCM : choisir le mot correct]
Entoure la lettre correspondant au mot ou à l'expression qui convient.
\begin{enumerate}
\item He was arrested for \_\_\_\_\_\_\_ drugs in the neighbourhood. \hfill (a) dealing \quad (b) treating \quad (c) curing \quad (d) preventing
\item Smoking is strictly \_\_\_\_\_\_\_ in all public buildings. \hfill (a) addicted \quad (b) banned \quad (c) cured \quad (d) aged
\item After years of drinking, he finally decided to \_\_\_\_\_\_\_ alcohol. \hfill (a) give up \quad (b) give in \quad (c) give away \quad (d) give back
\item The \_\_\_\_\_\_\_ of the population is a major challenge for hospitals. \hfill (a) youth \quad (b) ageing \quad (c) addiction \quad (d) withdrawal
\item The government has launched an \_\_\_\_\_\_\_ campaign against drug abuse. \hfill (a) awareness \quad (b) illness \quad (c) overdose \quad (d) recovery
\end{enumerate}
\end{exercice}

\begin{exercice}[Vrai ou faux, à justifier]
Lis attentivement le texte suivant, puis dis si chaque affirmation est \eng{true} ou \eng{false}. Justifie chaque réponse par une courte citation du texte.
\begin{textebox}[An awareness article]
Drug addiction is a disease, not a crime. Families should support drug addicts and encourage them to seek treatment instead of rejecting them. In Cameroon, health workers organise awareness campaigns in schools and markets to inform young people about the dangers of tobacco, alcohol and hard drugs. Experts also warn that the number of elderly people is increasing, and that old people need regular medical care, a balanced diet and the affection of their relatives.
\end{textebox}
\begin{enumerate}
\item According to the text, families should reject drug addicts.
\item Awareness campaigns take place only in hospitals.
\item The number of elderly people is going down.
\item Elderly people need medical care and affection.
\end{enumerate}
\end{exercice}

\begin{exercice}[Vocabulaire : définitions]
Associe chaque mot de la colonne de gauche à sa définition en anglais (colonne de droite). Écris les paires sur ta feuille.
\begin{center}
\begin{tabularx}{\linewidth}{|l|X|}
\hline
\rowcolor{popblueL} \textbf{Word} & \textbf{Definition} \\
\hline
1. withdrawal & a. a person who is physically unable to stop taking drugs \\
\hline
2. addict & b. the unpleasant physical effects of stopping a drug \\
\hline
3. overdose & c. a place where patients receive treatment for addiction \\
\hline
4. rehabilitation centre & d. taking too much of a drug at once, with danger of death \\
\hline
5. second-hand smoke & e. the smoke that non-smokers breathe in from other people's cigarettes \\
\hline
\end{tabularx}
\end{center}
\end{exercice}

\begin{exercice}[Conjugaison et formes à compléter]
Complète chaque phrase avec la forme correcte du verbe entre parenthèses (temps, voix ou forme nominale selon le contexte).
\begin{enumerate}
\item My grandfather \_\_\_\_\_\_\_ (smoke) two packets a day when he was young, but he stopped ten years ago.
\item Many lives \_\_\_\_\_\_\_ (save) every year thanks to awareness campaigns.
\item If young people \_\_\_\_\_\_\_ (know) the dangers of drugs, fewer of them would try them.
\item The doctor advised him \_\_\_\_\_\_\_ (reduce) his alcohol consumption.
\item She \_\_\_\_\_\_\_ (suffer) from back pain since she turned seventy.
\end{enumerate}
\end{exercice}

\courssub{Je m'entraîne}

\begin{exercice}[Traduction de phrases]
Traduis les phrases suivantes en anglais, en employant le vocabulaire de la leçon.
\begin{enumerate}
\item Le tabagisme est la première cause de décès évitable dans le monde.
\item Mon oncle est devenu dépendant de l'alcool après avoir perdu son emploi.
\item Les personnes âgées ont besoin d'une alimentation équilibrée et d'exercice régulier.
\item Le gouvernement devrait interdire la vente de cigarettes aux mineurs.
\item Après une cure de désintoxication, elle s'est complètement rétablie.
\end{enumerate}
\end{exercice}

\begin{exercice}[Correction de faux amis et de calques]
Chaque phrase contient une erreur typique de francophone (faux ami ou calque du français). Réécris chaque phrase en anglais correct.
\begin{enumerate}
\item He is very sensible to drugs. \emph{(sens : il réagit fortement aux drogues)}
\item She eventually stopped smoking last year. \emph{(sens : finalement)}
\item He consumes drugs every weekend.
\item My grandmother is an ancient person. \emph{(sens : une personne âgée)}
\item The medicine made him a big good. \emph{(sens : lui a fait beaucoup de bien)}
\end{enumerate}
\end{exercice}

\begin{exercice}[Texte à trous avec banque de mots]
Complète le texte suivant avec les mots de la banque. Attention : il y a deux mots en trop.
\begin{center}
\fbox{\eng{withdrawal — addicted — smoke-free — elderly — moderation — awareness — cure — harmless}}
\end{center}
\begin{textebox}[A healthier city]
The mayor of Buea has decided to make all schools \_\_\_\_\_\_\_ zones: nobody is allowed to smoke near the classrooms. A local association runs \_\_\_\_\_\_\_ campaigns to teach teenagers that alcohol should be consumed with \_\_\_\_\_\_\_, or not at all. Young people who become \_\_\_\_\_\_\_ to hard drugs are sent to a rehabilitation centre, where doctors help them through the painful period of \_\_\_\_\_\_\_. Special programmes also provide free medical check-ups for the \_\_\_\_\_\_\_ people of the town.
\end{textebox}
\end{exercice}

\begin{exercice}[Familles de mots : tableau à compléter]
Recopie et complète le tableau suivant avec la forme manquante de chaque famille de mots. Attention à l'orthographe.
\begin{center}
\begin{tabular}{|l|l|l|}
\hline
\rowcolor{popblueL} \textbf{Verb} & \textbf{Noun} & \textbf{Adjective} \\
\hline
to addict & \_\_\_\_\_\_\_\_\_\_\_\_ & addicted \\
\hline
\_\_\_\_\_\_\_\_\_\_\_\_ & alcohol & alcoholic \\
\hline
to smoke & \_\_\_\_\_\_\_\_\_\_\_\_ & \_\_\_\_\_\_\_\_\_\_\_\_ \\
\hline
to heal & health & \_\_\_\_\_\_\_\_\_\_\_\_ \\
\hline
\_\_\_\_\_\_\_\_\_\_\_\_ & care & careful \\
\hline
\end{tabular}
\end{center}
Puis utilise trois des mots du tableau dans des phrases complètes de ton invention sur le thème de la santé.
\end{exercice}

\courssub{Je me prépare au Bac}

\begin{exercice}[Compréhension de texte et questions de langue — type Bac]
Lis le texte suivant, puis réponds aux questions en anglais, par des phrases complètes.
\begin{textebox}[Keeping our communities healthy]
In many Cameroonian towns, the fight against drugs, tobacco and alcohol abuse has become a priority. In Douala, a group of nurses visits secondary schools every month to warn students about the dangers of smoking. “One cigarette may look harmless,” explains Nurse Etame, “but nicotine is one of the most addictive substances in the world. Many of my patients started smoking at fifteen and now, at forty, they cannot climb a flight of stairs without stopping to breathe.”

Alcohol is another serious problem. Although it is legal for adults, drinking too much damages the liver and the brain, and it destroys families. Health workers advise adults who drink to do so with moderation, and they encourage heavy drinkers to join support groups where they can share their experiences.

Meanwhile, Cameroon, like many countries, must also care for a growing number of elderly people. Old age often brings high blood pressure, diabetes and joint pain. Doctors recommend regular check-ups, a balanced diet, light exercise such as walking, and above all, the presence of the family. “An old person who feels loved,” says Nurse Etame with a smile, “is already half cured.”
\end{textebox}
\textbf{A. Comprehension questions}
\begin{enumerate}
\item What do the nurses of Douala do every month?
\item According to Nurse Etame, why is nicotine dangerous?
\item Mention two organs damaged by excessive drinking.
\item What solution is proposed to heavy drinkers?
\item List three health problems mentioned in connection with old age.
\item What, according to the text, is the best “medicine” for an elderly person?
\end{enumerate}
\textbf{B. Language work}
\begin{enumerate}
\item Find in the text: (a) a synonym of “dangerous”; (b) the opposite of “harmful”; (c) a word meaning “a medical examination”.
\item Put into the passive voice: “Drinking too much damages the liver.”
\item Complete with a word from the text: “Nicotine is highly \_\_\_\_\_\_\_: once you start, it is very hard to stop.”
\end{enumerate}
\end{exercice}

\begin{exercice}[Traduction — type Bac]
Traduis le passage suivant en anglais.
\begin{textebox}[Texte à traduire]
De nos jours, beaucoup de jeunes croient que fumer les rend plus attirants. C'est une grave erreur. Le tabac détruit les poumons, coûte cher et rend dépendant. Quant à l'alcool, il doit être consommé avec modération. Les autorités devraient organiser davantage de campagnes de sensibilisation dans les écoles et les quartiers. Enfin, n'oublions pas nos aînés : les personnes âgées méritent notre respect, nos soins et notre affection.
\end{textebox}
\end{exercice}

\begin{exercice}[Expression écrite — type Bac (80--100 mots)]
Your school is organising a health day. As the president of the Health Club, write a short speech (80--100 words) warning students against tobacco, alcohol and drug abuse. In your speech:
\begin{itemize}
\item greet the audience and introduce the topic;
\item present at least two dangers of tobacco, alcohol or drugs;
\item give two pieces of advice using \eng{should}, \eng{had better} or \eng{I advise you to};
\item end with an encouraging message about healthy living and caring for the elderly.
\end{itemize}
\end{exercice}

\newpage

\coursec{Corrigés des exercices}

\begin{corrige}[Exercice 1 — QCM : choisir le mot correct]
\textbf{Réponses :}
\begin{enumerate}
\item \textbf{(a) dealing} — \eng{He was arrested for dealing drugs in the neighbourhood.} \eng{To deal drugs} = vendre de la drogue (trafic). \eng{Treat} = soigner, \eng{cure} = guérir, \eng{prevent} = prévenir : aucun ne se construit avec \eng{drugs} dans ce sens.
\item \textbf{(b) banned} — \eng{Smoking is strictly banned in all public buildings.} \eng{Banned} = interdit. \eng{Addicted} = dépendant, \eng{cured} = guéri, \eng{aged} = âgé : sens impossibles ici.
\item \textbf{(a) give up} — \eng{He finally decided to give up alcohol.} \eng{To give up} = arrêter, renoncer à. Attention aux autres particules : \eng{give in} = céder, \eng{give away} = donner, révéler, \eng{give back} = rendre.
\item \textbf{(b) ageing} — \eng{The ageing of the population is a major challenge for hospitals.} \eng{Ageing} (orthographe britannique, \eng{aging} aux USA) = vieillissement. \eng{Youth} = jeunesse, \eng{addiction} = dépendance, \eng{withdrawal} = sevrage.
\item \textbf{(a) awareness} — \eng{The government has launched an awareness campaign against drug abuse.} Collocation à retenir : \eng{an awareness campaign} = une campagne de sensibilisation.
\end{enumerate}
\textbf{Point de vigilance :} ne jamais traduire « une campagne de sensibilisation » par \eng{a sensitization campaign} (calque du français) : l'anglais dit \eng{an awareness campaign}.
\end{corrige}

\begin{corrige}[Exercice 2 — Vrai ou faux, à justifier]
\begin{enumerate}
\item \textbf{False.} Justification : \eng{“Families should support drug addicts and encourage them to seek treatment instead of rejecting them.”} Le texte dit exactement le contraire : il faut soutenir les toxicomanes, pas les rejeter.
\item \textbf{False.} Justification : \eng{“Health workers organise awareness campaigns in schools and markets.”} Les campagnes ont lieu dans les écoles et les marchés, pas seulement dans les hôpitaux.
\item \textbf{False.} Justification : \eng{“The number of elderly people is increasing.”} Le nombre de personnes âgées augmente (\eng{increase}), il ne diminue pas (\eng{decrease}).
\item \textbf{True.} Justification : \eng{“Old people need regular medical care, a balanced diet and the affection of their relatives.”}
\end{enumerate}
\textbf{Méthode :} au Bac, une justification sans citation exacte du texte ne rapporte pas tous les points. Recopie toujours le segment du texte entre guillemets anglais “ … ”.
\end{corrige}

\begin{corrige}[Exercice 3 — Vocabulaire : définitions]
\textbf{Paires correctes :} 1–b ; 2–a ; 3–d ; 4–c ; 5–e.
\begin{itemize}
\item \eng{Withdrawal} = the unpleasant physical effects of stopping a drug (le sevrage : tremblements, nausées, angoisse quand on arrête une drogue).
\item \eng{An addict} = a person who is physically unable to stop taking drugs (un toxicomane).
\item \eng{An overdose} = taking too much of a drug at once, with danger of death (une surdose).
\item \eng{A rehabilitation centre} = a place where patients receive treatment for addiction (un centre de désintoxication).
\item \eng{Second-hand smoke} = the smoke that non-smokers breathe in from other people's cigarettes (le tabagisme passif).
\end{itemize}
\textbf{Point de vigilance :} \eng{overdose} se prononce « o-veur-dose » et \eng{withdrawal} « wi-zdro-ol » ; attention à l'orthographe de \eng{withdrawal} (un seul \emph{l} final, pas de \emph{e}).
\end{corrige}

\begin{corrige}[Exercice 4 — Conjugaison et formes à compléter]
\begin{enumerate}
\item \eng{My grandfather \textbf{smoked} two packets a day when he was young, but he stopped ten years ago.} — Prétérit simple : habitude révolue, repérée par \eng{when he was young}. (On accepte aussi \eng{used to smoke} / \eng{would smoke}.)
\item \eng{Many lives \textbf{are saved} every year thanks to awareness campaigns.} — Présent passif : \eng{every year} indique une vérité générale ; le sujet \eng{lives} subit l'action (\eng{to be} + participe passé).
\item \eng{If young people \textbf{knew} the dangers of drugs, fewer of them would try them.} — Conditionnel de type 2 (hypothèse irréelle au présent) : \eng{if} + prétérit, \eng{would} + base verbale.
\item \eng{The doctor advised him \textbf{to reduce} his alcohol consumption.} — Structure \eng{advise somebody to do something} : \eng{advise} + complément + infinitif en \eng{to}.
\item \eng{She \textbf{has suffered} from back pain since she turned seventy.} — Present perfect : action commencée dans le passé et qui continue, signalée par \eng{since}.
\end{enumerate}
\textbf{Point de vigilance :} après \eng{since}, le français utilise le présent (« elle souffre depuis… »), l'anglais exige le present perfect : jamais \eng{she suffers since}.
\end{corrige}

\begin{corrige}[Exercice 5 — Traduction de phrases]
\begin{enumerate}
\item \eng{Smoking is the leading cause of preventable death in the world.} — \eng{Preventable} = évitable ; \eng{leading cause} = première cause. Évite le calque \eng{the first cause of avoidable deaths}, lourd et peu naturel.
\item \eng{My uncle became addicted to alcohol after losing his job.} — \eng{Addicted to} + nom ; \eng{after} + gérondif (\eng{losing}).
\item \eng{Elderly people need a balanced diet and regular exercise.} — \eng{Elderly people} ou \eng{the elderly} ; jamais \eng{ancient people}.
\item \eng{The government should ban the sale of cigarettes to minors.} — \eng{To ban} = interdire ; \eng{minors} = mineurs.
\item \eng{After a detox programme, she made a full recovery.} — \eng{A detox programme} = une cure de désintoxication ; \eng{to make a full recovery} = se rétablir complètement. On accepte aussi : \eng{After treatment at a rehabilitation centre, she fully recovered.}
\end{enumerate}
\textbf{Point de vigilance :} « dépendant de » = \eng{addicted to} (ou \eng{dependent on}), jamais \eng{addicted of}.
\end{corrige}

\begin{corrige}[Exercice 6 — Correction de faux amis et de calques]
\begin{enumerate}
\item \eng{He is very \textbf{sensitive} to drugs.} — Faux ami : \eng{sensible} = raisonnable ; « sensible » (qui réagit) = \eng{sensitive}.
\item \eng{She \textbf{finally} stopped smoking last year.} — Faux ami : \eng{eventually} = finalement après une longue attente incertaine ; « finalement, en fin de compte » = \eng{finally} / \eng{in the end}.
\item \eng{He \textbf{takes} drugs every weekend.} — Calque : on ne dit pas \eng{consume drugs} ; la collocation correcte est \eng{to take drugs} (ou \eng{to use drugs}).
\item \eng{My grandmother is an \textbf{elderly} person.} — Faux ami : \eng{ancient} = antique (un monument, une civilisation) ; « âgé » = \eng{elderly} / \eng{old}.
\item \eng{The medicine \textbf{did him a lot of good}.} — Calque de « faire du bien » : \eng{to do somebody good} ; jamais \eng{make good} dans ce sens.
\end{enumerate}
\textbf{Point de vigilance :} ces cinq pièges reviennent très souvent au Bac ; mémorise les paires \eng{sensible/sensitive}, \eng{eventually/finally}, \eng{ancient/elderly}.
\end{corrige}

\begin{corrige}[Exercice 7 — Texte à trous avec banque de mots]
\textbf{Texte complété :}

\eng{The mayor of Buea has decided to make all schools \textbf{smoke-free} zones: nobody is allowed to smoke near the classrooms. A local association runs \textbf{awareness} campaigns to teach teenagers that alcohol should be consumed with \textbf{moderation}, or not at all. Young people who become \textbf{addicted} to hard drugs are sent to a rehabilitation centre, where doctors help them through the painful period of \textbf{withdrawal}. Special programmes also provide free medical check-ups for the \textbf{elderly} people of the town.}

\textbf{Mots en trop :} \eng{cure} et \eng{harmless}.

\textbf{Justifications :} \eng{smoke-free} s'accorde avec \eng{zones} (adjectif composé) ; \eng{awareness campaigns} est la collocation attendue ; \eng{with moderation} = avec modération ; \eng{become addicted to} + nom ; \eng{withdrawal} (nom) suit l'article défini \eng{the} ; \eng{elderly} qualifie \eng{people}.
\end{corrige}

\begin{corrige}[Exercice 8 — Familles de mots : tableau à compléter]
\begin{center}
\begin{tabular}{|l|l|l|}
\hline
\rowcolor{popblueL} \textbf{Verb} & \textbf{Noun} & \textbf{Adjective} \\
\hline
to addict & \textbf{addiction} & addicted \\
\hline
— & alcohol & alcoholic \\
\hline
to smoke & \textbf{smoke / smoking / smoker} & \textbf{smoky} \\
\hline
to heal & health & \textbf{healthy} \\
\hline
\textbf{to care} & care & careful \\
\hline
\end{tabular}
\end{center}
\textbf{Remarques :} la famille d'\eng{alcohol} n'a pas de verbe courant en anglais (on dit \eng{to drink alcohol}) ; pour \eng{smoke}, on accepte \eng{smoke} (la fumée), \eng{smoking} (le tabagisme) ou \eng{smoker} (le fumeur), et l'adjectif \eng{smoky} (enfumé) ; attention à l'orthographe : \eng{health} $\rightarrow$ \eng{healthy} (ajout de \emph{y}), \eng{addict} $\rightarrow$ \eng{addiction}. Le verbe \eng{to addict} existe mais est rare à la voix active (on utilise surtout la forme passive \eng{be addicted to}).

\textbf{Exemples de phrases (modèles) :}
\begin{itemize}
\item \eng{Addiction to tobacco kills thousands of people every year.}
\item \eng{A healthy lifestyle includes a balanced diet and regular exercise.}
\item \eng{Nurses must be careful when they give medicines to elderly patients.}
\end{itemize}
\end{corrige}

\begin{corrige}[Exercice 9 — Compréhension de texte et questions de langue — type Bac]
\textbf{Barème indicatif (sur 20) :} partie A : 6 questions $\times$ 2 = 12 points (1 point pour le sens, 1 point pour la correction grammaticale) ; partie B : 3 questions = 6 points ; présentation générale : 2 points.

\textbf{A. Comprehension questions — réponses modèles :}
\begin{enumerate}
\item \eng{Every month, the nurses of Douala visit secondary schools to warn students about the dangers of smoking.}
\item \eng{According to Nurse Etame, nicotine is dangerous because it is one of the most addictive substances in the world.}
\item \eng{Excessive drinking damages the liver and the brain.}
\item \eng{Heavy drinkers are encouraged to join support groups where they can share their experiences.}
\item \eng{Old age often brings high blood pressure, diabetes and joint pain.}
\item \eng{According to the text, the best “medicine” for an elderly person is the presence and love of the family: “An old person who feels loved is already half cured.”}
\end{enumerate}

\textbf{B. Language work :}
\begin{enumerate}
\item (a) synonyme de \eng{dangerous} : \eng{harmful} n'apparaît pas, mais le texte contient \eng{serious} (\eng{a serious problem}) — réponse attendue : \eng{serious} ; (b) contraire de \eng{harmful} : \eng{harmless} (\eng{“One cigarette may look harmless”}) ; (c) « examen médical » : \eng{check-up}.
\item \eng{The liver is damaged by drinking too much.} — Passif : le complément d'objet \eng{the liver} devient sujet ; \eng{damages} $\rightarrow$ \eng{is damaged} ; l'agent est introduit par \eng{by}.
\item \eng{Nicotine is highly \textbf{addictive}: once you start, it is very hard to stop.}
\end{enumerate}

\textbf{Points de vigilance :} ne pas confondre \eng{addictive} (qui crée la dépendance, qualité de la substance) et \eng{addicted} (dépendant, état de la personne) ; au passif, ne jamais oublier l'auxiliaire \eng{be} : jamais \eng{the liver damaged by…}.
\end{corrige}

\begin{corrige}[Exercice 10 — Traduction — type Bac]
\textbf{Barème indicatif (sur 10) :} 5 segments $\times$ 2 points (1 point pour le sens, 1 point pour la langue).

\textbf{Traduction modèle :}

\eng{Nowadays, many young people believe that smoking makes them more attractive. This is a serious mistake. Tobacco destroys the lungs, is expensive and is addictive. As for alcohol, it should be consumed in moderation. The authorities should organise more awareness campaigns in schools and neighbourhoods. Finally, let us not forget our elders: elderly people deserve our respect, our care and our affection.}

\textbf{Points de vigilance :}
\begin{itemize}
\item « rendre + adjectif » = \eng{make} + complément + adjectif : \eng{smoking makes them more attractive} (jamais \eng{makes them to be}).
\item « rendre dépendant » : \eng{is addictive} ou \eng{makes people addicted} ; éviter le calque \eng{makes dependent}.
\item « avec modération » : \eng{in moderation} ou \eng{with moderation} (les deux acceptés).
\item « nos aînés » : \eng{our elders} ou \eng{our elderly relatives} ; jamais \eng{our ancients}.
\item « mériter » : \eng{to deserve} (faux ami : \eng{to deserve} $\neq$ « desservir »).
\end{itemize}
\end{corrige}

\begin{corrige}[Exercice 11 — Expression écrite — type Bac (80--100 mots)]
\textbf{Barème indicatif (sur 10) :} respect de la tâche et des 4 consignes : 4 points ; richesse et exactitude du vocabulaire de la leçon : 2 points ; grammaire et orthographe : 3 points ; organisation et longueur (80--100 mots) : 1 point.

\textbf{Proposition de rédaction modèle (98 mots) :}

\eng{Dear friends,}

\eng{Good morning, and welcome to our Health Day. Today I want to warn you against tobacco, alcohol and drugs. Smoking destroys your lungs and makes you addicted, because nicotine is a highly addictive substance. Alcohol damages the liver and destroys families, while hard drugs can lead to overdose and death.}

\eng{You should say no to cigarettes, and you had better avoid bad company. I advise you to practise sport and eat a balanced diet instead.}

\eng{Finally, let us care for our elderly relatives: they deserve our love. Stay healthy, stay strong!}

\textbf{Points de vigilance :}
\begin{itemize}
\item Vérifie que les quatre tâches sont visibles : salutation + sujet, deux dangers, deux conseils avec \eng{should} / \eng{had better} / \eng{I advise you to}, message final sur les personnes âgées.
\item \eng{Had better} + base verbale sans \eng{to} : \eng{you had better avoid}, jamais \eng{you had better to avoid}.
\item Compte les mots : en dessous de 80 ou au-dessus de 100, tu perds le point d'organisation.
\item Réemploie le lexique de la leçon (\eng{addicted, addictive, damages, balanced diet, deserve}) : c'est ce que le correcteur attend.
\end{itemize}
\end{corrige}

\newpage

\coursec{Fiche bilan}

\begin{popfigure}
\begin{tikzpicture}[mindmap, grow cyclic, every node/.style={concept, font=\small}, root concept/.append style={concept color=popblue, font=\bfseries\small}, level 1/.append style={level distance=3.4cm, sibling angle=60}, level 2/.append style={level distance=2.1cm, font=\scriptsize}]
\node[root concept]{Keeping informed: health risks}
  child[concept color=poporange]{ node {Drugs} 
    child { node {addiction, overdose} }
    child { node {rehab, dealer} }
  }
  child[concept color=poppurple]{ node {Tobacco} 
    child { node {smoker, lung cancer} }
    child { node {quit smoking} }
  }
  child[concept color=popgreen]{ node {Alcohol} 
    child { node {alcoholic, drunk} }
    child { node {hangover, liver} }
  }
  child[concept color=popgold]{ node {Aging} 
    child { node {elderly, wrinkles} }
    child { node {retirement, care} }
  }
  child[concept color=poppink]{ node {Word families} 
    child { node {addict $\to$ addiction} }
  }
  child[concept color=popblue!60]{ node {False friends} 
    child { node {abuse, sensible} }
  };
\end{tikzpicture}
\legende{Carte mentale : le vocabulaire de la santé face aux drogues, au tabac, à l'alcool et au vieillissement.}
\end{popfigure}

\begin{bilan}
\textbf{1. Lexique clé — Drugs and drug abuse}
\begin{itemize}
\item \eng{a drug} : une drogue ; \eng{drug abuse} : l'abus de drogues ; \eng{a drug addict} : un toxicomane ; \eng{addiction} : la dépendance ; \eng{to be addicted to} : être dépendant de ; \eng{an overdose} : une overdose ; \eng{a drug dealer / trafficker} : un revendeur / trafiquant de drogue ; \eng{rehabilitation (rehab)} : la désintoxication ; \eng{to get hooked on drugs} : devenir accro à la drogue ; \eng{to come off drugs} : se sevrer.
\end{itemize}
\textbf{2. Lexique clé — Tobacco and smoking}
\begin{itemize}
\item \eng{tobacco} : le tabac ; \eng{a smoker / a non-smoker} : un fumeur / un non-fumeur ; \eng{a heavy smoker} : un gros fumeur ; \eng{second-hand (passive) smoking} : le tabagisme passif ; \eng{lung cancer} : le cancer du poumon ; \eng{to quit / give up smoking} : arrêter de fumer ; \eng{nicotine} : la nicotine ; \eng{a cigarette butt} : un mégot.
\end{itemize}
\textbf{3. Lexique clé — Alcohol abuse}
\begin{itemize}
\item \eng{alcohol} : l'alcool ; \eng{an alcoholic} : un alcoolique ; \eng{alcoholism} : l'alcoolisme ; \eng{to be drunk} : être ivre ; \eng{drink-driving} : la conduite en état d'ivresse ; \eng{a hangover} : la gueule de bois ; \eng{liver disease} : une maladie du foie ; \eng{to drink in moderation} : boire avec modération.
\end{itemize}
\textbf{4. Lexique clé — Aging and elderly health}
\begin{itemize}
\item \eng{to age / ageing} : vieillir / le vieillissement ; \eng{the elderly} : les personnes âgées ; \eng{an old people's home} : une maison de retraite ; \eng{retirement} : la retraite ; \eng{a pensioner} : un retraité ; \eng{wrinkles} : les rides ; \eng{to suffer from arthritis} : souffrir d'arthrite ; \eng{health care} : les soins de santé ; \eng{life expectancy} : l'espérance de vie.
\end{itemize}
\textbf{5. Familles de mots et faux amis}
\begin{center}
\begin{tabularx}{\linewidth}{|X|X|X|}\hline
\rowcolor{popblueL} Verbe & Nom & Adjectif \\ \hline
to addict & addiction, addict & addicted, addictive \\ \hline
to smoke & smoke, smoker, smoking & smoky \\ \hline
to drink & drink, drinker & drunk, drunken \\ \hline
to age & age, ageing & aged, elderly \\ \hline
to care & care & careful, careless \\ \hline
\end{tabularx}
\end{center}
\begin{itemize}
\item \eng{abuse} (d'une substance) = abus, consommation excessive ; « abuser quelqu'un » se dit \eng{to take advantage of someone}.
\item \eng{sensible} = raisonnable ; « sensible » se dit \eng{sensitive}.
\item \eng{a drug} = une drogue ou un médicament ; « une droguerie » n'a rien à voir.
\item \eng{actually} = en fait ; « actuellement » se dit \eng{currently}.
\end{itemize}
\textbf{6. Prononciation et orthographe}
\begin{itemize}
\item \eng{alcohol} « a-lko-rol », \eng{addict} « a-dikte », \eng{ageing} « èï-dji-ng », \eng{tough} « taf ».
\item Verbes en -ing : \eng{smoking} (e muet tombe), \eng{quitting} (doublement du t), \eng{ageing / aging} (les deux existent).
\end{itemize}
\textbf{7. Méthode express}
\begin{enumerate}
\item Repérer le thème du texte (drugs, tobacco, alcohol, aging) avant de répondre.
\item Employer les collocations exactes : \eng{to suffer from}, \eng{to be addicted to}, \eng{to give up smoking}.
\item À l'écrit, varier : \eng{the elderly} / \eng{old people} / \eng{senior citizens}.
\end{enumerate}
\end{bilan}

\begin{aretenir}[Checklist avant le Bac]
\begin{itemize}
\item[$\square$] Je connais le vocabulaire des drogues : \eng{addict, addiction, overdose, rehab, dealer}.
\item[$\square$] Je sais parler du tabac : \eng{heavy smoker, passive smoking, lung cancer, to quit smoking}.
\item[$\square$] Je maîtrise le lexique de l'alcool : \eng{alcoholic, drunk, hangover, liver disease}.
\item[$\square$] Je sais décrire le vieillissement : \eng{the elderly, retirement, wrinkles, life expectancy}.
\item[$\square$] J'utilise les bonnes prépositions : \eng{addicted to}, \eng{suffer from}, \eng{die of}, \eng{dependent on}.
\item[$\square$] Je connais les familles de mots : \eng{addict / addiction / addicted / addictive}.
\item[$\square$] J'évite les faux amis : \eng{abuse}, \eng{sensible}, \eng{actually}, \eng{drug}.
\item[$\square$] Je dis \eng{to give up smoking} et non « to stop to smoke ».
\item[$\square$] Je prononce correctement \eng{alcohol}, \eng{addict} et \eng{ageing}.
\item[$\square$] Je sais réemployer ce lexique dans une phrase ou un court paragraphe argumentatif.
\end{itemize}
\end{aretenir}

\findecours

\end{document}
